
%%%%%%%%%%%%%%%%%%[basicstyle=\footnotesize\ttfamily]%%%%%%%%%%%%%%%%%%%%%%%%%%%%%%%%%%%%%%%%%%%%%%%%%%%%%%%%%%%
%
% Package pgfplots.sty documentation.
%
% Copyright 2007-2009 by Christian Feuersaenger.
%
% This program is free software: you can redistribute it and/or modify
% it under the terms of the GNU General Public License as published by
% the Free Software Foundation, either version 3 of the License, or
% (at your option) any later version.
%
% This program is distributed in the hope that it will be useful,
% but WITHOUT ANY WARRANTY; without even the implied warranty of
% MERCHANTABILITY or FITNESS FOR A PARTICULAR PURPOSE.  See the
% GNU General Public License for more details.
%
% You should have received a copy of the GNU General Public License
% along with this program.  If not, see <http://www.gnu.org/licenses/>.
%
%
%%%%%%%%%%%%%%%%%%%%%%%%%%%%%%%%%%%%%%%%%%%%%%%%%%%%%%%%%%%%%%%%%%%%%%%%%%%%%

\ifdefined\pdfvariable
    \pdfvariable minorversion=5
    \pdfvariable objcompresslevel=2
\else\ifdefined\XeTeXversion
    \special{pdf:minorversion 5}
    \special{dvipdfmx:config z 2} % also sets normal compresslevel=2
\else
    \pdfminorversion=5 % to allow compression
    \pdfobjcompresslevel=2
\fi\fi
\documentclass[a4paper,doc2]{ltxdoc}

%%%%%%%%%%%%%%%%%%%%%%%%%%%%%%%%%%%%%%%%%%%%%%%%%%%%%%%%%%%%%%%%%%%%%%%%%%%%%
%
% Package pgfplots.sty documentation.
%
% Copyright 2007/2008 by Christian Feuersaenger.
%
% This program is free software: you can redistribute it and/or modify
% it under the terms of the GNU General Public License as published by
% the Free Software Foundation, either version 3 of the License, or
% (at your option) any later version.
%
% This program is distributed in the hope that it will be useful,
% but WITHOUT ANY WARRANTY; without even the implied warranty of
% MERCHANTABILITY or FITNESS FOR A PARTICULAR PURPOSE.  See the
% GNU General Public License for more details.
%
% You should have received a copy of the GNU General Public License
% along with this program.  If not, see <http://www.gnu.org/licenses/>.
%
%
%%%%%%%%%%%%%%%%%%%%%%%%%%%%%%%%%%%%%%%%%%%%%%%%%%%%%%%%%%%%%%%%%%%%%%%%%%%%%


\usepackage{makeidx}

% 中文支持：xelatex + ctex
\usepackage[UTF8]{ctex}

% DON't let hyperref overload the format of index and glossary.
% I want to do that on my own in the stylefiles for makeindex...
\makeatletter
\let\@old@wrindex=\@wrindex
\makeatother

\usepackage{ifpdf}
\usepackage[pdfborder=0 0 0]{hyperref}
    \hypersetup{%
        colorlinks=true,    % use true to enable colors below:
        linkcolor=blue,%red,
        filecolor=blue,%magenta,
        %pagecolor=blue,%red,
        urlcolor=blue,%cyan,
        citecolor=blue,
        %frenchlinks=false,    % small caps instead of colors
        pdfborder=0 0 0,    % PDF link-darstellung, falls colorlinks=false. 0 0 0: nix. 0 0 1: default.
        %plainpages=false,    % Das ist notwendig, wenn die Seitenzahlen z.T. in Arabischen und z.T. in römischen Ziffern gemacht werden.
        pdftitle=Package PGFPLOTS manual,
        pdfauthor=Dr. Christian Feuersänger,
        %pdfsubject=,
        pdfkeywords={pgfplots pgf tikz tex latex},
    }

\makeatletter
\let\@wrindex=\@old@wrindex
\makeatother

\date{\small{Revision \pgfplotsrevision\space (\pgfplotsrevisiondate)}}

\usepackage{pdftexcmds}

\usepackage{luatex85}

\makeatletter
% disables colorlinks for all following \ref commands
\def\pgfplotsmanualdisablecolorforref{%
    \pgfutil@ifundefined{pgfplotsmanual@oldref}{%
        \let\pgfplotsmanual@oldref=\ref
    }{}%
    \def\ref##1{%
        \begingroup
        \let\Hy@colorlink=\pgfplots@disabled@Hy@colorlink
        \let\Hy@endcolorlink=\pgfplots@disabled@Hy@endcolorlink
        \pgfplotsmanual@oldref{##1}%
        \endgroup
    }%
}%
\def\pgfplots@disabled@Hy@colorlink#1{\begingroup}%
\def\pgfplots@disabled@Hy@endcolorlink{\endgroup}%
\makeatother

% Formatiere Seitennummern im Index:
\newcommand{\indexpageno}[1]{%
    {\bfseries\hyperpage{#1}}%
}

\def\C{\mathbb{C}}
\newcommand{\R}{\mathbb{R}}
\newcommand{\N}{\mathbb{N}}
\newcommand{\Z}{\mathbb{Z}}

\long\def\COMMENTLOWLEVEL#1\ENDCOMMENT{}
\def\ENDCOMMENT{}

\usepackage{textcomp}
\usepackage{booktabs}

\usepackage{calc}
\usepackage[formats]{listings}
%\usepackage{courier} % don't use it -- the '^' character can't be copy-pasted in courier

\usepackage{array}
\lstset{%
    basicstyle=\ttfamily,
    language=[LaTeX]tex, % Seems as if \lstset{language=tex} must be invoked BEFORE loading tikz!?
    tabsize=4,
    breaklines=true,
    breakindent=0pt,
}

\ifpdf
%    \pdfinfo {
%        /Author    (Christian Feuersaenger)
%    }

\else
%    \def\pgfsysdriver{pgfsys-dvipdfm.def}
\fi
%\def\pgfsysdriver{pgfsys-pdftex.def}
\usepackage{pgfplots}
\usepackage{pgfplotstable}

\ifpdf
    % this allows to disable the clickable lib from command line using
    % \pdflatex '\def\pgfplotsclickabledisabled{1}
%%%%%%%%%%%%%%%%%%%%%%%%%%%%%%%%%%%%%%%%%%%%%%%%%%%%%%%%%%%%%%%%%%%%%%%%%%%%%
%
% Package pgfplots.sty documentation.
%
% Copyright 2007/2008 by Christian Feuersaenger.
%
% This program is free software: you can redistribute it and/or modify
% it under the terms of the GNU General Public License as published by
% the Free Software Foundation, either version 3 of the License, or
% (at your option) any later version.
%
% This program is distributed in the hope that it will be useful,
% but WITHOUT ANY WARRANTY; without even the implied warranty of
% MERCHANTABILITY or FITNESS FOR A PARTICULAR PURPOSE.  See the
% GNU General Public License for more details.
%
% You should have received a copy of the GNU General Public License
% along with this program.  If not, see <http://www.gnu.org/licenses/>.
%
%
%%%%%%%%%%%%%%%%%%%%%%%%%%%%%%%%%%%%%%%%%%%%%%%%%%%%%%%%%%%%%%%%%%%%%%%%%%%%%
% SEE pgfplots-macros.tex as well!
%\pdfminorversion=5 % to allow compression
%\pdfobjcompresslevel=2
\documentclass[a4paper,openany]{book}

% Promote `Missing character' reports to full errors
% NOTE: lualatex turns this into a fatal \shipout error (no PDF produced);
% keeping it commented keeps both engines compiling.
%\tracinglostchars=3

\let\bookmaketitle=\maketitle

% -----------------------------------
% this here is from ltxdoc:
\usepackage{doc}[=v2]
\AtBeginDocument{\MakeShortVerb{\|}}
% 缓解超宽 verb/键名引用导致的 Overfull hbox（允许断行时轻微伸展行宽）
\setlength{\emergencystretch}{10em}
%\setlength{\textwidth}{355pt}
%\addtolength\marginparwidth{30pt}
%\addtolength\oddsidemargin{20pt}
%\addtolength\evensidemargin{20pt}
\def\cmd#1{\cs{\expandafter\cmd@to@cs\string#1}}
\def\cmd@to@cs#1#2{\char\number`#2\relax}
\DeclareRobustCommand\cs[1]{\texttt{\char`\\#1}}
\providecommand\marg[1]{%
  {\ttfamily\char`\{}\meta{#1}{\ttfamily\char`\}}}
\providecommand\oarg[1]{%
  {\ttfamily[}\meta{#1}{\ttfamily]}}
\providecommand\parg[1]{%
  {\ttfamily(}\meta{#1}{\ttfamily)}}
\raggedbottom

% -----------------------------------

%%%%%%%%%%%%%%%%%%%%%%%%%%%%%%%%%%%%%%%%%%%%%%%%%%%%%%%%%%%%%%%%%%%%%%%%%%%%%
%
% Package pgfplots.sty documentation.
%
% Copyright 2007/2008 by Christian Feuersaenger.
%
% This program is free software: you can redistribute it and/or modify
% it under the terms of the GNU General Public License as published by
% the Free Software Foundation, either version 3 of the License, or
% (at your option) any later version.
%
% This program is distributed in the hope that it will be useful,
% but WITHOUT ANY WARRANTY; without even the implied warranty of
% MERCHANTABILITY or FITNESS FOR A PARTICULAR PURPOSE.  See the
% GNU General Public License for more details.
%
% You should have received a copy of the GNU General Public License
% along with this program.  If not, see <http://www.gnu.org/licenses/>.
%
%
%%%%%%%%%%%%%%%%%%%%%%%%%%%%%%%%%%%%%%%%%%%%%%%%%%%%%%%%%%%%%%%%%%%%%%%%%%%%%


\usepackage{makeidx}

% 中文支持：xelatex + ctex
\usepackage[UTF8]{ctex}

% DON't let hyperref overload the format of index and glossary.
% I want to do that on my own in the stylefiles for makeindex...
\makeatletter
\let\@old@wrindex=\@wrindex
\makeatother

\usepackage{ifpdf}
\usepackage[pdfborder=0 0 0]{hyperref}
    \hypersetup{%
        colorlinks=true,    % use true to enable colors below:
        linkcolor=blue,%red,
        filecolor=blue,%magenta,
        %pagecolor=blue,%red,
        urlcolor=blue,%cyan,
        citecolor=blue,
        %frenchlinks=false,    % small caps instead of colors
        pdfborder=0 0 0,    % PDF link-darstellung, falls colorlinks=false. 0 0 0: nix. 0 0 1: default.
        %plainpages=false,    % Das ist notwendig, wenn die Seitenzahlen z.T. in Arabischen und z.T. in römischen Ziffern gemacht werden.
        pdftitle=Package PGFPLOTS manual,
        pdfauthor=Dr. Christian Feuersänger,
        %pdfsubject=,
        pdfkeywords={pgfplots pgf tikz tex latex},
    }

\makeatletter
\let\@wrindex=\@old@wrindex
\makeatother

\date{\small{Revision \pgfplotsrevision\space (\pgfplotsrevisiondate)}}

\usepackage{pdftexcmds}

\usepackage{luatex85}

\makeatletter
% disables colorlinks for all following \ref commands
\def\pgfplotsmanualdisablecolorforref{%
    \pgfutil@ifundefined{pgfplotsmanual@oldref}{%
        \let\pgfplotsmanual@oldref=\ref
    }{}%
    \def\ref##1{%
        \begingroup
        \let\Hy@colorlink=\pgfplots@disabled@Hy@colorlink
        \let\Hy@endcolorlink=\pgfplots@disabled@Hy@endcolorlink
        \pgfplotsmanual@oldref{##1}%
        \endgroup
    }%
}%
\def\pgfplots@disabled@Hy@colorlink#1{\begingroup}%
\def\pgfplots@disabled@Hy@endcolorlink{\endgroup}%
\makeatother

% Formatiere Seitennummern im Index:
\newcommand{\indexpageno}[1]{%
    {\bfseries\hyperpage{#1}}%
}

\def\C{\mathbb{C}}
\newcommand{\R}{\mathbb{R}}
\newcommand{\N}{\mathbb{N}}
\newcommand{\Z}{\mathbb{Z}}

\long\def\COMMENTLOWLEVEL#1\ENDCOMMENT{}
\def\ENDCOMMENT{}

\usepackage{textcomp}
\usepackage{booktabs}

\usepackage{calc}
\usepackage[formats]{listings}
%\usepackage{courier} % don't use it -- the '^' character can't be copy-pasted in courier

\usepackage{array}
\lstset{%
    basicstyle=\ttfamily,
    language=[LaTeX]tex, % Seems as if \lstset{language=tex} must be invoked BEFORE loading tikz!?
    tabsize=4,
    breaklines=true,
    breakindent=0pt,
}

\ifpdf
%    \pdfinfo {
%        /Author    (Christian Feuersaenger)
%    }

\else
%    \def\pgfsysdriver{pgfsys-dvipdfm.def}
\fi
%\def\pgfsysdriver{pgfsys-pdftex.def}
\usepackage{pgfplots}
\usepackage{pgfplotstable}

\ifpdf
    % this allows to disable the clickable lib from command line using
    % \pdflatex '\def\pgfplotsclickabledisabled{1}\input{pgfplots.tex}'
    \expandafter\ifx\csname pgfplotsclickabledisabled\endcsname\relax
        \usepgfplotslibrary{clickable}
    \fi
\fi

%\usepackage{fp}
% ATTENTION:
% this requires pgf version NEWER than 2.00:
%\usetikzlibrary{fixedpointarithmetic}

\usepgfplotslibrary{dateplot,units,groupplots,statistics}

\pgfplotsset{
    table/search path={.,plotdata}
}

\usepackage[a4paper,left=2.25cm,right=2.25cm,top=2.5cm,bottom=2.5cm]{geometry}
\usepackage{amsmath,amssymb}
\usepackage{xxcolor}

\usepackage{ifluatex}

% tlmgr install zapfding symbol
\usepackage{pifont}

\def\lua{\textsc{lua}}

\ifluatex
    %\usepackage[utf8]{luainputenc}
\else
    \expandafter\ifx\csname XeTeXrevision\endcsname\relax
        \usepackage[utf8]{inputenc}
    \fi
\fi
\usepackage{amsmath}

% tlmgr eurosym
\usepackage{eurosym}

% tlmgr install units
\usepackage{nicefrac}
\input{pgfplots-macros}

\makeatletter
\@ifpackageloaded{tex4ht}{
}{%
    \IfFileExists{ocg.sty}{%
        \usepackage{ocg}%
    }{%
        \usepackage{pgfplots_ocg_copy}%
    }
}%
\makeatother

\usepackage{nicefrac}

\graphicspath{{figures/}}

\edef\preambleconfig{width=7cm,compat=\pgfkeysvalueof{/pgfplots/compat/mostrecent}}

% generate the following file such that it does not need to be updated
% manually for each release:
\makeatletter
\immediate\openout\w@pgf@writea pgfplots.preliminaries.compatcurrent.tmp
\immediate\write\w@pgf@writea{%
\string\begin{codeexample}[code only]^^J%
\string\usepackage{pgfplots}^^J%
\string\pgfplotsset{compat=\pgfkeysvalueof{/pgfplots/compat/mostrecent}}^^J%
\string\end{codeexample}%
}%
\immediate\closeout\w@pgf@writea


\expandafter\pgfplotsset\expandafter{\preambleconfig}

\pgfplotsset{lua debug=off and silent}

\def\pgfplotsmanualcurlibrary{}

% And now, invoke
%     /codeexample/typeset listing/.add={% Preamble:\pgfplotsset{\preambleconfig}}{}}
% since listings are VERBATIM, I need to do some low-level things
% here to get the correct \catcodes:
\pgfkeys{/codeexample/typeset listing/.add code={%
        \ifcode@execute
            \begingroup
            \def\b@show@header{0}%
            \pgfutil@in@{axis}{#1}%
            \ifpgfutil@in@
                \def\b@show@header{1}%
            \fi
            %
            \pgfutil@in@{addplot}{#1}%
            \ifpgfutil@in@
                \def\b@show@header{1}%
            \fi
            %
            \pgfutil@in@{smithchart}{#1}%
            \ifpgfutil@in@
                \def\b@show@header{1}%
            \fi
            \if1\b@show@header
                {\tiny
                    \% Preamble: \pgfmanualpdfref{\textbackslash pgfplotsset}{\pgfmanual@pretty@backslash pgfplotsset}%
                        \pgfmanual@pretty@lbrace \expandafter\pgfmanualprettyprintpgfkeys\expandafter{\preambleconfig}\pgfmanual@pretty@rbrace
                        \ifx\pgfplotsmanualcurlibrary\pgfutil@empty
                        \else
                            \pgfmanual@pretty@backslash usepgfplotslibrary\pgfmanual@pretty@lbrace \pgfplotsmanualcurlibrary\pgfmanual@pretty@rbrace
                        \fi
                }%
            \fi
            \endgroup
        \fi
    }{},%
    %/codeexample/typeset listing/.show code,
}%
\makeatother

\pgfplotsset{
    %every axis/.append style={width=7cm},
    filter discard warning=false,
}

\pgfqkeys{/codeexample}{%
    every codeexample/.append style={
        width=8cm,
        /pgfplots/legend style={fill=graphicbackground},
        /pgfplots/contour/every contour label/.append style={
            every node/.append style={fill=graphicbackground}
        },
    },
    tabsize=4,
}

\usetikzlibrary{backgrounds,patterns}
% Global styles:
\tikzset{
  shape example/.style={
    color=black!30,
    draw,
    fill=yellow!30,
    line width=.5cm,
    inner xsep=2.5cm,
    inner ysep=0.5cm}
}

\newcommand{\FIXME}[1]{\textcolor{red}{(FIXME: #1)}}

% fuer environment 'sidewaysfigure' bspw
% \usepackage{rotating}

\newcommand\Tikz{Ti\textit kZ}
\newcommand\PGF{\textsc{pgf}}
\newcommand\PGFPlots{\pgfplotsmakefilelinkifuseful{pgfplots}{\textsc{pgfplots}}}
\newcommand\PGFPlotstable{\pgfplotsmakefilelinkifuseful{pgfplotstable}{\textsc{PgfplotsTable}}}

\makeindex

% Fix overfull hboxes automatically:
\tolerance=2000
\emergencystretch=10em

\tikzset{prefix=gnuplot/pgfplots_} % prefix for 'plot function'

\author{%
    Dr.\ Christian Feuersänger\\
    {\footnotesize\texttt{cfeuersaenger@users.sourceforge.net}}}%


\usepackage{todonotes}
\newcommand\todosp[2][]{%        % Stefan Pinnow
    \todo[
        disable,
        color=blue!40,
        #1]{#2}
}

% ---------------------------------------------------------------
% 修复：doc 宏包 v2 的 \ltxdoc@meta 在 TeX Live 2026 中未定义，
% 导致 \marg/\oarg/\parg（内部调用 \meta）写 aux/idx 时报
% "Undefined control sequence"。
% 对策（双保险）：
%  1) 兜底提供 \ltxdoc@meta（doc 的 \meta 依赖它）；
%  2) 在 \AtBeginDocument 末尾覆盖 \meta 为鲁棒命令
%     （doc v2 会在 begin document 时重定义 \meta，必须最后覆盖）。
\makeatletter
\providecommand\ltxdoc@meta[1]{%
    {\normalfont\itshape\textlangle #1\textrangle}%
}
\AtBeginDocument{%
    \DeclareRobustCommand\meta[1]{%
        {\normalfont\itshape\textlangle #1\textrangle}%
    }%
}
\makeatother


\makeatletter
% I want a two-column index just as in pgfmanual styles. This here
% was the best way to get one:
\def\index@prologue{\section*{索引}\addcontentsline{toc}{chapter}{索引}
}
\makeatother

%\RequirePackage[german,english,francais]{babel}

\def\matlabcolormaptext{此色图与 Matlab$^\text{\textregistered}$ 中以类似名称发布的色图相似。}

\IfFileExists{tikzlibraryspy.code.tex}{%
\usetikzlibrary{spy}
}{%
    \message{ERROR: tikz SPY library NOT available. The manual will only compile partially.^^J}%
}%

\usepackage{xparse}% for colorbrewer manual

\usetikzlibrary{
    decorations.markings,
    decorations.footprints,
    shapes.arrows,
    matrix,
    positioning,
}

\usepgfplotslibrary{
    fillbetween,
    ternary,
    smithchart,
    patchplots,
    polar,
    colormaps,
    colorbrewer,
    colortol,           % docu not ready yet
}
\pgfqkeys{/codeexample}{%
    every codeexample/.append style={
        /pgfplots/every ternary axis/.append style={
            /pgfplots/legend style={fill=graphicbackground},
        }
    },
    tabsize=4,
}

\pgfplotsmanualenableexternalizationofexpensive

%\usetikzlibrary{external}
%\tikzexternalize[prefix=figures/]{pgfplots}

\title{%
    \PGFPlots{} 宏包手册\\
    {\small 在 \LaTeX{} 中绘制二维/三维图形，版本 \pgfplotsversion}\\
    {\small\url{https://github.com/pgf-tikz/pgfplots}}
    %\\{\small Attention: you are using an unstable development version.}
}

\makeatletter
\long\def\abstractsmuggle{%
    \centering
    \textbf{摘要}\\[0.5cm]

    \begin{minipage}{12cm}
        \PGFPlots{} 直接在 \TeX{} 中以常规或对数坐标绘制高质量函数图，
        界面友好。用户只需提供坐标轴标签、图例条目以及一个或多个图的坐标数据，
        \PGFPlots{} 即可完成坐标轴缩放、对数与刻度计算并绘制图形。
        它支持线图、散点图、分段常数图、柱状图、面积图、网格图与曲面图、
        面片图、等高线图、矢量图、直方图、箱线图、极坐标轴、三元相图、
        史密斯圆图等多种图形。它基于 Till Tantau 开发的 \PGF{}/\Tikz{} 宏包。
    \end{minipage}

}%

\expandafter\date\expandafter{\@date\\[2cm]
    \abstractsmuggle
}%
\makeatother

%\includeonly{pgfplots.reference}


\begin{document}

\def\plotcoords{%
\addplot coordinates {
(5,8.312e-02)    (17,2.547e-02)   (49,7.407e-03)
(129,2.102e-03)  (321,5.874e-04)  (769,1.623e-04)
(1793,4.442e-05) (4097,1.207e-05) (9217,3.261e-06)
};

\addplot coordinates{
(7,8.472e-02)    (31,3.044e-02)    (111,1.022e-02)
(351,3.303e-03)  (1023,1.039e-03)  (2815,3.196e-04)
(7423,9.658e-05) (18943,2.873e-05) (47103,8.437e-06)
};

\addplot coordinates{
(9,7.881e-02)     (49,3.243e-02)    (209,1.232e-02)
(769,4.454e-03)   (2561,1.551e-03)  (7937,5.236e-04)
(23297,1.723e-04) (65537,5.545e-05) (178177,1.751e-05)
};

\addplot coordinates{
(11,6.887e-02)    (71,3.177e-02)     (351,1.341e-02)
(1471,5.334e-03)  (5503,2.027e-03)   (18943,7.415e-04)
(61183,2.628e-04) (187903,9.063e-05) (553983,3.053e-05)
};

\addplot coordinates{
(13,5.755e-02)     (97,2.925e-02)     (545,1.351e-02)
(2561,5.842e-03)   (10625,2.397e-03)  (40193,9.414e-04)
(141569,3.564e-04) (471041,1.308e-04)
(1496065,4.670e-05)
};
}%


\bookmaketitle
\tableofcontents

\chapter{简介}

本宏包提供了方便地生成图和带标签坐标轴的工具。它可以绘制二维和三维的普通图、
对数图以及半对数图。坐标轴刻度、标签、图例（在叠加多幅图时）都可以通过键值
选项来添加。它支持线图、散点图、分段常数图、柱状图、面积图、网格图与曲面图、
补丁图、等高线图、矢量场图、直方图、箱线图、极坐标轴、三元图、史密斯图，
以及更多类型。它还可以循环使用一组预定义的线/标记/颜色规格。

总之，其目的是简化高质量函数图和/或数据图的生成，并解决以下问题：
%
\begin{itemize}
    \item 文档、字体类型与字号的一致性，
    \item 在坐标轴描述中直接使用 \TeX{} 数学模式，
    \item 数据与图的一致性（无需第三方工具），
    \item 通过导言区配置和样式实现文档之间的
        一致性。
\end{itemize}
%
尽管并非必需，也可以使用作为 \Tikz{} 一部分而开发的 |external| 库来生成
独立的 |.pdf| 或 |.eps| 图形。

欢迎使用 \PGFPlots{} 对二维和三维的中等规模数据集进行可视化。

它基于 Till Tantau 开发的 \PGF{}/\Tikz{} 宏包。


\chapter[关于 PGFPlots：预备知识]{关于 {\normalfont\PGFPlots{}}：预备知识}

本节包含关于升级、开发团队、安装（如果你需要手动安装的话）以及
故障排查的信息。除了升级说明之外，你可以完全跳过本节。

\PGFPlots{} 完全构建在 \Tikz{}/\PGF{} 之上。了解 \Tikz{} 会简化使用
\PGFPlots{} 的工作，尽管这并非必需。

\section{组成部分}

\PGFPlots{} 由两个组件构成：
%
\begin{enumerate}
    \item 绘图组件（即你当前正在阅读的部分），以及
    \item \PGFPlotstable{} 组件，它简化数字表格的数字格式化
        与后处理。它作为一个独立的
        宏包提供，并拥有自己的手册
        \href{file:pgfplotstable.pdf}{pgfplotstable.pdf}。
\end{enumerate}


\section{升级说明}

本版本提供了大量改进，感兴趣的读者可以在 \texttt{ChangeLog} 中找到
全部细节。不过，以下几项变更值得特别注意。

\PGFPlots{} 中一个共通的设置是键 |compat|：强烈建议始终把它写进你的
|.tex| 文件。虽然这个键给最终用户带来了一些工作，但它也解决了一个常见需求：
它确保你的 |.tex| 文件总是产生相同的输出，即使你安装了新版本的
\PGFPlots{}。另一方面，它让我们作为维护者能够修复软件缺陷并引入行为变更，
前提是这些变更只影响具有合适 |compat| 级别的文档。|compat| 键的确切影响、
可选取值及其含义将在以下各节中说明。


\subsection{新增可选功能}

\PGFPlots{} 在编写时就考虑了向后兼容：旧的 \TeX{}
文件应当无需修改即可编译，外观也不发生变化。不过，新功能偶尔会导致不同的行为。
在这种情况下，\PGFPlots{} 会停用新功能。\footnote{如果出现破坏向后兼容的情况，
我们深表歉意——并请你提交 bug 报告。我们会负责处理。}

任何导致向后兼容问题的新功能或错误修复都需要\emph{手动}且\emph{显式地}
启用。为此，你应当使用
%
\input pgfplots.preliminaries.compatcurrent.tmp

%
\noindent 放在你的导言区中。这将配置兼容层。

你至少应使用 |compat=1.3|。建议的取值会在运行 \TeX{}
之后打印到 |.log| 文件中。

以下是一份变更清单，除非使用合适的
|compat| 级别，否则这些变更处于停用状态：
%
\begin{enumerate}
    \item \PGFPlots{} 1.18 和 1.17 没有额外约束，在
        兼容级别方面与 1.16 相同。
    \item \PGFPlots{} 1.16 没有额外约束，在
        兼容级别方面与 1.15 相同。
    \item \PGFPlots{} 1.15 启用 |3d log sampling|，并修复
        柱状图的 |clip limits| 问题。
    \item \PGFPlots{} 1.14 改变了系统处理非均匀色图的方式，
        并启用高级 |colormap| 操作（参见
        |of colormap|）。
    \item \PGFPlots{} 1.13 修复极坐标轴中的坐标轴标签，并确保
        |shader=flat| 所选颜色与 |z buffer| 和
        |mesh/ordering| 无关。此外，它为所有堆叠柱状图启用
        |stack negative=separate|。保留较旧的兼容级别是为了
        保留最终用户的临时变通方案。
    \item \PGFPlots{} 1.12 启用 |lua backend| 并定义
        |boxplot/estimator=Excel|。
    \item \PGFPlots{} 1.11 改变了 |axis cs|：它现在是默认
        坐标系。如果你写 |\draw (1,2) -- (2,2);|，\PGFPlots{}
        会自动把它视为
        |\draw (axis cs:1,2) -- (axis cs:2,2);|。
    \item \PGFPlots{} 1.10 在兼容性方面与 1.9 没有区别。
    \item \PGFPlots{} 1.9 附带一个预设，用于组合 |ybar stacked| 和
        |nodes near coords|。此外，它会抑制堆叠柱状图中的空增量。
        要启用这个新预设，你必须
        使用 |compat=1.9| 或更高版本。
    \item \PGFPlots{} 1.8 对标签与刻度比例标签的对齐方式进行了
        新的修订。此外，它改进了
        |hide axis| 的包围盒。该修订通过 |compat=1.8| 或
        更高版本启用。

        配置 |compat=1.8| 是\emph{修复}
        三维坐标轴中 |axis lines=center| 的\emph{必要}条件。
    \item \PGFPlots{} 1.7 新增了以坐标轴单位定义柱宽的选项。
        这些通过 |compat=1.7| 或更高版本启用。
    \item \PGFPlots{} 1.6 新增了用于更精确缩放的选项，以及
        针对 |\addplot3 graphics| 的更多缩放选项。这些通过
        |compat=1.6| 或更高版本启用。
    \item \PGFPlots{} 1.5.1 将圆和椭圆的半径解释为
        \PGFPlots{} 坐标（旧版本使用 \pgfname{} 单位向量，
        它们与 \PGFPlots 没有直接关系）。换句话说：从 1.5.1 版开始，
        可以在坐标轴内部编写 |\draw circle[radius=5]|。这需要
        |\pgfplotsset{compat=1.5.1}| 或更高版本。

        若没有此兼容设置，圆和椭圆将像早期版本那样使用
        \pgfname{} 的底层画布单位。
    \item \PGFPlots{} 1.5 默认使用 |log origin=0|（这会影响
        对数柱状图或堆叠对数图）。较旧的版本
        保持 |log origin=infty|。这需要 |\pgfplotsset{compat=1.5}| 或
        更高版本。
    \item \PGFPlots{} 1.4 修复了若干较小的 bug，
        与早期版本相比可能产生约 $1$--$2\text{pt}$ 的差异。
        这需要 |\pgfplotsset{compat=1.4}| 或更高版本。
    \item \PGFPlots{} 1.3 带来了用户界面上的改进。
        旧版本中“行为选项”与“样式
        选项”之间的技术区分不再必要（不过仍
        完全支持）。

        这总是启用的。
    \item \PGFPlots{} 1.3 有一个新功能，允许自动\emph{把坐标轴标签
        紧贴刻度标签移动}。强烈推荐使用。它需要
        |\pgfplotsset{compat=1.3}| 或更高版本。

        由于这会影响间距，默认情况下并不启用。
    \item \PGFPlots{} 1.3 支持反向坐标轴。不再需要
        使用带负单位的变通方案。\pgfkeys{/pdflinks/search key
        prefixes in/.add={/pgfplots/,}{}}

        请查看 |x dir=reverse| 键。

        现有的变通方案仍能正常工作。使用
        |\pgfplotsset{compat=1.3}| 或更高版本，与 |x dir=reverse| 一起
        切换到新版本。
\end{enumerate}


\subsection{可能需要注意的旧功能}

\begin{enumerate}
    \item |scatter/classes| 功能从 1.3 版起会生成恰当的图例。
        这可能会改变已有使用 |scatter/classes| 的图的图例
        外观。
    \item 从 \PGFPlots{} 1.1 开始，\emph{不应再使用} |\tikzstyle|
        来设置 \PGFPlots{} 选项。

        尽管 |\tikzstyle| 仍被某些较旧的 \PGFPlots{}
        选项支持，但你应当把任何出现的 |\tikzstyle| 替换为
        |\pgfplotsset{|\meta{样式名称}|/.style={|\meta{键值列表}|}}|
        或相应的 |/.append style| 变体。更多细节参见
        Section~\ref{sec:styles}。
\end{enumerate}
%
对于这些变更带来的任何不便，我深表歉意。

\begin{pgfplotskey}{%
    compat=\mchoice{
        1.18,
        1.17,
        1.16,
        1.15,
        1.14,
        1.13,
        1.12,
        1.11,
        1.10,
        1.9,
        1.8,
        1.7,
        1.6,
        1.5.1,
        1.5,
        1.4,
        1.3,
        pre 1.3,
        default
    } (initially default)}
    导言区配置
    %
\input pgfplots.preliminaries.compatcurrent.tmp
    %
    允许在向后兼容与最新功能之间进行选择。
    这个键被设计为文档中遇到的第一个 \PGFPlots{} 键，因为它会准备
    全局选项。

    有时，你可能想在同一
    文档中使用不同的版本。这时，提供
    %
\begin{codeexample}[code only]
\begin{figure}
    \pgfplotsset{compat=1.4}
    ...
    \caption{...}
\end{figure}
\end{codeexample}
    %
    \noindent 以便把兼容设置限制在当前
    上下文中（本例中即 |figure| 环境）。

    查看你的 |.log| 文件输出，即可得到 |compat| 的建议取值。

    使用 |\pgfplotsset{compat=default}| 恢复出厂设置。

    尽管通常没有必要，也可以只启用
    选定的变更，而在总体上保持与旧版本的兼容：
    %
    \begin{pgfplotskeylist}{%
        compat/path replacement=\meta{version},
        compat/labels=\meta{version},
        compat/scaling=\meta{version},
        compat/scale mode=\meta{version},
        compat/empty line=\meta{version},
        compat/plot3graphics=\meta{version},
        compat/bar nodes=\meta{version},
        compat/BB=\meta{version},
        compat/bar width by units=\meta{version},
        compat/pgfpoint substitution=\meta{version},
        compat/general=\meta{version}
    }
        假设我们有一个使用 |\pgfplotsset{compat=1.3}|
        的文档，并且你想保持这种方式。

        此外，你发现 1.5.1 版支持圆和
        椭圆。这时，使用
\begin{codeexample}[]
% preamble:
\pgfplotsset{
    compat=1.3,
    compat/path replacement=1.5.1,
}
\begin{tikzpicture}
\begin{axis}[
    extra x ticks={-2,2},
    extra y ticks={-2,2},
    extra tick style={grid=major}]
    \addplot {x};
    \draw (axis cs:0,0) circle[radius=2];
\end{axis}
\end{tikzpicture}
\end{codeexample}

        所有这些键都接受 |compat| 键的可选值。

        |compat/path replacement| 键控制圆和
        椭圆半径的解释方式。

        |compat/labels| 键控制坐标轴标签的对齐方式：要么
        紧邻刻度，要么使用绝对偏移。从 |1.8| 起，它还
        启用了一套全新的坐标轴标签样式修订。在大多数
        情况下，你看不出区别——但它修复了
        三维坐标轴中的 |axis lines=center|。

        |compat/scaling| 键控制 1.4 和 1.6 版中引入的
        一些 bug 修复：为了提高精度，它们可能带来轻微的
        缩放差异。

        |compat/plot3graphics| 控制
        |\addplot3 graphics| 的新功能。

        |compat/scale mode| 允许启用/禁用警告“The
        content of your 3d axis has CHANGED compared to previous versions”，
        因为 |axis equal| 和 |unit vector ratio| 功能在
        1.6 之前的所有版本中都是坏的，并在 1.6 中得到了修复。

        |compat/empty line| 允许在输入文件中写入空行以
        产生一次跳跃。这需要 |compat=1.4| 或更新版本。详见
        |empty line|。

        |compat/BB| 改变包围盒，使其即使在
        |hide axis| 的情况下也紧凑贴合。

        |compat/bar width by units| 允许用坐标轴单位来表示
        |bar width=1|（即以坐标轴单位计）。

        |compat/bar nodes| 启用 |ybar stacked| 和
        |nodes near coords| 的预设。此外，它为
        堆叠柱状图启用 |stacked ignores zero|。

        |compat/general| 键控制 |log origin|、|lua backend|、
        |enable tick line clipping| 和 |boxplot/estimator|。

        |compat/pgfpoint substitution| 键决定 |axis cs| 是否为
        默认坐标系（自 1.11 起）。

        详细效果可参见本节开头。
    \end{pgfplotskeylist}

    取值 \meta{version} 可以是 |default|、一个版本号，或 |newest|。
    取值 |default| 与 |pre 1.3| 相同（除无关紧要的差异外）。
    强烈\emph{不建议}使用 |newest|：它可能根据 \PGFPlots{}
    的当前版本而改变你的文档。请查看你的 |.log| 文件，
    以获得关于最佳版本的建议。
\end{pgfplotskey}


\section{开发团队}

\PGFPlots{} 主要由 Christian Feuersänger 编写，Pascal Wolkotte 和
Nick Papior Andersen 作为业余项目贡献了许多改进。我们希望它有用，
并能提供有价值的图。

如果你有兴趣编写一些东西但不知从何入手，可以参阅
随 \PGFPlots{} 提供的辅助手册
\href{file:TeX-programming-notes.pdf}{TeX-programming-notes.pdf}。它
远非完整，但也许是一个好的
起点（至少能引导你找到更多文献）。


\section{致谢}

感谢上帝赐予我享受编程的无数时光。我感谢 Pascal Wolkotte 和
Nick Papior Andersen 作为开发团队一员所付出的编程努力与贡献。我感谢 Francesco Poli 为 |contour lua| 实现 |lua| 算法。
我感谢 Jürnjakob Dugge 贡献了
|hist/density|、用于 \verbpdfref{\addplot3} |graphics| 的 matlab 脚本、出色的
用户论坛帮助以及有价值的 bug 报告。我感谢 Stefan Tibus，他贡献了
|\addplot shell| 功能。我感谢 Tom Cashman 贡献了
|reverse legend| 功能。特别感谢 Stefan Pinnow 持续
测试 \PGFPlots{}、讨论需求、提出功能请求和
bug 修复（从而带来众多质量改进），以及移植并
集成 |colorbrewer| 库。此外，我感谢 Schweitzer 教授
富有成效的讨论，以及他最初鼓励我开始这样一个
宏包。感谢 Meine 博士的想法和建议。特别感谢
Markus Böhning 校对了 \PGF{}、\PGFPlots{} 和
\PGFPlotstable{} 的所有手册。感谢 Vincent A. Traag 把 |colorbrewer| 配色
带给 \PGFPlots{}。也感谢众多国际贡献者，他们
提出了功能请求、发现了 bug，或仅仅是改进了
手册！

最后但同样重要的是，我感谢 Till Tantau 和 Mark Wibrow 打造了出色的
图形（及更多）宏包 \PGF{} 和 \Tikz{}，它们正是
\PGFPlots{} 的基础。


\input{pgfplots.install}

% main=pgfplots.tex

\chapter{循序渐进教程}
{%
\tikzset{external/figure name/.add={}{tutorials_}}%


\section{简介}

为了分析和传达研究成果、论文结果，或者仅仅是某种结果，数据可视化往往
是必要且方便的。

\PGFPlots{} 是一个可视化工具。设计 \PGFPlots{} 的动机在于：作为最终用户，你
只需提供数据和描述作为输入，而 \PGFPlots{} 会处理其余一切，例如选择合适的
缩放因子、缩放到指定的目标尺寸、选择合适的显示范围、分配刻度位置、绘制
坐标轴并把描述文字放置在合适的位置。

\PGFPlots{} 解决了 \LaTeX{} 中一个长期存在的可视化问题：它的
描述文字使用与嵌入正文相同的字体，字号也完全相同。它直接嵌入 \LaTeX{}，
使得尽可能方便地使用 \LaTeX{} 强大的数学模式：从各种坐标轴描述一直到
用户自定义的标注。它提供文档级统一的线型、配色方案、标记\ldots{}
所有这些共同保证了一致性。

\PGFPlots{} 提供高质量的输出。同时，它是一种嵌入式解决方案：
它在很大程度上独立于第三方工具，不过也提供导入功能，以利用现有的工具。

它的主要目标是：你提供数据和描述——然后
\PGFPlots{} 无需更多输入即可运行。如果你愿意，也可以按自己的喜好进行
定制。


\input{pgfplots.tutorial1.tex}
\input{pgfplots.tutorial2.tex}
\input{pgfplots.tutorial3.tex}
\input{pgfplots.tutorial4.tex}
}%

% main=manual.tex

\chapter{参考手册}
\label{cha:pgfplots:reference}

{
\tikzset{external/figure name/.add={}{reference_}}

\input{pgfplots.reference.texdialects.tex}
\input{pgfplots.reference.axis-addplot.tex}
\input{pgfplots.reference.preliminaryoptions.tex}
\input{pgfplots.reference.2dplots.tex}
\input{pgfplots.reference.3dplots.tex}
\input{pgfplots.reference.markers-meta.tex}
\input{pgfplots.reference.axisdescription.tex}
\input{pgfplots.reference.scaling.tex}
\input{pgfplots.reference.3dconfiguration.tex}
\input{pgfplots.reference.errorbars.tex}
\input{pgfplots.reference.numberformatting.tex}
\input{pgfplots.reference.specifyrange.tex}
\input{pgfplots.reference.tickoptions.tex}
\input{pgfplots.reference.gridoptions-axiscoordinates.tex}
\input{pgfplots.reference.styleoptions.tex}
\input{pgfplots.reference.alignment.tex}
\input{pgfplots.reference.bb-clip.tex}
\input{pgfplots.reference.symbolic-transformations.tex}
\input{pgfplots.reference.coordfiltering.tex}
\input{pgfplots.reference.transformations.tex}
\input{pgfplots.reference.linefitting.tex}
\input{pgfplots.reference.miscellaneous.tex}
\input{pgfplots.reference.tikzinteroperability.tex}
\input{pgfplots.reference.layers.tex}
\input{pgfplots.reference.technicalinternals.tex}

}


\chapter{相关库}

本章介绍一些随 \PGFPlots{} 提供的库，但它们多少有些特殊，需要单独
激活。
\pgfmanualpdflabel{\textbackslash usepgfplotslibrary}{}

% \input{pgfplots.libs.statistics.tex}
% \endinput

\input{pgfplots.libs.clickable.tex}
\input{pgfplots.libs.colorbrewer.tex}
\input{pgfplots.libs.colormaps.tex}
%\input{pgfplots.libs.colortol.tex}         % not ready yet
\input{pgfplots.libs.dateplot.tex}
\input{pgfplots.libs.decorations.softclip.tex}
\input{pgfplots.libs.external.tex}
\input{pgfplots.libs.fillbetween.tex}
%\section{FIXME} \endinput
\input{pgfplots.libs.groupplots.tex}
\input{pgfplots.libs.patchplots.tex}
\input{pgfplots.libs.polar.tex}
\input{pgfplots.libs.smithchart.tex}
\input{pgfplots.libs.statistics.tex}
\input{pgfplots.libs.ternary.tex}
\input{pgfplots.libs.units.tex}


\chapter{内存与速度考量}

{
\tikzset{external/figure name/.add={}{memspeed_}}

\section{\TeX{} 的内存上限}
\label{sec:pgfplots:optimization}

只要正确配置 \TeX{} 的内存上限，\PGFPlots{} 就能排版包含数千个数据点的图。其运行时间大致与输入点数成正比。\footnote{事实上，运行时间是伪线性的：当点数大约达到 $100{,}000$ 时，会退化为二次增长。这一限制同样适用于 \PGF{} 路径的路径长度。此外，对堆叠图（stacked plots）而言，线性运行时间目前尚不可实现。}
%
\pgfplotsexpensiveexample
\begin{codeexample}[]
\begin{tikzpicture}
\begin{axis}[
    enlargelimits=0.01,
    title style={yshift=5pt},
    title=Scatter plot with $2250$ points,
]
    \addplot [
        blue,
        mark=*,
        only marks,
        mark options={scale=0.3},
    ] file [skip first]
        {plotdata/pgfplots_scatterdata3.dat};
\end{axis}
\end{tikzpicture}
\end{codeexample}

\pgfplotsexpensiveexample
\begin{codeexample}[]
\begin{tikzpicture}
\begin{axis}[
    enlarge x limits=0.03,
    title=Ornstein--Uhlenbeck sample
        ($13000$ time steps),
    xlabel=$t$,
]
    \addplot [blue] file {plotdata/ou.dat};
\end{axis}
\end{tikzpicture}
\end{codeexample}

\pgfplotsexpensiveexample
\begin{codeexample}[]
% huger graphs are possible; consider lualatex
\begin{tikzpicture}
\begin{axis}[
    title=$80 \times 80$ Smooth Surface,
    xlabel=$x$,
    ylabel=$y$,
]
  \addplot3 [surf,samples=80,shader=interp,domain=0:1]
        {sin(deg(8*pi*x))* exp(-20*(y-0.5)^2)
        + exp(-(x-0.5)^2*30
            - (y-0.25)^2 - (x-0.5)*(y-0.25))};
\end{axis}
\end{tikzpicture}
\end{codeexample}

\PGFPlots{} 完全依赖 \TeX{} 完成所有排版工作。它使用 \PGF{} 的前端层（front-end layer）和基础层（basic layer）来执行所有绘图操作。对于复杂的图，这可能需要花费一些时间，你可以参阅第~\ref{cha:pgfplots:importexport} 章，了解如何将单个图导出到外部图形文件。外部化（externalization）是减少排版时间的最佳方式。

然而，对于包含大量数据点的大规模图，很容易触及 \TeX{} 自身能力的上限。


\section{内存限制}

大多数 \TeX{} 发行版的默认设置都相当保守，因此可能需要调整它们。

通常，日志文件或最终的错误消息会包含一段关于已使用资源的摘要，从而提示需要增大哪一个参数。


\subsection{Lua\LaTeX{}}

一个效果相当不错的解决方案是切换 \LaTeX{} 可执行程序：如果你使用的是一个较完善的 \TeX{} 发行版，那么它应该已经附带了 |lualatex| 可执行程序。它采用动态内存分配，因此通常对任何 \PGFPlots{} 坐标轴都有足够的内存可用。

Lua\LaTeX{} 可执行程序 |lualatex| 与 |pdflatex| 几乎完全兼容。

这种方法适用于任何平台。


\subsection{MiK\TeX{}}

如果你正在使用 MiK\TeX{}，并且不想（或不能）切换到 |lualatex|，可以按如下步骤操作。

对于 MiK\TeX{}，可以通过两种方式增大内存上限。第一种是使用命令行开关：
%
\begin{codeexample}[code only]
pdflatex
    --stack-size=n --save-size=n
    --main-memory=n --extra-mem-top=n --extra-mem-bot=n
    --pool-size=n --max-strings=n
\end{codeexample}
%
\noindent 如果 MiK\TeX{} 出现内存不足，请调整这些设置。通常，我们不会手动调用 |pdflatex|：有一个开发辅助工具会负责所有调用，因此需要调整的是它。

有时，永久修改 MiK\TeX{} 配置文件可能更好，例如可以避免每次重新配置 \TeX{} 开发程序。这可以通过下面的命令实现
%
\begin{codeexample}[code only]
initexmf --edit-config-file=pdflatex
\end{codeexample}
%
\noindent 该命令既可以在 Windows 的命令提示符中输入，也可以通过「开始 $\gg$ 运行」执行。执行后，系统会用编辑器打开相应的配置文件。一个示例配置文件可以是
%
\begin{codeexample}[code only]
main_memory=90000000
save_size=80000
\end{codeexample}
%
或者也可以输入下面列出的任意一项配置。感谢 ``LeSpocky'' 提供的文档：

\url{http://blog.antiblau.de/2009/04/21/speicherlimits-von-miktex-erhoehen}.


\subsection{\TeX{}Live 或类似发行版}

除了切换到 |lualatex| 之外，你还可以按如下方式操作，以保留现有的 |dvips| 或 |pdflatex| 工作流。

对于 Unix 安装，需要调整配置文件。可以按如下步骤进行：
%
\begin{enumerate}
    \item 在你的系统上找到 |texmf.cnf|。在我的 Ubuntu 安装中，它位于

        |/usr/share/texmf/web2c/texmf.cnf|。
    \item 可以直接修改 |texmf.cnf|，也可以把它复制到某个方便的位置。如果你选择复制，步骤如下：
        %
        \begin{itemize}
            \item 在你本地的副本中只保留被修改的条目，以减少冲突。\TeX{} 总是会读取其搜索路径中找到的 \emph{所有} 配置文件。
            \item 调整搜索路径，使其能找到你的本地副本。这可以通过环境变量 |TEXMFCNF| 来完成。假设你的本地副本位于 |~/texmf/mytexcnf/texmf.cnf|，则可以写
                %
\begin{codeexample}[code only]
export TEXMFCNF=~/texmf/mytexcnf:
\end{codeexample}
                %
                这样就会先在你的目录中搜索，然后再搜索所有其他系统目录。
        \end{itemize}
    \item 你应当修改以下条目
        %
\begin{codeexample}[code only]
main_memory = n
extra_mem_top = n
extra_mem_bot = n
max_strings = n
param_size = n
save_size = n
stack_size = n
\end{codeexample}
        %
        日志文件通常会包含关于需要增大哪一个参数的信息。
\end{enumerate}
%
下面给出一个此类配置文件的示例。它会修改内存上限。
%
\begin{enumerate}
    \item 创建文件 |~/texmf/mytexcnf/texmf.cnf|（如必要，也一并创建相应的路径）。
        %
\begin{codeexample}[code only]
% newly created file ~/texmf/mytexcnf/texmf.cnf:
% If you want to change some of these sizes only for a certain TeX
% variant, the usual dot notation works, e.g.,
% main_memory.hugetex = 20000000
main_memory = 230000000 % words of inimemory available; also applies to inimf&mp
extra_mem_top = 10000000     % extra high memory for chars, tokens, etc.
extra_mem_bot = 10000000     % extra low memory for boxes, glue, breakpoints, etc.
save_size = 150000    % for saving values outside current group
stack_size = 150000    % simultaneous input sources

% Max number of characters in all strings, including all error messages,
% help texts, font names, control sequences.  These values apply to TeX and MP.
%pool_size = 1250000
% Minimum pool space after TeX/MP's own strings; must be at least
% 25000 less than pool_size, but doesn't need to be nearly that large.
%string_vacancies = 90000
% Maximum number of strings.
%max_strings = 100000
% min pool space left after loading .fmt
%pool_free = 47500
\end{codeexample}
        %
    \item 运行 |texhash|，让 \TeX{} 更新它的 |~/texmf/ls-R| 数据库。
    \item 创建环境变量 |TEXMFCNF|，并将其值设为 `|~/texmf/mytexcnf:|'（注意末尾的 `|:|'！）。在我的 Linux 系统上，可以通过向 |~/.bashrc| 添加
        %
\begin{codeexample}[code only]
export TEXMFCNF=~/texmf/mytexcnf:
\end{codeexample}
        %
        来实现。
\end{enumerate}

遗憾的是，\TeX{} 不允许任意设置内存上限：可执行程序中硬编码了一个上界。


\section{减少排版时间}

\PGFPlots{} 会进行大量计算，从抽象的坐标计算一直到由 \PGF{} 实现的底层 |.pdf| 绘图命令。对于复杂的图，这可能需要相当长的时间——尤其是 3D 图。


\subsection{LUA}

如果你使用 |compat=1.12|（或更新版本），并通过 |lualatex| 编译文档，\PGFPlots{} 就会启用其 |lua backend|。这一开关能减少生成输出文件所需的时间，对 3D 图尤其明显。

\begin{pgfplotskey}{lua backend=\mchoice{true,false} (initially true)}
    如果 |lua backend| 处于活动状态，并且文档是通过 |lualatex| 处理的，\PGFPlots{} 就会启用可扩展性与性能改进；它产生的输出与不启用时相同，但耗时更少，并且具有更智能的内存管理。

    该特性依赖于用快速脚本语言 \lua 对 \PGFPlots{} 的部分重新实现。要从中受益，你需要
    %
    \begin{enumerate}
        \item 在导言区写入 |compat=1.12|（或更新版本），并且
        \item 使用 |lualatex| 来编译你的 |.tex| 文件（或者至少编译那些包含 \PGFPlots{} 图的文件）。
    \end{enumerate}

    使用该特性后，生成曲面图所需的时间可以缩减到纯 \TeX{} 实现（即 |pdflatex|）所需时间的 25\%--50\%。其他类型的图同样会从中受益。

    |lua backend| 以两种方式工作：第一，用更快的 \lua 对应实现替换若干独立的子程序。这相对通用。第二，用等价的 \lua 实现替换整个处理步骤。%
    %
    \footnote{
        第一种方式的一个例子是：每当 \PGFPlots{} 在其 |colormap| 中进行查找时，它就把控制权转交给 |lua backend|，并将结果立即传回给 \TeX{}。主控制流仍然位于较慢的 \TeX{} 实现中。第二种方式的一个例子是 \texttt{\textbackslash addplot expression}：\PGFPlots{} 会把数学表达式以及所有相关的输入参数（包括 \texttt{samples} 和 \texttt{domain}）拷贝到 \texttt{lua backend}。然后，\texttt{lua backend} 会执行所有循环并收集坐标，最终把这些坐标交给 \pgfname{} 的（\TeX{}）实现。这样，控制流就位于快速的 \lua{} 解释器中。
    }
    \PGFPlots{} 1.12 为重要且耗时的操作提供了 \lua{} 实现。只要可能，就会使用这一实现。然而，它目前只覆盖了部分功能：\PGFPlots{} 的某些特性尚不被 \lua{} 后端支持。在这种情况下，它会自动关闭，并回退使用 \TeX{} 实现。

    |lua backend| 还会继续扩展，以获得更多性能提升并覆盖更多特性。如果你想查看 |lua backend| 是否在你的图中被（去）激活，请参考 |lua debug|。

    请注意，|lua backend| 使用了一个不同的、精度更高的数学引擎。因此，与 |pdflatex| 生成的输出相比，图的外观可能会出现一些无关紧要的像素差异。一般建议固定使用其中一种方式来生成文档，即要么使用 |lualatex|，\emph{要么}使用 |pdflatex|。由于 |lua backend| 的存在，从一种方式迁移到另一种方式可能会改变你文档的外观。

    最终，\PGFPlots{} 可能会加入一些计算开销较大的特性，这些特性只能作为 |lua backend| 的一部分实现，没有它将不可用。

    请注意，|lua backend| 需要一个至少支持 \lua~5.2 的 \TeX{} 发行版（例如 \TeX{}Live 2014）。如果你的 |lualatex| 附带的是更老版本的 \lua{} 解释器，它会自动停用。

    下面是一些如何从 |lua backend| 受益的指南：
    %
    \begin{itemize}
        \item 涉及宏定义的数学表达式在 |lua backend| 中不可用。最佳实践：优先使用 |declare function| 而不是宏常量：
            %
\begin{codeexample}[]
\begin{tikzpicture}
\begin{axis}[
    small,
    title=Do,
    samples=2,
]
    \addplot+ [
        declare function={C=4;}
    ] {C*x};
\end{axis}
\end{tikzpicture}
\end{codeexample}

\begin{codeexample}[]
\begin{tikzpicture}
\begin{axis}[
    small,
    title=Don't,
    samples=2,
]
        \def\constantC{4}
    \addplot {\constantC*x};
\end{axis}
\end{tikzpicture}
\end{codeexample}

            两者在语义上是等价的，但由于 \lua{} 无法解释 \TeX{} 宏，它会拒绝处理「Don't」这一情形。「Don't」情形会产生一条日志消息
            %
\begin{verbatim}
Package pgfplots info on input line 16: Deactivating LUA version of
plot expression for plot 0 (type 'pgfplothandlerlineto'):
y expression '\constantC *x' contains a TeX macro.
\end{verbatim}
            %
        \item 任何需要「原生」\TeX{} 代码的操作在 |lua backend| 中都不可用。这包括 |x filter/.code={...}|，因为 |/.code| 无法映射到 \lua。最佳实践：优先使用 |x filter/.expression| 而不是 |x filter/.code|：
            %
\begin{codeexample}[]
\begin{tikzpicture}
    \begin{axis}[small,grid=major,title=Do]
    \addplot+ [
        unbounded coords=jump,
        y filter/.expression={y==3 ? nan : y},
    ] table {
        x   y
        1   2
        1.5 2.5
        2   3
        2.5 3.5
        3   4
    };
    \end{axis}
\end{tikzpicture}
\end{codeexample}
            %
        \item 绘图表达式可以完全在 \lua 中处理；其他所有输入坐标类型仍由 \TeX{} 读取数值后再交给 \lua。
        \item 还有若干操作计划在未来版本中加入 |lua backend| 支持。
    \end{itemize}
\end{pgfplotskey}

\begin{pgfplotskeylist}{%
    lua debug,
    lua debug=\mchoice{false,off,off and silent,verbose,compileerror}%
}
    通常，|lua backend| 是静默工作的：\PGFPlots{} 自行决定它是否处于活动状态、是否支持当前操作，并据此继续执行。如果两个条件都满足，就会把控制权转交给 |lua backend|。如果其中一个不满足，它会自动回退到旧的 \TeX{} 实现。因此，你无法立即感知到它的存在，只有编译图片所需的时间会有所不同。

    该键用于控制：当 |lua backend| 处于活动状态、但由于遇到了不支持的特性而无法使用时，\PGFPlots{} 应如何处理。

    选项 \declaretext{false} 和 \declaretext{off} 会关闭 \lua{} 调试。在这种情况下，\PGFPlots{} \emph{仅}在使用了 \TeX{} 实现而不是 |lua backend| 时才写入日志消息。此外，它只会把消息写入 |.log| 文件（而不会写入控制台输出）。

    选项 \declaretext{off and silent} 既不写信息消息，也不写调试消息：它与 |off| 相同，只是在使用了 \TeX{} 实现时不再写日志消息。这是信息量最少的选项。

    选项 \declaretext{verbose} 还会把成功消息写入 |.log| 文件（不写入控制台）。

    选项 \declaretext{compileerror} 会在 |lua backend| 处于活动状态、但不支持当前图或绘图时中止编译。

    请注意，\PGFPlots{} 可能附带一些开销较大的受限操作，这些操作只能在 |lua backend| 中使用。如果由于某种原因无法使用 |lua backend|，这些条目将导致编译错误。
\end{pgfplotskeylist}


\subsection{每个图只编译一次}

减少排版时间的一种方法是：让 \PGF{} 在一次运行中为一部分（或全部）图形生成独立的临时 |.pdf|（或 |.eps|）文档，并在后续运行中复用这些临时图片。对于 \PGFPlots{} 而言，这是减少较大文档排版时间最有效的方式。可以通过第~\ref{sec:pgfplots:export} 节中描述的 |external| 库来实现。

}


\chapter{从其他格式导入/导出}
\label{cha:pgfplots:importexport}

{
\tikzset{external/figure name/.add={}{importexport_}}

本章介绍如何把单个图导出为独立的 \pdf{} 图形文件（或 \eps{} 图形文件）。此外，还介绍了一个 Matlab$^\text{\textregistered}$ 脚本，可用于把 Matlab 图转换为 \PGFPlots{}。


\section[导出为 PDF/EPS]{导出为 {\normalfont\pdf{}/\eps{}}}
\label{sec:pgfplots:export}

可以使用 \pgfname{} 和/或 \Tikz{} 的例程，把图导出为独立的 \pdf{} 文档。


\subsection{使用 \Tikz 的自动外部化框架}

\begin{pgfplotslibrary}{external}
\pgfkeys{/pdflinks/search key prefixes in/.add={/tikz/external/,}{}}
    |external| 库提供了一种便捷的方法，把每个独立的 |tikzpicture| 导出为单独的~|.pdf|（或~|.eps|）。后续的 \LaTeX{} 运行将直接包含这些图形，从而显著减少排版时间。

    该库也可以用于把文档提交给那些甚至没有安装 \PGFPlots{} 或 \Tikz{} 的作者。


    \paragraph{技术前言：}

    |external| 库由 Christian Feuersänger（\PGFPlots 的作者）编写。它已作为通用库贡献给 \Tikz{}，因此完整的参考文档以及所有调节选项可参见~\cite[Section ``Externalization Library'']{tikz}。命令 |\usepgfplotslibrary{external}| 实际上只是一个包装，它加载 |\usetikzlibrary{external}|；如果因为所安装的 \pgfname{} 版本至多为 $2.00$ 而该库尚不存在，则加载 \PGFPlots{} 自带的一个副本。

    |external| 库的设计使得对文档本身\emph{无需做任何改动}。其思路如下：
    %
    \begin{enumerate}
        \item 每一个 |\begin{tikzpicture}| $\dotsc$ |\end{tikzpicture}| 都会得到一个文件名。文件名既可以用 |\tikzsetnextfilename|\marg{output file name} 手动指定，也可以自动生成；自动生成时使用 \meta{tex file name}|-figure|\meta{number}，其中 \meta{number} 递增。

        \item 该库通过 \TeX{} 的 write18 系统调用，自动以形如 |pdflatex --jobname |\meta{output file name} 的命令写出生成的图片。这与调用 |gnuplot| 所用的是同一套框架。
    \end{enumerate}
    %
    唯一必要的步骤是使用

    \pgfmanualpdflabel{\textbackslash tikzexternalize}{}%
    |\usepgfplotslibrary{external}|

    |\tikzexternalize|

    \noindent 在文档导言区的某处（与系统相关的配置选项见下文）。文档无需再做其他修改。假设我们有一个名为 |test.tex| 的文件：
    %
\begin{codeexample}[code only]
\documentclass{article}
\usepackage{pgfplots}
    \usepgfplotslibrary{external}
    \tikzexternalize% activate externalization!
\begin{document}
    \begin{figure}
        \begin{tikzpicture}
            \begin{axis}
                \addplot {x^2};
            \end{axis}
        \end{tikzpicture}
        \caption{Our first external graphics example}
    \end{figure}

    \begin{figure}
        \begin{tikzpicture}
            \begin{axis}
                \addplot {x^3};
            \end{axis}
        \end{tikzpicture}
        \caption{A second graphics}
    \end{figure}
\end{document}
\end{codeexample}
    %
    \noindent 要启用系统调用，我们输入
    %
\begin{codeexample}[code only]
pdflatex -shell-escape test
\end{codeexample}
    %
    \noindent 之后 \LaTeX{} 会自动生成所需的图形文件 |test-figure0.pdf| 和 |test-figure1.pdf|。之后每次调用 |pdflatex| 都会直接使用 |\includegraphics|，不再重新处理 |tikzpicture| 本身（对于 MiK\TeX{}，需要不同的命令行开关，参见 |shell escape| 选项）。

    如果某个图需要重新生成，只需删除全部或部分图形文件并重新生成即可。另一种方式是在文档的某处使用 |\tikzset{external/force remake}|，之后每一个图都会自动重新生成。

    外部化图的文件名有三种修改方式：
    %
    \begin{itemize}
        \item 使用 |prefix| 修改整体文件名前缀，
        \item 使用 |\tikzsetnextfilename| 修改单个图的文件名，
        \item 使用 |figure name| 修改一组图的文件名。
    \end{itemize}

    \begin{key}{/tikz/external/prefix=\marg{file name prefix} (initially empty)}
        |\tikzsetexternalprefix|\marg{file name prefix} 的快捷方式，见下文。
    \end{key}

    \begin{command}{\tikzsetexternalprefix\marg{file name prefix}}
        为所有文件名指定一个公共前缀。例如，
        %
\begin{codeexample}[code only]
\tikzsetexternalprefix{figures/}
\end{codeexample}
        %
        会把 |figures/| 加到每一个外部图形文件名之前。
    \end{command}

    \begin{command}{\tikzsetnextfilename\marg{file name}}
        为\emph{下一个} \tikzname{} 图或 |\tikz| 短命令设置文件名。它\emph{只}对下一个图生效。

        没有显式指定文件名的图会自动生成文件名。

        请注意，|prefix| 仍然会被加到 \meta{file name} 之前。
        %
\begin{codeexample}[code only]
\documentclass{article}
% main document, called main.tex
\usepackage{tikz}
    \usepgfplotslibrary{external}
    \tikzexternalize[prefix=figures/]% activate with a name prefix
\begin{document}
\tikzsetnextfilename{firstplot}
\begin{tikzpicture} % will be written to 'figures/firstplot.pdf'
    \begin{axis}
        \addplot {x};
    \end{axis}
\end{tikzpicture}

\begin{tikzpicture} % will be written to 'figures/main-figure0.pdf'
   \draw [help lines] (0,0) grid (5,5);
\end{tikzpicture}
\end{document}
\end{codeexample}

\begin{codeexample}[code only]
pdflatex -shell-escape main
\end{codeexample}
        %
    \end{command}

    \begin{key}{/tikz/external/figure name=\marg{name}}
        与 |\tikzsetfigurename|\marg{name} 相同。
    \end{key}

    \begin{command}{\tikzsetfigurename\marg{name}}
        修改\emph{之后所有}图的名称。在文档编写过程中，可以通过
        |\tikzset{external/figure name|=\marg{name}|}| 来修改 |figure name|。每一个不同的 \meta{name} 都会使用一个独立的计数器，每个计数器都从
        $0$ 开始。\footnote{这些计数器存储在不同的\emph{宏}中。换句话说：不需要使用 \TeX{} 寄存器。}

        |prefix| 的值会在 |figure name| 求值之后再应用。
        %
\begin{codeexample}[code only]
\documentclass{article}
% main document, called main.tex
\usepackage{tikz}
    \usepgfplotslibrary{external}
    \tikzexternalize% activate externalization!
\begin{document}
% will be written to 'main-figure0.pdf'
\begin{tikzpicture}
    \begin{semilogyaxis}
        \addplot {exp(x)};
    \end{semilogyaxis}
\end{tikzpicture}

{
    \tikzset{external/figure name={subset_}}
    A simple image is \tikz \fill (0,0) circle(5pt);. % will be written to 'subset_0.pdf'

    \begin{tikzpicture} % will be written to 'subset_1.pdf'
        \begin{axis}
            \addplot {x^2};
        \end{axis}
    \end{tikzpicture}
}% here, the old file name will be restored:

\begin{tikzpicture} % will be written to 'main-figure1.pdf'
    \begin{axis}
        \addplot[domain=1e-3:100] {1/x};
    \end{axis}
\end{tikzpicture}
\end{document}
\end{codeexample}
        %
        |figure name| 的作用域在下一个右花括号处结束（与 |\tikzset| 设置的所有值一样）。

            \medbreak
        备注：使用
        |\tikzset{external/figure name/.add=|\marg{prefix}\marg{suffix}|}| 可以在 |figure name| 的实际值前添加一个 \meta{prefix}、在末尾追加一个 \meta{suffix}。这对于类似下面这样的场景可能有用：
        %
\begin{codeexample}[code only]
\tikzset{external/figure name=main}

% uses main_0.pdf, main_1.pdf, ...

\chapter{The first section}
{\tikzset{external/figure name/.add={}{_firstsection}}
    ...
    % uses main_firstsection_0.pdf, main_firstsection_1.pdf, ...
}

\chapter{The second section}
{\tikzset{external/figure name/.add={}{secondsection_}}
    ...
    % uses main_secondsection_0.pdf, main_secondsection_1.pdf, ...
    \section{Second subsection}
    {\tikzset{external/figure name/.add={}{sub_}}
        ...
        % uses main_secondsection_sub_0.pdf, main_secondsection_sub_1.pdf, ...
    }
    % uses main_secondsection_2.pdf, main_secondsection_3.pdf, ...
}
\end{codeexample}
    \end{command}

    \begin{command}{\tikzappendtofigurename\marg{suffix}}
        在 |figure name| 的实际值之后追加 \meta{suffix}。

        它是
        |\tikzset{external/figure name/.add={}|\marg{suffix}|}| 的快捷方式（即使没有安装 \tikzname{} 也支持此快捷方式，见下文）。
    \end{command}


    \paragraph{\textsc{eps} 输出或 MiK\TeX{} 的配置选项：}

    由于 |external| 库通过系统调用来工作，因此需要根据本地系统进行调整。对于 MiK\TeX{} 这是必要的，因为它使用不同的选项来启用这些系统调用。对于 \textsc{eps} 输出也是必要的，因为这涉及到一组不同的工具。

    请注意，\emph{最重要的部分}是启用系统调用。通常的做法是用 |pdflatex -shell-escape| 或 |pdflatex -enable-write18|（MiK\TeX{}）来排版你的文档。这些选项\emph{需要在你的 \TeX{} 编辑器中配置}。除这一步之外，你可能还希望配置系统调用：

    \begin{key}{/tikz/external/system call=\marg{template}}
    \label{extlib:systemcall:option}
        用于生成系统调用的模板字符串。在 \meta{template} 内部，宏 |\image| 可以作为占位符，表示即将生成的图片；而 |\texsource| 包含主文件名（实际上它包含 |\input|\marg{main file name}，但这无关紧要）。

        默认值为
        %
\begin{codeexample}[code only]
\tikzset{external/system call={pdflatex \tikzexternalcheckshellescape -halt-on-error
    -interaction=batchmode -jobname "\image" "\texsource"}
\end{codeexample}
        %
        \noindent 其中 \declareandlabel{\tikzexternalcheckshellescape} 仅在当前文档是用 |-shell-escape| 排版时，才会插入配置键 |shell escape| 的值。\footnote{请注意，对于默认配置而言这总是成立的。这一安全考虑主要针对 \texttt{mode=list and make}，它\emph{不}需要 shell escape 也能工作。}

        对于 |eps| 输出，你可以（并且需要）使用
        %
\begin{codeexample}[code only]
\tikzset{external/system call={latex \tikzexternalcheckshellescape -halt-on-error
    -interaction=batchmode -jobname "\image" "\texsource" &&
    dvips -o "\image".ps "\image".dvi}}
\end{codeexample}

        参数 \meta{template} 会用 |\edef| 展开，因此所有控制序列都会被展开。在这一求值过程中，`|\\|' 会变成一个普通反斜杠 `|\|'。此外，如果某些宏包把双引号 `|"|'、单引号 `|'|'、分号和短横线 `|-|' 用作宏，它们也会被转换为普通字符。例如，这保证了与 |german| 宏包的兼容性。
    \end{key}

    \begin{key}{/tikz/external/shell escape=\marg{command-line arg} (initially -shell-escape)}
        包含启用 |\write18| 特性的 |latex| 命令行选项。

        对于 \TeX{}Live，是 |-shell-escape|。对于 MiK\TeX{}，系统使用 |-enable-write18|。
    \end{key}


    \paragraph{对外部文件中标签与引用的支持}

    |external| 库对外部文件中的 |\label| 和 |\ref|（以及其他通常会在 |.aux| 文件中保存信息的命令）提供了额外的支持。

    然而，当你尝试使用下面两种情况时，仍有几点需要注意：
    %
    \begin{enumerate}
        \item[a)] 在外部化图形中引用主文档中的某个对象；或
        \item[b)] 在外部化图形中使用 |\label|，而该标签在主文档中被引用。
    \end{enumerate}

    对于情况 a)，外部化图形中的 |\ref| \emph{仅}在你\emph{手动}发出所需的系统调用、或通过 |make| 时才能工作。初始配置 |mode=convert with system call| \emph{不}支持 |\ref|。但你可以复制并粘贴由 |mode=convert with system call| 生成的系统调用，然后手动执行。原因在于 |\ref| 信息存储在主 |.aux| 文件中——但在调用 |mode=convert with system call| 时，该辅助文件还没有完全写完（存在竞争条件）。请注意，|\pageref| 不被支持（抱歉）。因此：如果外部图形中含有 |\ref|，请考虑使用 |mode=list and make|，或复制并粘贴该图的系统调用并手动执行。

    情况 b) 由外部库自动实现。具体来说，外部化图形中的 |\label| 会使外部库为每一张外部图生成单独的辅助文件。这些文件名为 \meta{imagename}|.dpth|。扩展名 |.dpth| 表示该文件还包含图片的深度（即 \tikzname 的 |baseline| 键）。此外，本应写入 |.aux| 文件的任何内容都会被重定向到 |.dpth| 文件——但仅限发生在外部化 |tikzpicture| 环境内部的内容。当主文档加载该图时，会把 |.dpth| 文件的内容拷贝到主 |.aux| 文件中。之后对主文档的后续编译就会包含外部的 |\label| 信息。换句话说，外部图形中的 |\label| 需要遵循下面的工作流程：
    %
    \begin{enumerate}
        \item 外部图形需要连同其 |.dpth| 一起生成（通常由 \tikzname 自动完成）。
        \item 主文档包含该外部图形，并把 |.dpth| 的内容拷贝到自己的主 |.aux| 文件中。
        \item 主文档需要再编译一次，以重新读取其 |.aux| 文件。\footnote{请注意，在 \LaTeX{} 中，无法在 \texttt{\textbackslash begin\{document\}} 之后激活辅助文件的内容。}
    \end{enumerate}
    %
    还有一种特殊情况：当 |\label|/|\ref| 是由 |tikzpicture| 绘制出来的时候。例如图例中的 |\ref| 图片，或 |\pgfplotslegendfromname| 特性。在这种情况下，你需要按照情况 a) 的方式处理，因为 |mode=convert with system call| 本身无法处理这些内容。

    换句话说：外部文档中的 |\label| 可以自动工作，只要把主文档多编译几次即可。|\ref| 则可能需要像上面情况 a) 那样手动调整。


    \paragraph{操作模式}

    \begin{key}{/tikz/external/mode=\mchoice{convert with system call,list and make,$\dotsc$} (initially convert with system call)}
        该键用于修改默认的操作模式。有若干种选项，全部详见
        ~\cite[Section ``Externalization Library'']{tikz}。其中最常用的大概是初始配置 |convert with system call| 以及专门选项 |list and make|。

        选项 |list and make| 配置该库去检查外部图形是否已经存在，并使用已存在的图形。如果图形不存在，该库会\emph{跳过}该图。但它同时还会生成一个 |makefile| 用于生成这些图形，以及一份所有必需图形文件的清单。

        不强制要求使用 |make|：该库只期待你以某种方式生成图片，并不关心「如何生成」。使用
        |make -f |\meta{name-of-tex-file}|.makefile -j 2| 可以并行执行，这确实可能是一个不错的选择。此外，makefile 还支持文件依赖：如果你的某个数据表已更新，外部图形会自动重新生成。\PGFPlots{} 会把任何文件依赖（输入文件和数据表）告知外部库。

        这两种模式具有如下特点：
        %
        \begin{enumerate}
            \item |convert with system call| 是全自动的，会即时完成所有工作。但它\emph{无法}处理外部图片中的 |\ref| 和/或 |\label| 信息。
            \item |list and make| 需要手动（手动发出系统调用）或半自动转换（使用生成的 \meta{main}|.makefile|），并多次运行 |pdflatex|。生成的 Makefile 可以并行处理。此外，|list and make| 对 |\ref| 和 |\label| 提供\emph{完整支持}：在外部化图形中定义的任何 |\label| 仍然可供主文档使用。

                如果你有使用 |legend to name| 的图例，或使用 |\label|/|\ref|，那么你需要先生成定义该 |\label|（或 |legend to name|）的图形，然后在主文档上运行 |pdflatex| 两次。之后就可以把该图例图形外部化了。
        \end{enumerate}
    \end{key}

    完整的参考文档以及其余选项记录在~\cite[``Externalization Library'']{tikz} 中。该参考还包含关于以下内容的信息：
    %
    \begin{itemize}
        \item 如何使用
            |\tikzset{external/|\declareandlabel{force remake}|}| 和
            |\tikzset{external/|\declareandlabel{remake next}|}| 来重新生成选定的图，
        \item 如何用
            |\tikzset{external/|\declareandlabel{export}|=false}| 部分禁用外部化，或用
            |\tikzexternaldisable| 完全禁用，
        \item 如何使用
            |\tikzset{external/optimize command away=\myExpensiveMacro}| 优化转换过程的速度，
        \item 如何用
            |\tikzpicturedependsonfile|\marg{name} 和/或
            |\tikzexternalfiledependsonfile|\marg{external file}\marg{name} 添加更多重新生成的依赖，
        \item 如何启用 |png| 导出的示例，
        \item 如何在未安装 \pgfname{} 的情况下排版此类文档，或
        \item 如何应对 |.pdf| 图片与边界框限制的变通方法。 \index{External Graphics!Bounding Box Issues}
            \index{Bounding Box Control!Image Externalization Problems}
    \end{itemize}


    \paragraph{在未安装 {\normalfont\pgfname} 或 {\normalfont\PGFPlots} 的情况下使用该库}

    有一个小型替换宏包 \declareandlabel{tikzexternal.sty}，可以在每个图都导出之后使用。思路是注释掉 |\usepackage{tikz}| 和 |\usepackage{pgfplots}|，改为写 |\usepackage{tikzexternal}|：
    %
\begin{codeexample}[code only]
% \usepackage{tikz}
% \usepackage{pgfplots}
\usepackage{tikzexternal}
    \tikzexternalize% activate externalization
\begin{document}
    \begin{tikzpicture}
        ...
    \end{tikzpicture}
...
\end{document}
\end{codeexample}
    %
    你不需要安装 \pgfname{}、\tikzname{} 或 \PGFPlots{}。你需要的是 |tikzexternal.sty| 以及所有生成的图（包括图片文件 `|.pdf|'，以及包含 |baseline| 选项信息的 `|.dpth|' 文件）。|tikzexternal.sty| 随 \pgfname{} 一起发布，位于目录
    %
\begin{codeexample}[code only]
latex/pgf/utilities/tikzexternal.sty
\end{codeexample}
    %
    同时 \PGFPlots{} 也附带了一份副本：
    %
\begin{codeexample}[code only]
tex/generic/pgfplots/oldpgfcompatib/pgfplotsoldpgfsupp_tikzexternal.sty
\end{codeexample}
    %
    只需把该文件拷贝到你的目录下，并将其重命名为 |tikzexternal.sty|。


    \paragraph{注意：}

    这个小型替换宏包不支持键值接口。因此，必须用 |\tikzsetexternalprefix| 代替 |prefix| 选项，用 |\tikzsetfigurename| 代替 |figure name| 选项，因为在这种上下文中 |\tikzset| 不可用。另外，如果你之前用过 |\pgfplotsset|，可能还需要定义一个空的宏 |\pgfplotsset|。
\end{pgfplotslibrary}


\subsection[手动使用 PGF 外部化框架]{手动使用 {\normalfont\pgfname} 的外部化框架}

把 \TeX{} 图导出为独立图形文件的另一种方法是使用 \pgfname{} 的外部化框架。它需要更多工作，但比 |external| 库更通用。基本思路是用如下命令把希望导出的部分封装起来：

\declareandlabel{\beginpgfgraphicnamed}\marg{output file name}

\meta{picture contents}

\declareandlabel{\endpgfgraphicnamed}.

\noindent 此外，还需要在导言区用如下命令告知 \pgfname{} 主文档的名称：

\declareandlabel{\pgfrealjobname}\marg{the real job's name}

\noindent 在导言区。这会启用两种不同的模式：
%
\begin{enumerate}
    \item 第一种是正常排版模式。\LaTeX{} 会检查是否存在名为 \meta{output file name}、且带有某种可接受扩展名的文件——如果存在，则用 |\pgfimage| 包含该图形文件，并跳过 \meta{picture contents}。如果该文件不存在，则正常排版 \meta{picture contents}。当 |\jobname| 等于 \meta{the real job's name} 时，使用此模式。
    \item 第二种模式在 |\jobname| 等于 \meta{output file name} 时启用，它进入「转换模式」，用于写出图形文件 \meta{output file name}。此时，\emph{只}把 \meta{picture contents} 写入 |\jobname|，\LaTeX{} 的其余部分照常处理，但会被静默丢弃。

        对于每一个外部化图形文件，都需要手动用
        |pdflatex --jobname |\meta{output file name} 来启动此模式。
\end{enumerate}
%
下面是一个完整的示例。
%
\begin{codeexample}[code only]
\documentclass{article}
\usepackage{pgfplots}
    \pgfrealjobname{test}
\begin{document}
    \begin{figure}
        \beginpgfgraphicnamed{testfigure}
            \begin{tikzpicture}
                \begin{axis}
                    \addplot {x^2};
                \end{axis}
            \end{tikzpicture}
        \endpgfgraphicnamed
        \caption{Our first external graphics example}
    \end{figure}

    \begin{figure}
        \beginpgfgraphicnamed{testfigure2}
            \begin{tikzpicture}
                \begin{axis}
                    \addplot {x^3};
                \end{axis}
            \end{tikzpicture}
        \endpgfgraphicnamed
        \caption{A second graphics}
    \end{figure}
\end{document}
\end{codeexample}
%
\noindent 该文件名为 |test.tex|，我们用如下命令处理它（例如）：
%
\begin{codeexample}[code only]
pdflatex test
\end{codeexample}
%
\noindent 现在我们输入
%
\begin{codeexample}[code only]
pdflatex --jobname testfigure test
pdflatex --jobname testfigure2 test
\end{codeexample}
%
\noindent 进入转换模式。最后这两次调用\emph{只}会把我们命名的图形环境内容写入各自的输出文件 |testfigure.pdf| 和 |testfigure2.pdf|：一个对应 \meta{testfigure}，另一个对应 \meta{testfigure2}。

总之，需要 |\pgfrealjobname|，并对每一个外部化图形环境调用
|pdflatex --jobname |\meta{graphics file}。请注意，必须使用上面给出的语法，\emph{不要}使用 |\begin{pgfgraphicnamed}|。

这些步骤在~\cite{tikz} 的 Section``Externalizing Graphics'' 中有更详细的说明。


\paragraph{注意：}

不要忘记正确的 |\pgfrealjobname| 语句！如果缺失，外部化将根本无法工作。如果写错，每次调用 \LaTeX{} 都会生成空的输出文件。

需要注意的是，这种图片外部化方法并不局限于 \Tikz{} 图环境。事实上，它会把起始与结束语句之间的所有内容都收集到外部文件中。这里隐含的假设是被封装的内容是一个盒子，但你也可以用类似 \LaTeX{} minipage 的方式（或使用 |\vbox|，它功能没那么强大，但对文档其余部分影响也更小）封装完整的段落。

\begin{key}{/pgf/images/aux in dpth=\mchoice{true,false} (initially false)}
    如果该布尔键被设为 |true|，则在外部图片内部生成的任何 |\label| 信息都会存储到前面提到的 |.dpth| 文件中。这样主文档就可以引用文档中被外部化部分的标签信息（尽管你可能需要多次运行 |latex|）。

    对 |\ref| 提供了标签支持，可能也包括 |\cite|。|\pageref| 命令只被部分支持。
\end{key}


\paragraph{在未安装 {\normalfont\pgfname} 的情况下使用该库}

只需把宏包 |\usepackage{tikz}| 和 |\usepackage{pgfplots}| 取消注释，改为使用
%
\begin{codeexample}[code only]
\long\def\beginpgfgraphicnamed#1#2\endpgfgraphicnamed{%
    \begingroup
    \setbox1=\hbox{\includegraphics{#1}}%
    \openin1=#1.dpth
    \ifeof1 \box1
    \else
        \read1 to\pgfincludeexternalgraphicsdp \closein1
        \dimen0=\pgfincludeexternalgraphicsdp\relax
        \hbox{\lower\dimen0 \box1 }%
    \fi
    \endgroup
}
\end{codeexample}
%
代替。这将包含生成的图形文件（并且会遵循 |.dpth| 文件中存储的 |baseline| 信息）。这样一来，你就不需要安装 \pgfname{} 或 \PGFPlots{}。详见~\cite{tikz} 的 Section``Externalizing Graphics''。


\section{从 Matlab 导入}

\subsection{从 Matlab 导入网格数据到 PGFPlots}
\label{sec:pgfplots:import:matlab}

虽然把 Matlab 向量写入文件很容易（使用 |save P.dat data -ASCII|），但导出网格数据就要复杂得多。

主要问题在于如何把网格结构告诉 \PGFPlots{}。

下面是实现这一任务的一个例子：在 Matlab 中，我们有网格数据 |X|、|Y| 和 |Z|，它们是大小相同的矩阵。例如，假设我们有
%
\begin{codeexample}[code only]
[X,Y] = meshgrid( linspace(-1,1,5), linspace(4,5,10) );
Z = X + Y;
surf(X,Y,Z)
\end{codeexample}
%
\noindent 作为数据。然后，我们可以生成一个 $N \times 3$ 的表，按列优先顺序包含所有单独的元素：
%
\begin{codeexample}[code only]
data = [ X(:) Y(:) Z(:) ]
% or -ascii
save P.dat data -ASCII
size(X)
ans =
   10.00    5.00
\end{codeexample}
%
\noindent 其中第二条命令把 $N \times 3$ 的表存入 |P.dat|。最后，我们可以在 \PGFPlots{} 中使用

|\addplot3[surf,mesh/rows=10,mesh/ordering=colwise,shader=interp] file {P.dat};|

来读取这些数据。我们需要提供行数（这里是 $10$）或列数——以及排序方式（对于 Matlab 矩阵是 |colwise|）。

在 \PGFPlots{} 中更快的另一种做法是：在 Matlab 中转置矩阵，并告诉 \PGFPlots{} 它们是按 |rowwise| 顺序排列的。这样，最后一步就变为
%
\begin{codeexample}[code only]
XX=X'; YY=Y'; ZZ=Z';
data = [ XX(:) YY(:) ZZ(:) ]
save P.dat data -ASCII
\end{codeexample}
%
\noindent 对应的 \PGFPlots{} 命令为

|\addplot3[surf,mesh/cols=10,mesh/ordering=rowwise,shader=interp] file {P.dat};|.


\subsection{matlab2pgfplots.m}

这是一个 Matlab$^\text{\textregistered}$ 脚本，尝试把 Matlab 图转换为 \PGFPlots{}。它要求 Matlab 版本 7.4（或更高）。


\paragraph{注意：}

该脚本已经严重过时，只支持 \PGFPlots{} 非常小的一个子集。你可以看看 |matlab2tikz|，这是 Nico Schlömer 编写的转换脚本，地址为

\url{http://www.mathworks.com/matlabcentral/fileexchange/22022-matlab2tikz}

\noindent 它同样使用 \PGFPlots{} 进行 \LaTeX{} 转换。

\medskip
|matlab2pgfplots.m| 的思路是
%
\begin{itemize}
    \item 以完整的 Matlab 图作为输入，
    \item 获取坐标轴标签、坐标轴缩放（对数或普通）以及图例条目，
    \item 获取所有绘图坐标
\end{itemize}
%
并写出一个等价的 \texttt{.pgf} 文件，用 \PGFPlots{} 来排版该图。

其目的\emph{不是}模拟 Matlab。它只是转换工作的第一步。在 Matlab 命令行输入
%
\begin{lstlisting}
> help matlab2pgfplots
\end{lstlisting}
%
在 Matlab 提示符下，可获得关于其特性与局限的更多信息。

该脚本仍处于实验阶段。


\subsection{matlab2pgfplots.sh}

一个 \texttt{bash} 脚本，简单地启动 Matlab 并运行
%
\begin{lstlisting}
    f=hgload( 'somefigure.fig' );
    matlab2pgfplots( 'outputfile.pgf', 'fig', f );
\end{lstlisting}
%
参见上面的 matlab2pgfplots.m。


\subsection{从 Matlab 导入颜色映射表}

有时你可能希望在 \PGFPlots{} 中复用 Matlab 的 |colormap|。下面是一个把 Matlab 颜色映射表转换为 \PGFPlots{} 的小脚本：
%
\begin{codeexample}[code only]
C = colormap;  % gets data of the current colormap.
% C = colormap(jet) % gets data of "jet"
eachnth = 1;
I = 1:eachnth:size(C,1); % this is nonsense for eachnth=1 -- but perhaps you don't want each color.
CC = C(I,:);
TeXstring = [ ...
 sprintf('\\pgfplotsset{\n\tcolormap={matlab}{\n') ...
 sprintf('\t\trgb=(%f,%f,%f)\n',CC') ...
 sprintf('\t}\n}\n') ]
\end{codeexample}


\section{SVG 输出}

可以把每一个 \Tikz{} 图都写入一个可缩放矢量图形（\texttt{.svg}）文件。这与 \PGFPlots{} 无关，它是 \PGF{} 的一个独立驱动。请参见~\cite[Section ``Producing HTML / SVG
Output'']{tikz}。


\section{在 Python 中生成 \PGFPlots{} 图形}

Mario Orne DÍAZ ANADÓN 贡献了一个小型 python 脚本 |pgfplots.py|，它提供了一个简单的接口，可以在 python 中生成 \PGFPlots{} 图。该脚本位于 \PGFPlots{} 安装目录下的
|pgfplots/scripts/pgfplots/pgfplots.py|；其文档可在该文件中找到。

}


\chapter{工具与基础层命令}
\label{cha:pgfplots:lowlevel}

本章记录的命令提供了对 \PGFPlots{} 更基础元素的访问。其中大多数与 \pgfname{} 的基础层密切相关，尤其是各种与坐标轴相关的点命令。还有一些是通用工具，例如循环。

不过，本节中的大多数元素只对高级用户有意义——而且可能只在特殊情况下才用得到。


\section{工具命令}

\begin{command}{\foreach \meta{variables} |in| \meta{list} \marg{commands}}
    \Tikz{} 提供的一个功能强大的循环命令，参见~\cite[Section
    ``Utilities'']{tikz}。
    %
\begin{codeexample}[]
\foreach \x in {1,2,...,4} {Iterating \x. }%
\end{codeexample}

    一个与 \PGFPlots{} 相关的例子可以是
    %
\begin{codeexample}[code only]
\foreach \i in {1,2,...,10} {\addplot table {datafile\i}; }%
\end{codeexample}
\end{command}

\begin{command}{\pgfplotsforeachungrouped \meta{variable} |in| \meta{list} \marg{command}}
    |\foreach| 的一个特殊变体，它有两点不同：在执行 \meta{command} 时不引入额外的分组；并且允许对（简单的！）
    \meta{$x_0$}|,|\meta{$x_1$}|,...,|\meta{$x_n$} 表达式调用数学解析器。
    %
\begin{codeexample}[]
\def\allcollected{}
\pgfplotsforeachungrouped \x in {1,2,...,4} {Iterating \x. \edef\allcollected{\allcollected, \x}}%
All collected = \allcollected.
\end{codeexample}

    一个更有用的例子是与数据表配合使用。下面的例子取自 \PGFPlotstable{}：
    %
\begin{codeexample}[code only]
\pgfplotsforeachungrouped \i in {1,2,...,10} {%
    \pgfplotstablevertcat{\output}{datafile\i} % appends `datafile\i' -> `\output'
}%
% since it was ungrouped, \output is still defined (would not work
% with \foreach)
\end{codeexample}

    \paragraph{备注：}

    特殊语法
    \meta{list}=\meta{$x_0$}|,|\meta{$x_1$}|,...,|\meta{$x_n$}，即以两个元素开头、中间为省略号、末尾为一个元素，会为循环调用数学解析器。因此，当 |/pgf/fpu| 处于活动状态时，它支持的数值范围比其他任何语法都更大。在其他情况下，
    |\pgfplotsforeachungrouped| 会调用 |\foreach|，并在不使用 \TeX{} 分组的情况下提供结果。

    请记住：在坐标轴环境内，所有循环构造（包括自定义循环、|\foreach| 和 |\pgfplotsforeachungrouped|）都需要小心处理：循环变量只能在被立即求值的地方使用；但 \PGFPlots{} 会推迟许多宏的求值。例如，要遍历某些内容并生成形如 |\node at (axis cs:\i,0.5)...| 的坐标轴描述，循环宏 |\i| 会在 |\end{axis}| 处才被求值——但那时循环已经结束，它的值已经丢失。正确的做法是\emph{显式展开}循环变量。例如：
    %
\begin{codeexample}[code only]
\pgfplotsforeachungrouped \i/\j in {
    1 / a,
    2 / b,
    3 / c
}{
    \edef\temp{\noexpand\node at (axis cs: \i,0.5) {\j};}
    % \show\temp % lets TeX show you what \temp contains
    \temp
}
\end{codeexample}
    %
    该示例生成三次循环迭代：|\i=1|、|\j=a|；然后 |\i=2|、|j=b|；然后 |\i=3|、|\j=c|。在循环体内部，它把它们展开，并用「展开定义」|\edef| 把结果赋给一个宏。结果中不再包含 |\i| 或 |\j|（因为它们已经被展开）。然后，调用得到的宏。关于 \TeX{} 命令 |\edef| 以及展开控制的详细信息，请参见随 \PGFPlots{} 一起提供的文档
    \href{file:TeX-programming-notes.pdf}{TeX-programming-notes.pdf}。
\end{command}

\begin{command}{\pgfplotsinvokeforeach\marg{list} \marg{command}}
    |\pgfplotsforeachungrouped| 的一个变体（因而也是 |\foreach| 的变体），它会对 \meta{list} 中的每一个元素，把 \meta{command} 内出现的 |#1| 各替换一次。因此，它实际上假设 \marg{command} 像一个 |\newcommand| 的命令体。

    换句话说，对 \meta{list} 的每一个元素都会调用一次 \meta{command}。\meta{list} 的当前元素以 |#1| 的形式可用。

    与 |\pgfplotsforeachungrouped| 一样，本命令\emph{不}引入额外的作用域（即它也是不分组的）。

    它与 |\foreach \x in |\meta{list}\marg{command} 的区别很细微：|\x| \emph{不会}被展开，而 |#1| 会。
    %
\begin{codeexample}[]
\pgfkeys{
    otherstyle a/.code={[a]},
    otherstyle b/.code={[b]},
    otherstyle c/.code={[c]},
    otherstyle d/.code={[d]}}
\pgfplotsinvokeforeach{a,b,c,d}
    {\pgfkeys{key #1/.style={otherstyle #1}}}
Invoke them:
\pgfkeys{key a} \pgfkeys{key b}
\pgfkeys{key c} \pgfkeys{key d}
\end{codeexample}
相反的例子则使用一个宏（这里是 |\x|）作为循环参数：
\begin{codeexample}[]
\pgfkeys{
    otherstyle a/.code={[a]},
    otherstyle b/.code={[b]},
    otherstyle c/.code={[c]},
    otherstyle d/.code={[d]}}
\pgfplotsforeachungrouped \x in {a,b,c,d}
    {\pgfkeys{key \x/.style={otherstyle \x}}}
Invoke them:
\pgfkeys{key a} \pgfkeys{key b}
\pgfkeys{key c} \pgfkeys{key d}
\end{codeexample}

    \paragraph{限制：}

    目前还不能嵌套使用此命令（因为它不通过作用域提供保护）。
\end{command}

\begin{command}{\pgfmathparse\marg{expression}}
    对 \meta{expression} 调用 \pgfname{} 的数学解析器，并把
    \declareandlabel{\pgfmathresult} 定义为结果。
    %
\begin{codeexample}[]
\pgfmathparse{1+41}

The result is `\pgfmathresult'.
\end{codeexample}
    %
    \noindent \pgfname{} 中的数学引擎通常使用 \TeX{} 内部的算术。也就是说：它适合处理范围在
    $[-16384,16384]$ 内的数字，精度为 $5$ 位。

    对于绘图应用而言，这个数值范围通常太小。\PGFPlots{} 通过
    |\pgfkeys{/pgf/fpu}\pgfmathparse{1+41}| 启用「浮点单元」（fpu），从而扩展数值范围，并让之后所有运算都以浮点数进行。

    在 \PGFPlots{} 中，键 |/pgfplots/use fpu| 默认是打开的，这意味着所有坐标运算都用 |fpu| 进行。然而，所有与 \pgfname{} 相关的绘图操作仍然使用标准数学引擎。

    如果你曾经需要处理扩展精度的数字，可以使用
    %
\begin{codeexample}[]
\pgfkeys{/pgf/fpu}%
\pgfmathparse{1000*1000}

The result is `\pgfmathprintnumber{\pgfmathresult}'.
\end{codeexample}
    %
    请注意，|fpu| 的结果通常不是人类可读格式，因此 |\pgfmathprintnumber| 是排版这类数字的首选方式。

    更多细节请参见 \cite{tikz}。
\end{command}

\begin{pgfplotskey}{use fpu=\mchoice{true,false} (initially true)}
    \PGFPlots{} 提供了多种计算数学表达式的方式，其中 |use fpu| 是最强大的。它的数学运算要么在 |lua backend| 中实现，要么在纯 \TeX{} 实现中完成，数值范围大、精度也足够。

    然而，存储在 |\pgfmathresult| 中的值是晦涩难懂的，需要通过专门的宏来处理。|use fpu| 这一开关只在该数值格式导致问题时才有用——也就是说，它是一个调试开关，在正常操作中不应使用。
\end{pgfplotskey}

\begin{key}{/pgf/declare function=\meta{function definitions}}
    允许定义一个或多个函数。

    参数 \meta {function definitions} 可以包含一个或多个定义，每一个\emph{必须}以分号结尾：
    %
\begin{codeexample}[]
\begin{tikzpicture}
\begin{axis}[
    declare function={
        C=4;
        square(\t)=(\t)^2 + C;
    },
]
    \addplot+ [samples=2] {C*x};

    \addplot {square(x)};
\end{axis}
\end{tikzpicture}
\end{codeexample}
    %
    这些定义本身的形式为 \meta{function}\meta{argument list} =
    \meta{definition}，其中 \meta{argument list} 是用逗号分隔的参数列表，例如 |\t| 或 |\t,\a,\b|。\meta{definition} 是一个数学表达式，它使用这些参数。

    更多细节请参见 \cite{tikz}。
\end{key}

\begin{command}{\pgfplotstableread\marg{file}}
    请参见 \PGFPlotstable{} 的手册 |pgfplotstable.pdf|，它是 \PGFPlots{} 套件的一部分。
\end{command}

\begin{command}{\pgfplotstabletypeset\marg{\textbackslash macro}}
    请参见 \PGFPlotstable{} 的手册 |pgfplotstable.pdf|，它是 \PGFPlots{} 套件的一部分。
\end{command}

\begin{command}{\pgfplotsiffileexists\marg{filename}\marg{true code}\marg{false code}}
    如果 \meta{filename} 存在则调用 \meta{true code}，否则调用 \meta{false code}。例如，可以在循环宏中使用它，依次绘制每一个数据文件，直到没有更多文件为止。
\end{command}

\begin{command}{\pgfplotsutilifstringequal\marg{first}\marg{second}\marg{true code}\marg{false code}}
    一个简单的「字符串比较」工具：当 \meta{first}
    $=$\meta{second} 时调用 \meta{true code}，否则调用 \meta{false code}。它不展开宏。
\end{command}

\begin{commandlist}{\pgfkeys,\pgfeov,\pgfkeysvalueof,\pgfkeysgetvalue}
    这些命令是 \Tikz{} 指定选项方式的一部分，属于其子宏包 |pgfkeys|。命令 |\pgfplotsset| 实际上只是对 |\pgfkeys| 的一层包装。

    关于 |\pgfkeys| 的简短介绍可参见~\cite{keyvalintro}，完整的参考当然是 \Tikz{} 手册~\cite{tikz}。

    键 |\pgfkeysvalueof|\marg{key name} 展开为某个键的值；
    |\pgfkeysgetvalue|\marg{key name}\marg{\textbackslash macro} 把 \meta{key name} 的值存入 \meta{\textbackslash macro}。宏 |\pgfeov| 用于在 |\pgfkeys| 中为代码键分隔参数，请参见上面提到的参考。
\end{commandlist}


\section[PGFPlots 坐标轴内部的命令]
        {{\normalfont\PGFPlots{}} 坐标轴内部的命令}

\begin{command}{\autoplotspeclist}
    此命令不应再使用，但会作为技术实现细节保留下来。请使用 `|cycle list|' 选项，见第~\ref{sec:cycle:list} 节。
\end{command}

\begin{command}{\logten}
    展开为常数 $\log(10)$。对于对数图很有用，因为 $\log(10^i)
    = i\log(10)$。此命令仅在 \Tikz{} 图内部可用。
\end{command}

\begin{command}{\pgfmathprintnumber\marg{number}}
    为 \meta{number} 生成排版精美的输出\footnote{此方法以前叫
    \texttt{\textbackslash prettyprintnumber}。它的功能已并入 \PGF{}，旧命令现已弃用。}。每个刻度标签都使用此方法。

    数字按当前数字打印选项输出；关于不同的数字样式、舍入精度与舍入方法，请参见随本宏包提供的 \PGFPlotstable{} 手册。
\end{command}

\begin{command}{\numplots}
    在任一坐标轴环境、相关样式、选项或命令内部，|\numplots| 展开为图的总数。
\end{command}

\begin{command}{\numplotsofactualtype}
    与 |\numplots| 类似，但此宏返回具有相同绘图处理器的图的总数。因此，如果启用了 |sharp plot|，它返回所有 |sharp plots| 的数量；如果启用了 |ybar|，它返回 |ybar| 图的数量，依此类推。
\end{command}

\begin{command}{\plotnum}
    在 |\addplot| 或任何相关样式、选项或命令内部，|\plotnum| 展开为当前图的编号，从~$0$ 开始。
\end{command}

\begin{command}{\plotnumofactualtype}
    与 |\plotnum| 类似，但它返回在所有同类型图中的编号。所有此类图的总数可用 |\numplotsofactualtype| 获取。
\end{command}

\begin{command}{\coordindex}
    在 |\addplot| 命令内部，此宏展开为当前坐标的编号（从~$0$ 开始）。

    它与 |x filter| 或 |y filter| 配合使用，用于（反）选坐标。
\end{command}


\section{路径操作}

\begin{commandlist}{\path,\draw,\fill,\node,\matrix}
    这些都是 \Tikz{} 绘图命令，全部记录在~\cite{tikz} 中。它们用于绘制或填充路径、生成文本节点或对齐的文本矩阵。它们分别等价于
    \pgfmanualpdflabel{/tikz/draw}{}|\path[draw]|、
    \pgfmanualpdflabel{/tikz/fill}{}|\path[fill]|、
    \pgfmanualpdflabel{/tikz/node}{}|\path[node]|、
    \pgfmanualpdflabel{/tikz/matrix}{}|\path[matrix]|。
\end{commandlist}

\begin{pathoperation}{--}{\meta{coordinate}}
    一个 \Tikz{} 路径操作，用直线连接当前点（即 |--| 之前的最后一个点）与 \meta{coordinate}。
\end{pathoperation}

{\catcode`\|=12
\begin{pathoperation}[noindex]{|-}{\meta{coordinate}}
\pgfmanualpdflabel[\catcode`\|=12 ]{|-}{}%
    一个 \Tikz{} 路径操作，用\emph{两}段直线连接当前点与
    \meta{coordinate}：先竖直，后水平。
\end{pathoperation}

\begin{pathoperation}[noindex]{-|}{\meta{coordinate}}
\pgfmanualpdflabel[\catcode`\|=12 ]{-|}{}%
    一个 \Tikz{} 路径操作，用\emph{两}段直线连接当前点与
    \meta{coordinate}：先水平，后竖直。
\end{pathoperation}
}

\begin{keylist}{/tikz/xshift=\marg{dimension},/tikz/yshift=\marg{dimension}}
    这些 \Tikz{} 键允许把某个对象沿 \meta{dimension} 平移，\meta{dimension} 可以是任何 \TeX{} 长度（或表达式）。
\end{keylist}

\begin{command}{\pgfplotsextra\marg{low-level path commands}}
    在 \PGFPlots{} 坐标轴中执行 \meta{low-level path commands} 的命令。由于坐标轴内部的任何绘图命令都必须推迟到坐标轴完成、缩放初始化之后才能执行，所以不能简单地直接绘制路径。取而代之，必须在坐标轴完成时再绘制它们。对每一个 \Tikz{} 路径这会自动完成——如果你写 |\pgfplotsextra|\marg{commands}，也是手动触发这一过程。
    %
\begin{codeexample}[]
\begin{tikzpicture}
\begin{axis}[xmin=0,xmax=3,ymin=0,ymax=5]
    \pgfplotsextra{
        \pgfpathmoveto{\pgfplotspointaxisxy{1}{2}}
        \pgfpathlineto{\pgfplotspointaxisxy{2}{4}}
        \pgfusepath{stroke}
    }
\end{axis}
\end{tikzpicture}
\end{codeexample}
    %
    上面的例子初始化了一个坐标轴，并在坐标轴就绪后执行基础层路径命令。多个 |\path|、|\addplot| 和 |\pgfplotsextra| 命令的执行顺序与它们在环境中出现的顺序相同。\footnote{堆叠图例外，其顺序可能相反，参见键 \texttt{reverse stack plots}。}
\end{command}

\begin{command}{\pgfplotspathaxisoutline}
    生成一条与当前坐标轴外轮廓相似的路径。该路径用于裁剪路径和背景路径（如果存在）。
\end{command}


\section{指定基础坐标}
\label{sec:basic:coordinates}

\begin{commandlist}{%
    \pgfplotspointaxisxy\marg{x coordinate}\marg{y coordinate},
    \pgfplotspointaxisxyz\marg{x coordinate}\marg{y coordinate}\marg{z coordinate}%
}
    类似 |\pgfpointxy| 的点命令，接受逻辑的绝对坐标，返回一个底层点。所有从用户变换到对数的变换都会被应用。

    由于变换是在坐标轴完成后才初始化的，因此此命令需要被推迟（参见 |\pgfplotsextra|）。

    此命令是基础层版本的 |axis cs:|\meta{x
    coordinate}|,|\meta{y coordinate}|,|\meta{z coordinate}。

    请注意，这也是可视化阶段的默认坐标系；换句话说：如果你写 |\draw (1,2) -- (1,4)|，\PGFPlots{} 会自动使用 |(axis cs:1,2) -- (axis cs:1,4)|。
\end{commandlist}

\begin{commandlist}{%
    \pgfplotspointaxisdirectionxy\marg{x coordinate}\marg{y coordinate},
    \pgfplotspointaxisdirectionxyz\marg{x coordinate}\marg{y coordinate}\marg{z coordinate}%
}
    类似 |\pgfpointxy| 的点命令，接受逻辑\emph{相对}坐标，返回一个底层点。所有从用户变换到对数的变换都会被应用。与
    |\pgfplotspointaxisxy| 的区别在于：这里会跳过线性变换的平移部分（对比 |disabledatascaling|）。

    此命令是基础层版本的 |axis direction cs:|\meta{x
    coordinate}|,|\meta{y coordinate}|,|\meta{z coordinate}。更多细节请参见 |axis direction cs| 的文档。

    当需要提供具有\emph{相对}性质的量（如方向或长度）时，使用此命令。一个用例是绘制椭圆：
    %
\begin{codeexample}[]
\begin{tikzpicture}
\begin{axis}[
    xmin=-3,   xmax=3,
    ymin=-3,   ymax=3,
    extra x ticks={-1,1},
    extra y ticks={-2,2},
    extra tick style={grid=major},
]
    \draw [red] \pgfextra{
        \pgfpathellipse{\pgfplotspointaxisxy{0}{0}}
            {\pgfplotspointaxisdirectionxy{1}{0}}
            {\pgfplotspointaxisdirectionxy{0}{2}}
    % see also the documentation of
    % 'axis direction cs' which
    % allows a simpler way to draw this ellipse
    };
    \draw [blue] \pgfextra{
        \pgfpathellipse{\pgfplotspointaxisxy{0}{0}}
            {\pgfplotspointaxisdirectionxy{1}{1}}
            {\pgfplotspointaxisdirectionxy{0}{2}}
    };
    \addplot [only marks,mark=*] coordinates
        { (0,0) };
\end{axis}
\end{tikzpicture}
\end{codeexample}

    由于变换是在坐标轴完成后才初始化的，因此此命令需要放在 \tikzname{} |\path| 命令内部（如上面例子中的 |\draw|），或者放在 |\pgfplotsextra| 内部。
\end{commandlist}

\begin{commandlist}{%
    \pgfplotspointrelaxisxy\marg{rel x coordinate}\marg{rel y coordinate},
    \pgfplotspointrelaxisxyz\marg{rel x coordinate}\marg{rel y coordinate}\marg{rel z coordinate}%
}
    接受\emph{相对}坐标的点命令：$x=0$ 对应坐标轴的\emph{下}限，$x=1$ 对应\emph{上}限。

    这些命令用于 |rel axis cs|。

    请注意，变换只有在坐标轴完成后才会初始化！这意味着你需要使用 |\pgfplotsextra|。
\end{commandlist}

\begin{commandlist}{%
    \pgfplotspointdescriptionxy\marg{$x$ fraction}\marg{$y$ fraction},
    \pgfplotsqpointdescriptionxy\marg{$x$ fraction}\marg{$y$ fraction}%
}
    点命令：|{0}{0}| 是坐标轴包围盒的左下角，|{1}{1}| 是右上角；其余值位于二者之间。`|q|' 版本更快，因为它不对参数调用数学解析器。

    它们用于 |axis description cs|，见
    Section~\ref{pgfplots:sec:axis:description:cs}。
\end{commandlist}

\begin{commandlist}{\pgfplotspointaxisorigin}
    原点处的点坐标，即 $(0,0,0)$。如果原点不在坐标轴范围内，则返回边界上最近的点。

    这与 |origin| 锚点返回的坐标相同。
\end{commandlist}

\begin{commandlist}{%
    \pgfplotstransformcoordinatex\marg{x coordinate of an axis},
    \pgfplotstransformcoordinatey\marg{y coordinate of an axis},
    \pgfplotstransformcoordinatez\marg{z coordinate of an axis}%
}
    把 |\pgfmathresult| 定义为与输入参数对应的 \PGF{} 底层坐标。

    该命令会应用任何 |[xyz] coord trafo| 键、数据缩放和/或对数，以及 \PGFPlots{} 用于把输入坐标映射为内部坐标的任何其他处理。

    结果可以在 |\pgfpointxy| 语句内部使用（也就是说，它仍然需要与相应的 \PGF{} 单位向量相乘缩放）。
    %
\begin{codeexample}[]
\begin{tikzpicture}
\pgfplotsset{compat/pgfpoint substitution=1.3}
\begin{axis}[xmin=0,xmax=2,ymin=0,ymax=5]
    \pgfplotsextra{
        \pgfplotstransformcoordinatex{1}
        \let\xcoord=\pgfmathresult
        \pgfplotstransformcoordinatey{1}
        \let\ycoord=\pgfmathresult
        \pgfpathcircle
            {\pgfqpointxy{\xcoord}{\ycoord}}
            {5pt}
        \pgfusepath{fill}
    }
\end{axis}
\end{tikzpicture}
\end{codeexample}
    %
    请注意，\PGFPlots{} 默认会把 |\pgfqpointxy| 替换为 |\pgfplotspointaxisxyz|——而此命令本身已经隐式完成了坐标变换。为了看出区别，前面的例子先通过
    |compat/pgfpoint substitution=1.3| 关闭了这种坐标系的自动替换。
    %
    此命令的结果也可以作为数学方法
    |transformcoordinatex| 使用（参见 |axis cs| 的文档）。

    请注意，变换只有在坐标轴完成后才会初始化。这意味着你需要像上面例子那样使用 |\pgfplotsextra|。
\end{commandlist}

\begin{commandlist}{%
    \pgfplotstransformdirectionx\marg{x direction of an axis},
    \pgfplotstransformdirectiony\marg{y direction of an axis},
    \pgfplotstransformdirectionz\marg{z direction of an axis}%
}
    把 |\pgfmathresult| 定义为 \PGF{} 底层\emph{方向向量分量}。

    方向向量需要与某个坐标\emph{相加}才能得到一个坐标，参见
    |\pgfplotspointaxisdirectionxy| 和 |axis direction cs| 的文档。

    参数 \meta{x direction of an axis} 的处理方式（几乎）与处理绝对位置的宏
    |\pgfplotstransformcoordinatex| 相同。唯一区别是：\emph{方向}不需要平移变换。

    此命令的结果也可以作为数学方法
    |transformdirectionx| 使用（参见 |axis direction cs| 的文档）。

    关于此命令的细节与示例，请参见 |axis direction cs|。
\end{commandlist}

% this command is for internal use only:
%--------------------------------------------------
% \begin{command}{\pgfplotsconvertunittocoordinate\marg{x, y or z}\marg{dimension}}
%     Converts a dimension (with unit!) to a corresponding $x$-, $y$- or $z$-coordinate. The result will be written to |\pgfmathresult| (without units).
%
%     It is possible to use the result as arguments for the |\pgfpointxyz| commands.
%
%     The effect is to multiply \meta{dimension} with the inverse length of the unit vector for the specified axis. These lengths are precomputed in \PGFPlots{} so the operation is fast.
% \begin{codeexample}[code only]
% \pgfplotsconvertunittocoordinate{x}{5pt}
% % now, the command uses exactly 5pt in x direction:
% \pgfqpointxyz{\pgfmathresult}{4}{3}
% \end{codeexample}
% \end{command}
%--------------------------------------------------

\begin{commandlist}{%
    \pgfplotspointunitx,
    \pgfplotspointunity,
    \pgfplotspointunitz%
}
    返回画布 $x$、$y$ 或 $z$ 单位向量的底层点命令。

    |\pgfplotspointunitx| 是 $x$ 方向上的 \pgfname{} 单位向量。

    这些向量本质上分别等同于\linebreak
    |\pgfqpointxyz{1}{0}{0}|、|\pgfqpointxyz{0}{1}{0}|\linebreak
    和 |\pgfqpointxyz{0}{0}{1}|。

    单位 $z$ 向量只在三维坐标轴中才有定义。
\end{commandlist}

\begin{commandlist}{%
    \pgfplotsunitxlength,
    \pgfplotsunitylength,
    \pgfplotsunitzlength,
    \pgfplotsunitxinvlength,
    \pgfplotsunityinvlength,
    \pgfplotsunitzinvlength%
}
    展开为相应单位向量 $x_i$ 的长度 $\lVert x_i \rVert$，或其倒数 $1/\lVert x_i \rVert$ 的宏。这些宏可以在 |\pgfmathparse| 内部使用。

    $x_i$ 即 |\pgfplotspointunitx| 系列。
\end{commandlist}

\begin{command}{\pgfplotsqpointoutsideofaxis%
        \marg{three-char-string}\marg{coordinate}\marg{normal distance}%
}
    对于二维图，返回四条可用坐标轴之一上的一个点坐标；对于三维图，则返回十二条可用坐标轴之一上的点坐标。

    所指定的坐标轴由一个三字符字符串唯一标识，作为命令的第一个参数传入。三个字符中的第一个为 `|0|'，表示该坐标轴的 $x$ 坐标经过坐标轴的下限；为 `|1|'，表示经过上限；为 `|2|'，表示经过原点。第二个字符同样是 |0|、|1| 或 |2|，表示 $y$ 轴上的位置。第三个字符对应第三维，即 $z$ 轴；二维图应保留为 `|0|'。但是，三个字符中\emph{必须}有一个为 `|v|'，表示该坐标轴\underline{是变化的}。例如，|v01| 表示 $\{ (x,y_{\min},z_{\max}) \vert x \in \R \}$。

    第二个参数 \meta{coordinate} 是该坐标轴上的逻辑坐标。由于该坐标轴的两个坐标是固定的，\meta{coordinate} 指的是坐标轴中\underline{变化}的那一分量。它必须是一个不带单位的数字；此处不支持数学表达式。

    第三个参数 \meta{normal distance} 是一个长度，例如 |10pt|。它把坐标沿外法向量方向从指定坐标轴向外平移。外法向量总是指向远离坐标轴的方向。它由 |\pgfplotspointouternormalvectorofaxis| 计算。

    此命令有若干变体，都记录在源代码中。其中一个特别有用：
\end{command}

\begin{command}{\pgfplotsqpointoutsideofaxisrel%
        \allowbreak\marg{three-char-string}\allowbreak\marg{axis fraction}\allowbreak\marg{normal distance}%
}
    此点坐标是 |\pgfplotsqpointoutsideofaxis| 的一个变体，它允许传入 \meta{axis fraction} 而不是绝对坐标。该分数是 $0$（坐标轴下限）到 $1$（坐标轴上限）之间的一个数，即以坐标轴总长的百分比表示。允许使用负值或大于 $1$ 的值。

    |\pgfplotsqpointoutsideofaxisrel| 命令在思路上类似于
    |rel axis cs|。

    与反向坐标轴配合使用时有一点特殊：如果坐标轴已通过 |x dir=reverse| 反向，且 |allow reversal of rel axis cs| 为 true，那么 $0$ 表示\emph{上}限，$1$ 表示\emph{下}限。这样，仅因坐标轴反向，坐标值不会发生改变。
        \index{allow reversal of rel axis cs}%
\end{command}

\begin{command}{\pgfplotspointouternormalvectorofaxis\marg{three-char-string}}
    一个点命令，返回相应坐标轴的外法向量。该法向量长度为 $1$（用
    |\pgfpointnormalised| 计算）。它与 |\pgfplotsqpointoutsideofaxis| 及其变体中使用的是同一个法向量。

    此命令的输出会被缓存，并在坐标轴的整个生命周期内复用。
\end{command}

\begin{command}{\pgfplotsticklabelaxisspec\marg{x, y or z}}
    展开为指定坐标轴（\meta{x}、\meta{y} 或 \meta{z}）所在的、用于承载刻度标签的坐标轴的三字符标识。
\end{command}

\begin{command}{\pgfplotsvalueoflargesttickdimen\marg{x, y or z}}
    展开为刻度位置到其刻度标签包围盒沿外单位法向量方向的最大距离。它还包括 |ticklabel shift| 键的值。

    此值用于 |ticklabel cs|。
\end{command}

\begin{commandlist}{
    \pgfplotsmathfloatviewdepthxyz\marg{x}\marg{y}\marg{z},
    \pgfplotsmathviewdepthxyz\marg{x}\marg{y}\marg{z}%
}
    这两个宏都把 |\pgfmathresult| 定义为三维点 $\bar x = (x,y,z)$ 的「深度」。深度定义为 $\bar x$ 与当前坐标轴视图方向 $\vec d$ 的内积。

    对于 |\pgfplotsmathfloatviewdepthxyz|，参数按浮点数解析，结果也以浮点数编码。可以用
    |\pgfmathfloattofixed{\pgfmathresult}| 生成定点表示。

    对于 |\pgfplotsmathviewdepthxyz|，内积使用 \TeX{} 算术完成，结果以定点数赋值。它略快，但数值范围小得多。

    这两个命令都只能在三维 \PGFPlots{} 坐标轴\emph{内部}使用（坐标轴初始化之后，参见
    |\pgfplotsextra|）。
\end{commandlist}

\begin{texif}{pgfplotsthreedim}
    一个 \TeX{} |\if|：如果坐标轴是三维的，求 \meta{true code}；否则求 \meta{else code}。
\end{texif}


\section{访问坐标轴范围}

也可以在可视化阶段，即在 |\end{axis}| 期间，访问坐标轴范围。请参见
第~\pageref{page:access:limits} 页上 |xmin| 的参考文档。


\section{访问点坐标的值}

在可视化阶段，\PGFPlots{} 提供了对当前正在处理的坐标及其值的访问。

这种访问需要调用特定的宏。这些宏把坐标值写入一些公开的键值对中。然后，就可以访问当前点的 $x$、$y$、$z$ 和颜色数据。

\begin{commandlist}{\pgfplotspointgetcoordinates,\pgfplotspointgetcoordinates\marg{point}}

    调用此宏之后，会设置以下键：

    \declaretext{/data point/x} 将包含当前点的 $x$ 坐标。

    \declaretext{/data point/y} 将包含当前点的 $y$ 坐标。

    \declaretext{/data point/z} 将包含当前点的 $z$ 坐标（如果适用）。

    \declaretext{/data point/meta} 将包含当前点的
    |point meta| 值（如果适用）。

    \declaretext{/data point/index} 将包含当前点在坐标流中的索引。它实际上与 |\coordindex| 相同。

    此命令实际上支持两种工作模式：
    %
    \begin{enumerate}
        \item 不带参数。此时，它返回当前绘图处理器即将处理的点的值。
        \item 带花括号参数。此时，它接受一个坐标并相应地为各键赋值。请注意，此命令仅支持二维坐标轴，并且只赋值
            |/data point/x| 和 |/data point/y|。
    \end{enumerate}

    返回的值与你在坐标轴上读到的相同，也与你写进 |axis cs| 的相同。

    这意味着任何 |x coord inv trafo| 都已经作用在该值上。也就是说，即使内部坐标是以对数形式存储的，这里也已经调用了指数函数。

    对于任何 |scatter| 图（包括 |nodes near coords|），此函数都会被隐式调用。这样就可以一次性访问\emph{所有}坐标值：
    %
\begin{codeexample}[]
\begin{tikzpicture}
\begin{axis}
    \addplot+ [
        domain=0:360,
        samples=9,
        only marks,
        nodes near coords={%
            \footnotesize
            $(\pgfmathprintnumber
                {\pgfkeysvalueof{/data point/x}},
               \pgfmathprintnumber
                {\pgfkeysvalueof{/data point/y}})$%
        },
    ] {sin(x)};
\end{axis}
\end{tikzpicture}
\end{codeexample}
    %
    这个例子之所以能工作，是因为 |\pgfplotspointgetcoordinates| 是 |nodes near coords| 标准实现的一部分；得到的值可以直接使用。请注意，如果只打印一个值，前面的例子还可以更简单：|nodes near coords| 会退回到使用 |point meta|。而 |point meta| 默认情况下本来就包含 $y$ 坐标。

    一个更高级的例子是：在 |ybar| 图中，希望把节点放在坐标轴下端，并从节点到相应柱条画一些虚线：
    %
\begin{codeexample}[]
\begin{tikzpicture}
\begin{axis}[
    ybar,
    nodes near coords,
    %
    % we want to provide absolute 'at' values
    % for the nodes:
    scatter/position=absolute,
    every node near coord/.style={
        at={(\pgfkeysvalueof{/data point/x},-1)},
        % pretty printing:
        anchor=north,
        /pgf/number format/fixed,
        /pgf/number format/precision=1,
        % assign a name which can be referenced below:
        name=NNC\pgfkeysvalueof{/data point/index},
    },
    % ... draw a dotted line between
    % the marker and the bar:
    /pgfplots/scatter/@post marker code/.add code={}{
        \draw [dotted,help lines]
            (NNC\pgfkeysvalueof{/data point/index})
            -- (\pgfkeysvalueof{/data point/x},
            {min(0,\pgfkeysvalueof{/data point/y})});
    },
    % assign suitable tick labels:
    xtick=data,
]
    % some dummy data:
    \addplot+ [
        domain=0:360,
        bar width=360/9,
        samples=9,
    ] {sin(x)};
\end{axis}
\end{tikzpicture}
\end{codeexample}
    %
    同样，该命令是作为 |nodes near coords| 的一部分隐式使用的，因此在例子本身中并未出现。

    \paragraph{另请参见} 在线相关示例：
    \url{https://tex.stackexchange.com/a/141006}。它演示了如何根据数值（放在柱条内部或柱条上方）放置由 |nodes near coords| 生成的节点。

    下面的例子使用一个花括号参数，我们希望获取该参数的坐标值：
    %
\begin{codeexample}[]
% requires \usetikzlibrary{intersections}
\begin{tikzpicture}
\begin{axis}
    \addplot [name path=A] {x^3};

    \draw [red,name path=HelperLine]
        (axis cs:-6,-100) -- (axis cs:6,130);

    \draw [
        font=\footnotesize,
        name intersections={of=A and HelperLine},
    ]
    node [pin={right:
      \pgfplotspointgetcoordinates{(intersection-1)}
        $(\pgfmathprintnumber[fixed]{
            \pgfkeysvalueof{/data point/x}},
          \pgfmathprintnumber[fixed]{
            \pgfkeysvalueof{/data point/y}})$
      }
    ] at (intersection-1) {}
    node [pin={
      \pgfplotspointgetcoordinates{(intersection-2)}
        $(\pgfmathprintnumber[fixed]{
            \pgfkeysvalueof{/data point/x}},
          \pgfmathprintnumber[fixed]{
            \pgfkeysvalueof{/data point/y}})$
      }
    ] at (intersection-2) {};
\end{axis}
\end{tikzpicture}
\end{codeexample}
    %
    该例子计算并命名了 |A| 与 |HelperLine| 的交点。下面的代码在交点处生成 pin。为此，它使用
    |\pgfplotspointgetcoordinates|\marg{point}，该宏会定义 |/data point/x|
    和 |/data point/y|。然后，用 |\pgfmathprintnumber| 把这些值格式化输出。

    在第二种模式下，|\pgfplotspointgetcoordinates|\marg{point} 期望 \meta{point} 为以下二者之一：
    %
    \begin{enumerate}
        \item 一个基础层 \PGF{} 点，例如 |\pgfpointanchor{A}{center}|；或
        \item 一个 \Tikz{} 点，例如 |(A.base)| 或 |(3,5)|。
    \end{enumerate}
\end{commandlist}

\begin{command}{\pgfplotspointgetnormalizedcoordinates}
    与 |\pgfplotspointgetcoordinates| 非常相似的一个宏。同样，它也应在可视化阶段调用。

    它赋值的输出宏完全相同，但值不同。更准确地说，它定义的宏为：

    \declaretext{/data point/x} 将包含当前点的\emph{归一化} $x$ 坐标。

    \declaretext{/data point/y} 将包含当前点的\emph{归一化} $y$ 坐标。

    \declaretext{/data point/z} 将包含当前点的\emph{归一化} $z$ 坐标（如果适用）。

    \declaretext{/data point/meta} 将包含当前点的
    |point meta| 值（如果适用）。

    \declaretext{/data point/index} 将包含当前点在坐标流中的索引。它实际上与 |\coordindex| 相同。

    关键字 \declaretext{normalized} 表示这些值采用适合坐标轴使用的数值形式。更具体地说：任何用户 |x coord inv trafo| 都会被\emph{忽略}。一个重要的例子是 |symbolic x coords|：归一化坐标将是某些关联数字，而不是符号。而
    |\pgfplotspointgetcoordinates| 返回的结果则是符号。对于对数坐标轴，归一化值是对数。

    通常，当你希望执行某种数学运算（如求平均或做差）时，归一化值会更有用。

    此函数需要显式调用。目前 |ybar stacked| 用它来对齐 |nodes near coords|。
\end{command}


\section{图层访问}

\begin{command}{\pgfplotsonlayer\marg{layer name}}
    一个底层命令，它会检查当前坐标轴是否启用了图层支持；如果启用，则调用 |\pgfonlayer|\marg{layer name}。

    必须有一个 |\endpgfplotsonlayer| 来结束该环境。
\end{command}

\begin{command}{\endpgfplotsonlayer}
    |\pgfplotsonlayer| 的结束。
\end{command}

\begin{command}{\pgfonlayer\marg{layer name}}
    \PGF{} 的一个底层命令，它会把直到匹配的 |\endpgfonlayer| 之前的所有内容收集到图层 \meta{layer name} 中。

    \meta{layer name} 必须是活动的，即它必须是 |set layers| 图层名集合的一部分。

    唯一的特殊情况是：如果你在某处调用 |\pgfdeclarelayer{discard}|，这个特殊图层具有一个「神奇的名字」，当用 |\pgfonlayer{discard}| 启用时，它充当 |/dev/null|：它不必是活动的，分配给该图层的所有内容都会被丢弃——只要它不在图层名配置之中。

    必须有一个 |\endpgfonlayer| 来结束该环境。
\end{command}

\begin{command}{\endpgfonlayer}
    |\pgfonlayer| 的结束。
\end{command}


\begin{command}{\pgfsetlayers\marg{layer list}}
    这是 \PGF{} 的一个底层命令。在撰写本文时，它是告诉 \PGF{} 当前/下一张图应使用哪些图层的唯一方式。它由 |set layers| 隐式使用。
\end{command}


\printindex

\bibliographystyle{abbrv} %gerapali} %gerabbrv} %gerunsrt.bst} %gerabbrv}% gerplain}
\nocite{pgfplotstable}
\nocite{programmingnotes}
\bibliography{pgfplots}
\end{document}

% -----------------------------------------------------------------------------
% For Stefan Pinnow as reminder on what to look for when editing the manual
% -----------------------------------------------------------------------------
% There should be no line breaks in the following environments
% - |...|
% - \declareandlabel{...}
% - \verbpdfref{...}
% If "MakeTikzPictures" isn't running through check one of the externalized
% LOG files what is the cause of that.
% -----------------------------------------------------------------------------
'
    \expandafter\ifx\csname pgfplotsclickabledisabled\endcsname\relax
        \usepgfplotslibrary{clickable}
    \fi
\fi

%\usepackage{fp}
% ATTENTION:
% this requires pgf version NEWER than 2.00:
%\usetikzlibrary{fixedpointarithmetic}

\usepgfplotslibrary{dateplot,units,groupplots,statistics}

\pgfplotsset{
    table/search path={.,plotdata}
}

\usepackage[a4paper,left=2.25cm,right=2.25cm,top=2.5cm,bottom=2.5cm]{geometry}
\usepackage{amsmath,amssymb}
\usepackage{xxcolor}

\usepackage{ifluatex}

% tlmgr install zapfding symbol
\usepackage{pifont}

\def\lua{\textsc{lua}}

\ifluatex
    %\usepackage[utf8]{luainputenc}
\else
    \expandafter\ifx\csname XeTeXrevision\endcsname\relax
        \usepackage[utf8]{inputenc}
    \fi
\fi
\usepackage{amsmath}

% tlmgr eurosym
\usepackage{eurosym}

% tlmgr install units
\usepackage{nicefrac}

\def\eps{\textsc{eps}}

% Copyright 2006 by Till Tantau
%
% This file may be distributed and/or modified
%
% 1. under the LaTeX Project Public License and/or
% 2. under the GNU Free Documentation License.
%
% See the file doc/generic/pgf/licenses/LICENSE for more details.

% $Header: /cvsroot/pgfplots/pgfplots/doc/latex/pgfplots/pgfmanual-en-macros.tex,v 1.4 2009/03/05 22:23:22 ludewich Exp $


% This file is part of the pgf manual; I reuse it here with
% permission.
%
% I have changed very few things (marked with 'CF')

\providecommand\href[2]{\texttt{#1}}


\colorlet{examplefill}{yellow!80!black}
\definecolor{graphicbackground}{rgb}{0.96,0.96,0.8}
\definecolor{codebackground}{rgb}{0.8,0.8,1}

\newenvironment{pgfmanualentry}{%
  \pgfmanualresetlabel
  \list{}{\leftmargin=2em\itemindent-\leftmargin\def\makelabel##1{\hss##1}}%
}{\endlist}
\newcounter{pgfmanualentry}
\newcommand\pgfmanualentryheadline[1]{%
  \itemsep=0pt\parskip=0pt\relax
  {\raggedright\item\refstepcounter{pgfmanualentry}\pgfmanualstorelabel
   \strut{#1}\par}%
  \topsep=0pt}
\newcommand\pgfmanualbody{\parskip3pt}

\makeatletter
\newif\if@pgfmanualfirstlabel

\newcommand\pgfmanualresetlabel{%
  \@pgfmanualfirstlabeltrue}

% Store info of first label in \pgfmanual@labelinfo for latter use.
\newcommand\pgfmanualstorelabel{%
  \if@pgfmanualfirstlabel
    \begingroup
    \label@hook
    \xdef\pgfmanual@labelinfo
      {{\@currentlabel}{\thepage}{\@currentlabelname}{\@currentHref}{}}%
    \endgroup
    \global\@pgfmanualfirstlabelfalse
  \fi
}

% Used like a prifix of \label, \labelprefix makes the following 
% \label use label info stored in \pgfmanual@labelinfo. Based on the 
% 
% Usage:
%     \labelprefix\label{<label name>}
% Possible extension:
%     \labelprefix[<offset>]\label{<label name>}
%
\newcommand\labelprefix[2]{%
  \def\@pgfmanualtemp{#1}%
  \ifx\@pgfmanualtemp\pgfmanual@labelcmd
    \@bsphack
    \begingroup
      \def\label@name{#2}%
      % \label@hook
      \protected@write\@auxout{}{%
        \string\newlabel{#2}{\pgfmanual@labelinfo}}%
    \endgroup
    \@esphack
  \else
    \PackageError{pgfmanual-macros}{\string\labelprefix should use as prefix of \string\label.}{}%
  \fi
}
\def\pgfmanual@labelcmd{\label}
\makeatother

\newenvironment{pgflayout}[1]{
  \begin{pgfmanualentry}
    \pgfmanualentryheadline{\texttt{\string\pgfpagesuselayout\char`\{\declare{#1}\char`\}}\oarg{options}}
    \index{#1@\protect\texttt{#1} layout}%
    \index{Page layouts!#1@\protect\texttt{#1}}%
    \pgfmanualbody
}
{
  \end{pgfmanualentry}
}


\newenvironment{command}[1]{
  \begin{pgfmanualentry}
    \extractcommand#1\@@
    \pgfmanualbody
}
{
  \end{pgfmanualentry}
}

%% MW: START MATH MACROS
\def\mvar#1{{\rmfamily\textit{#1}}}

\makeatletter

\def\extractmathfunctionname#1{\extractmathfunctionname@#1(,)\tmpa\tmpb}
\def\extractmathfunctionname@#1(#2)#3\tmpb{\def\mathname{#1}}

\def\extractmathoperatorname{\begingroup\def\mvar##1{}\def\ {}\extractmathoperatorname@}
\def\extractmathoperatorname@#1{\xdef\mathname{#1}\endgroup}

\makeatother

\newenvironment{math-function}[1]{
    \begin{pgfmanualentry}
        \extractmathfunctionname{#1}
        \pgfmanualentryheadline{\texttt{#1}}%
        \index{\mathname @\protect\texttt{\mathname} math function}%
        \index{Math functions!\mathname @\protect\texttt{\mathname}}
        \pgfmanualbody
}
{
    \end{pgfmanualentry}
}

\newenvironment{math-operator}[1]{
    \begin{pgfmanualentry}
        \extractmathoperatorname{#1}
        \pgfmanualentryheadline{\texttt{#1}}%
        \index{\mathname @\protect\texttt{\mathname} math operator}%
        \index{Math operators!\mathname @\protect\texttt{\mathname}}
        \pgfmanualbody
}
{%
    \end{pgfmanualentry}
}

\newenvironment{math-constant}[1]{
    \begin{pgfmanualentry}
        \pgfmanualentryheadline{\texttt{#1}}%
        \index{#1@\protect\texttt{#1} math constant}%
        \index{Math constants!#1@\protect\texttt{#1}}
        \pgfmanualbody
}
{
    \end{pgfmanualentry}
}
\def\calcname{\textsc{calc}}
%% MW: END MATH MACROS


\def\extractcommand#1#2\@@{%
  \pgfmanualentryheadline{\declare{\texttt{\string#1}}#2}%
  \removeats{#1}%
  \index{\strippedat @\protect\myprintocmmand{\strippedat}}}


% \begin{environment}{{name}\marg{arguments}}
\renewenvironment{environment}[1]{
  \begin{pgfmanualentry}
    \extractenvironement#1\@@
    \pgfmanualbody
}
{
  \end{pgfmanualentry}
}

\def\extractenvironement#1#2\@@{%
  \pgfmanualentryheadline{{\ttfamily\char`\\begin\char`\{\declare{#1}\char`\}}#2}%
  \pgfmanualentryheadline{{\ttfamily\ \ }\meta{environment contents}}%
  \pgfmanualentryheadline{{\ttfamily\char`\\end\char`\{\declare{#1}\char`\}}}%
  \index{#1@\protect\texttt{#1} environment}%
  \index{Environments!#1@\protect\texttt{#1}}}


\newenvironment{plainenvironment}[1]{
  \begin{pgfmanualentry}
    \extractplainenvironement#1\@@
    \pgfmanualbody
}
{
  \end{pgfmanualentry}
}

\def\extractplainenvironement#1#2\@@{%
  \pgfmanualentryheadline{{\ttfamily\declare{\char`\\#1}}#2}%
  \pgfmanualentryheadline{{\ttfamily\ \ }\meta{environment contents}}%
  \pgfmanualentryheadline{{\ttfamily\declare{\char`\\end#1}}}%
  \index{#1@\protect\texttt{#1} environment}%
  \index{Environments!#1@\protect\texttt{#1}}}


\newenvironment{contextenvironment}[1]{
  \begin{pgfmanualentry}
    \extractcontextenvironement#1\@@
    \pgfmanualbody
}
{
  \end{pgfmanualentry}
}

\def\extractcontextenvironement#1#2\@@{%
  \pgfmanualentryheadline{{\ttfamily\declare{\char`\\start#1}}#2}%
  \pgfmanualentryheadline{{\ttfamily\ \ }\meta{environment contents}}%
  \pgfmanualentryheadline{{\ttfamily\declare{\char`\\stop#1}}}%
  \index{#1@\protect\texttt{#1} environment}%
  \index{Environments!#1@\protect\texttt{#1}}}


\newenvironment{shape}[1]{
  \begin{pgfmanualentry}
      \pgfmanualentryheadline{形状 {\ttfamily\declare{#1}}}%
    \index{#1@\protect\texttt{#1} shape}%
    \index{Shapes!#1@\protect\texttt{#1}}
    \pgfmanualbody
}
{
  \end{pgfmanualentry}
}


\newenvironment{handler}[1]{
  \begin{pgfmanualentry}
    \extracthandler#1\@nil%
    \pgfmanualbody
}
{
  \end{pgfmanualentry}
}

\def\gobble#1{}
\def\extracthandler#1#2\@nil{%
  \pgfmanualentryheadline{键处理器 \meta{key}{\ttfamily/\declare{#1}}#2}%
  \index{\gobble#1@\protect\texttt{#1} handler}%
  \index{Key handlers!#1@\protect\texttt{#1}}
}


\makeatletter


\newenvironment{stylekey}[1]{
  \begin{pgfmanualentry}
    \def\extrakeytext{样式， }
    \extractkey#1\@nil%
    \pgfmanualbody
}
{
  \end{pgfmanualentry}
}


% \begin{key}{/path/x=value}
% \begin{key}{/path/x=value (initially XXX)}
% \begin{key}{/path/x=value (default XXX)}
\newenvironment{key}[1]{
  \begin{pgfmanualentry}
    \def\extrakeytext{}
    %\def\altpath{\emph{\color{gray}or}}%
    \extractkey#1\@nil%
    \pgfmanualbody
}
{
  \end{pgfmanualentry}
}

\def\extractkey#1\@nil{%
  \pgfutil@in@={#1}%
  \ifpgfutil@in@%
    \extractkeyequal#1\@nil
  \else%
    \pgfutil@in@{(initial}{#1}%
    \ifpgfutil@in@%
      \extractequalinitial#1\@nil%
    \else
      \pgfmanualentryheadline{{\ttfamily\declarekey{#1}}\hfill(\extrakeytext{}无值)}%
      \def\mykey{#1}%
      \def\mypath{}%
      \def\myname{}%
      \firsttimetrue%
      \decompose#1/\nil%
    \fi
  \fi%
}

\pgfkeyssetvalue{/pgfmanual/gray key prefixes}{/pgfplots/,/tikz/}

\def\declarekey#1{%
    \begingroup
    \gdef\declarekey@didit{0}%
    \pgfkeysgetvalue{/pgfmanual/gray key prefixes}\declarekey@gray
    \def\declarekey@loop{\foreach \prefix in }%
    \expandafter\declarekey@loop\expandafter{\declarekey@gray} {%
        \expandafter\pgfutil@in@\expandafter{\prefix}{#1}%
        \ifpgfutil@in@
            \expandafter\def\expandafter\declarekey@handleprefix\prefix####1\relax{%
                \keyprefix{\prefix}\declare{####1}%
            }%
            \declarekey@handleprefix#1\relax
            \gdef\declarekey@didit{1}%
            \breakforeach
        \fi
    }%
    \if0\declarekey@didit
        \declare{#1}%
    \fi
    \endgroup
}%
\def\keyprefix#1{\textcolor{gray}{#1}}

\def\extractkeyequal#1=#2\@nil{%
  \pgfutil@in@{(default}{#2}%
  \ifpgfutil@in@%
    \extractdefault{#1}#2\@nil%
  \else%
    \pgfutil@in@{(initial}{#2}%
    \ifpgfutil@in@%
      \extractinitial{#1}#2\@nil%
    \else
%     \pgfmanualentryheadline{{\ttfamily\declarekey{#1}=}#2\hfill(\extrakeytext no default)}%
      \pgfmanualentryheadline{{\ttfamily\declarekey{#1}=}#2\hfill}%
    \fi%
  \fi%
  \def\mykey{#1}%
  \def\mypath{}%
  \def\myname{}%
  \firsttimetrue%
  \decompose#1/\nil%
}

\def\extractdefault#1#2(default #3)\@nil{%
  \pgfmanualentryheadline{{\ttfamily\declarekey{#1}\opt{=}}\opt{#2}\hfill (\extrakeytext{}默认 {\ttfamily#3})}%
}

\def\extractinitial#1#2(initially #3)\@nil{%
%  \pgfmanualentryheadline{{\ttfamily\declarekey{#1}=}#2\hfill (\extrakeytext no default, initially {\ttfamily#3})}%
  \pgfmanualentryheadline{{\ttfamily\declarekey{#1}=}#2\hfill (\extrakeytext{}初始 {\ttfamily#3})}%
}

\def\extractequalinitial#1 (initially #2)\@nil{%
  \pgfmanualentryheadline{{\ttfamily\declarekey{#1}}\hfill (\extrakeytext{}初始 {\ttfamily#2})}%
  \def\mykey{#1}%
  \def\mypath{}%
  \def\myname{}%
  \firsttimetrue%
  \decompose#1/\nil%
}

% Introduces a key alias '/#1/<name of current key>'
% to be used inside of \begin{key} ... \end{key}
\def\keyalias#1{\vspace{-3pt}\item{\small 别名 {\ttfamily/#1/\myname}}\vspace{-2pt}\par}

\newif\iffirsttime

\makeatother

\def\decompose/#1/#2\nil{%
  \def\test{#2}%
  \ifx\test\empty%
    % aha.
    \index{#1@\protect\texttt{#1} key}%
    \index{\mypath#1@\protect\texttt{#1}}%
    \def\myname{#1}%
  \else%
    \iffirsttime
%    CF : disabled for /pgfplots/ prefix.
%        \def\mypath{#1@\protect\texttt{/#1/}!}%
%        \firsttimefalse
        \def\pgfplotslocTMPa{pgfplots}%
        \edef\pgfplotslocTMPb{#1}%
        \ifx\pgfplotslocTMPb\pgfplotslocTMPa
            \def\mypath{}%
        \else
            \def\mypath{#1@\protect\texttt{/#1/}!}%
        \fi
        \firsttimefalse
    \else
      \expandafter\def\expandafter\mypath\expandafter{\mypath#1@\protect\texttt{#1/}!}%
    \fi
    \def\firsttime{}
    \decompose/#2\nil%
  \fi%
}

\newenvironment{predefinedmethod}[1]{
  \begin{pgfmanualentry}
    \extractpredefinedmethod#1\@nil
    \pgfmanualbody
}
{
  \end{pgfmanualentry}
}
\def\extractpredefinedmethod#1(#2)\@nil{%
  \pgfmanualentryheadline{方法 \declare{\ttfamily #1}\texttt(#2\texttt) \hfill（为所有类预定义）}
  \index{#1@\protect\texttt{#1} method}%
  \index{Methods!#1@\protect\texttt{#1}}
}


\newenvironment{ooclass}[1]{
  \begin{pgfmanualentry}
    \def\currentclass{#1}
    \pgfmanualentryheadline{类 \declare{\texttt{#1}}}
    \index{#1@\protect\texttt{#1} class}%
    \index{Class #1@Class \protect\texttt{#1}}%
    \index{Classes!#1@\protect\texttt{#1}}
    \pgfmanualbody
}
{
  \end{pgfmanualentry}
}

\newenvironment{method}[1]{
  \begin{pgfmanualentry}
    \extractmethod#1\@nil
    \pgfmanualbody
}
{
  \end{pgfmanualentry}
}
\def\extractmethod#1(#2)\@nil{%
  \def\test{#1}
  \ifx\test\currentclass
    \pgfmanualentryheadline{构造函数 \declare{\ttfamily #1}\texttt(#2\texttt)}
  \else
    \pgfmanualentryheadline{方法 \declare{\ttfamily #1}\texttt(#2\texttt)}
  \fi
  \index{#1@\protect\texttt{#1} method}%
  \index{Methods!#1@\protect\texttt{#1}}
  \index{Class \currentclass!#1@\protect\texttt{#1}}%
}

\iffalse
\newenvironment{attribute}[1]{
  \begin{pgfmanualentry}
    \extractattribute#1\@nil
    \pgfmanualbody
}
{
  \end{pgfmanualentry}
}
\fi
\def\extractattribute#1=#2;\@nil{%
  \def\test{#2}%
  \ifx\test\@empty
    \pgfmanualentryheadline{私有属性 \declare{\ttfamily #1} \hfill（初始为空）}
  \else
    \pgfmanualentryheadline{私有属性 \declare{\ttfamily #1} \hfill（初始为 {\ttfamily #2}）}
  \fi
  \index{#1@\protect\texttt{#1} attribute}%
  \index{Attributes!#1@\protect\texttt{#1}}
  \index{Class \currentclass!#1@\protect\texttt{#1}}%
}



\newenvironment{predefinednode}[1]{
  \begin{pgfmanualentry}
    \pgfmanualentryheadline{预定义节点 {\ttfamily\declare{#1}}}%
    \index{#1@\protect\texttt{#1} node}%
    \index{Predefined node!#1@\protect\texttt{#1}}
    \pgfmanualbody
}
{
  \end{pgfmanualentry}
}

\newenvironment{coordinatesystem}[1]{
  \begin{pgfmanualentry}
    \pgfmanualentryheadline{坐标系 {\ttfamily\declare{#1}}}%
    \index{#1@\protect\texttt{#1} coordinate system}%
    \index{Coordinate systems!#1@\protect\texttt{#1}}
    \pgfmanualbody
    \pgfmanualpdflabel{#1}{}%
}
{
  \end{pgfmanualentry}
}

\newenvironment{snake}[1]{
  \begin{pgfmanualentry}
    \pgfmanualentryheadline{蛇形 {\ttfamily\declare{#1}}}%
    \index{#1@\protect\texttt{#1} snake}%
    \index{Snakes!#1@\protect\texttt{#1}}
    \pgfmanualbody
}
{
  \end{pgfmanualentry}
}

\newenvironment{decoration}[1]{
  \begin{pgfmanualentry}
    \pgfmanualentryheadline{装饰 {\ttfamily\declare{#1}}}%
    \index{#1@\protect\texttt{#1} decoration}%
    \index{Decorations!#1@\protect\texttt{#1}}
    \pgfmanualbody
}
{
  \end{pgfmanualentry}
}


\def\pgfmanualbar{\char`\|}
\makeatletter
\newenvironment{pathoperation}[3][]{
  \begin{pgfmanualentry}
    \pgfmanualentryheadline{\textcolor{gray}{{\ttfamily\char`\\path}\
        \ \dots}
      \declare{\texttt{#2}}#3\ \textcolor{gray}{\dots\texttt{;}}}%
    \def\pgfmanualtest{#1}%
    \ifx\pgfmanualtest\@empty%
      \index{#2@\protect\texttt{#2} path operation}%
      \index{Path operations!#2@\protect\texttt{#2}}%
      \pgfmanualpdflabel{#2}{}%
    \fi%
    \pgfmanualbody
}
{
  \end{pgfmanualentry}
}
\makeatother

\def\extractcommand#1#2\@@{%
  \pgfmanualentryheadline{\declare{\texttt{\string#1}}#2}%
  \removeats{#1}%
  \index{\strippedat @\protect\myprintocmmand{\strippedat}}}

\def\doublebs{\texttt{\char`\\\char`\\}}


\newenvironment{package}[1]{
  \begin{pgfmanualentry}
    \pgfmanualentryheadline{{\ttfamily\char`\\usepackage\char`\{\declare{#1}\char`\}\space\space \char`\%\space\space  \LaTeX}}
    \index{#1@\protect\texttt{#1} package}%
    \index{Packages and files!#1@\protect\texttt{#1}}%
    \pgfmanualentryheadline{{\ttfamily\char`\\input \declare{#1}.tex\space\space\space \char`\%\space\space  plain \TeX}}
    \pgfmanualentryheadline{{\ttfamily\char`\\usemodule[\declare{#1}]\space\space \char`\%\space\space  Con\TeX t}}
    \pgfmanualbody
}
{
  \end{pgfmanualentry}
}


\newenvironment{pgfmodule}[1]{
  \begin{pgfmanualentry}
    \pgfmanualentryheadline{{\ttfamily\char`\\usepgfmodule\char`\{\declare{#1}\char`\}\space\space\space
        \char`\%\space\space  \LaTeX\space and plain \TeX\space and pure pgf}}
    \index{#1@\protect\texttt{#1} module}%
    \index{Modules!#1@\protect\texttt{#1}}%
    \pgfmanualentryheadline{{\ttfamily\char`\\usepgfmodule[\declare{#1}]\space\space \char`\%\space\space  Con\TeX t\space and pure pgf}}
    \pgfmanualbody
}
{
  \end{pgfmanualentry}
}

\newenvironment{pgflibrary}[1]{
  \begin{pgfmanualentry}
    \pgfmanualentryheadline{{\ttfamily\char`\\usepgflibrary\char`\{\declare{#1}\char`\}\space\space\space
        \char`\%\space\space  \LaTeX\space and plain \TeX\space and pure pgf}}
    \index{#1@\protect\texttt{#1} library}%
    \index{Libraries!#1@\protect\texttt{#1}}%
    \pgfmanualentryheadline{{\ttfamily\char`\\usepgflibrary[\declare{#1}]\space\space \char`\%\space\space  Con\TeX t\space and pure pgf}}
    \pgfmanualentryheadline{{\ttfamily\char`\\usetikzlibrary\char`\{\declare{#1}\char`\}\space\space
        \char`\%\space\space  \LaTeX\space and plain \TeX\space when using \tikzname}}
    \pgfmanualentryheadline{{\ttfamily\char`\\usetikzlibrary[\declare{#1}]\space
        \char`\%\space\space  Con\TeX t\space when using \tikzname}}
    \pgfmanualbody
}
{
  \end{pgfmanualentry}
}

\newenvironment{tikzlibrary}[1]{
  \begin{pgfmanualentry}
    \pgfmanualentryheadline{{\ttfamily\char`\\usetikzlibrary\char`\{\declare{#1}\char`\}\space\space \char`\%\space\space  \LaTeX\space and plain \TeX}}
    \index{#1@\protect\texttt{#1} library}%
    \index{Libraries!#1@\protect\texttt{#1}}%
    \pgfmanualentryheadline{{\ttfamily\char`\\usetikzlibrary[\declare{#1}]\space \char`\%\space\space Con\TeX t}}
    \pgfmanualbody
}
{
  \end{pgfmanualentry}
}



\newenvironment{filedescription}[1]{
  \begin{pgfmanualentry}
    \pgfmanualentryheadline{文件 {\ttfamily\declare{#1}}}%
    \index{#1@\protect\texttt{#1} file}%
    \index{Packages and files!#1@\protect\texttt{#1}}%
    \pgfmanualbody
}
{
  \end{pgfmanualentry}
}


\newenvironment{packageoption}[1]{
  \begin{pgfmanualentry}
    \pgfmanualentryheadline{{\ttfamily\char`\\usepackage[\declare{#1}]\char`\{pgf\char`\}}}
    \index{#1@\protect\texttt{#1} package option}%
    \index{Package options for \textsc{pgf}!#1@\protect\texttt{#1}}%
    \pgfmanualbody
}
{
  \end{pgfmanualentry}
}



\newcommand\opt[1]{{\color{black!50!green}#1}}
\newcommand\ooarg[1]{{\ttfamily[}\meta{#1}{\ttfamily]}}

\def\opt{\afterassignment\pgfmanualopt\let\next=}
\def\pgfmanualopt{\ifx\next\bgroup\bgroup\color{black!50!green}\else{\color{black!50!green}\next}\fi}



\def\beamer{\textsc{beamer}}
\def\pdf{\textsc{pdf}}
\def\pgfname{\textsc{pgf}}
\def\tikzname{Ti\emph{k}Z}
\def\pstricks{\textsc{pstricks}}
\def\prosper{\textsc{prosper}}
\def\seminar{\textsc{seminar}}
\def\texpower{\textsc{texpower}}
\def\foils{\textsc{foils}}

{
  \makeatletter
  \global\let\myempty=\@empty
  \global\let\mygobble=\@gobble
  \catcode`\@=12
  \gdef\getridofats#1@#2\relax{%
    \def\getridtest{#2}%
    \ifx\getridtest\myempty%
      \expandafter\def\expandafter\strippedat\expandafter{\strippedat#1}
    \else%
      \expandafter\def\expandafter\strippedat\expandafter{\strippedat#1\protect\printanat}
      \getridofats#2\relax%
    \fi%
  }

  \gdef\removeats#1{%
    \let\strippedat\myempty%
    \edef\strippedtext{\stripcommand#1}%
    \expandafter\getridofats\strippedtext @\relax%
  }

  \gdef\stripcommand#1{\expandafter\mygobble\string#1}
}

\def\printanat{\char`\@}

\def\declare{\afterassignment\pgfmanualdeclare\let\next=}
\def\pgfmanualdeclare{\ifx\next\bgroup\bgroup\color{red!75!black}\else{\color{red!75!black}\next}\fi}


\let\textoken=\command
\let\endtextoken=\endcommand

\def\myprintocmmand#1{\texttt{\char`\\#1}}

\def\example{\par\smallskip\noindent\textit{例： }}
\def\themeauthor{\par\smallskip\noindent\textit{主题作者： }}


\def\indexoption#1{%
  \index{#1@\protect\texttt{#1} option}%
  \index{Graphic options and styles!#1@\protect\texttt{#1}}%
}

\def\itemcalendaroption#1{\item \declare{\texttt{#1}}%
  \index{#1@\protect\texttt{#1} date test}%
  \index{Date tests!#1@\protect\texttt{#1}}%
}



\def\class#1{\list{}{\leftmargin=2em\itemindent-\leftmargin\def\makelabel##1{\hss##1}}%
\extractclass#1@\par\topsep=0pt}
\def\endclass{\endlist}
\def\extractclass#1#2@{%
\item{{{\ttfamily\char`\\documentclass}#2{\ttfamily\char`\{\declare{#1}\char`\}}}}%
  \index{#1@\protect\texttt{#1} class}%
  \index{Classes!#1@\protect\texttt{#1}}}

\def\partname{部分}

\makeatletter
\def\index@prologue{\section*{索引}\addcontentsline{toc}{section}{索引}
  本索引仅包含自动生成的条目。一份好的索引还应包含经过仔细挑选的关键词。本索引并不是一份好的索引。
  \bigskip
}
\c@IndexColumns=2
  \def\theindex{\@restonecoltrue
    \columnseprule \z@  \columnsep 29\p@
    \twocolumn[\index@prologue]%
       \parindent -30pt
       \columnsep 15pt
       \parskip 0pt plus 1pt
       \leftskip 30pt
       \rightskip 0pt plus 2cm
       \small
       \def\@idxitem{\par}%
    \let\item\@idxitem \ignorespaces}
  \def\endtheindex{\onecolumn}
\def\noindexing{\let\index=\@gobble}



\newcommand\symarrow[1]{
  \index{#1@\protect\texttt{#1} arrow tip}%
  \index{Arrow tips!#1@\protect\texttt{#1}}
  \texttt{#1}& 产生粗线
  \begin{tikzpicture}[arrows={#1-#1},thick,baseline]
    \useasboundingbox (0pt,-0.5ex) rectangle (1cm,2ex);
    \draw (0pt,.5ex) -- (1cm,.5ex);
  \end{tikzpicture} 和细线
  \begin{tikzpicture}[arrows={#1-#1},thin,baseline]
    \useasboundingbox (0pt,-0.5ex) rectangle (1cm,2ex);
    \draw (0pt,.5ex) -- (1cm,.5ex);
  \end{tikzpicture}
}
\newcommand\symarrowdouble[1]{
  \index{#1@\protect\texttt{#1} arrow tip}%
  \index{Arrow tips!#1@\protect\texttt{#1}}
  \texttt{#1}& 产生粗线
  \begin{tikzpicture}[arrows={#1-#1},thick,baseline]
    \useasboundingbox (0pt,-0.5ex) rectangle (1cm,2ex);
    \draw (0pt,.5ex) -- (1cm,.5ex);
  \end{tikzpicture}
  和细线
  \begin{tikzpicture}[arrows={#1-#1},thin,baseline]
    \useasboundingbox (0pt,-0.5ex) rectangle (1cm,2ex);
    \draw (0pt,.5ex) -- (1cm,.5ex);
  \end{tikzpicture}，双线
  \begin{tikzpicture}[arrows={#1-#1},thick,baseline]
    \useasboundingbox (0pt,-0.5ex) rectangle (1cm,2ex);
    \draw[double,double equal sign distance] (0pt,.5ex) -- (1cm,.5ex);
  \end{tikzpicture} 和
  \begin{tikzpicture}[arrows={#1-#1},thin,baseline]
    \useasboundingbox (0pt,-0.5ex) rectangle (1cm,2ex);
    \draw[double,double equal sign distance] (0pt,.5ex) -- (1cm,.5ex);
  \end{tikzpicture}
}

\newcommand\sarrow[2]{
  \index{#1@\protect\texttt{#1} arrow tip}%
  \index{Arrow tips!#1@\protect\texttt{#1}}
  \index{#2@\protect\texttt{#2} arrow tip}%
  \index{Arrow tips!#2@\protect\texttt{#2}}
  \texttt{#1-#2}& 产生粗线
  \begin{tikzpicture}[arrows={#1-#2},thick,baseline]
    \useasboundingbox (0pt,-0.5ex) rectangle (1cm,2ex);
    \draw (0pt,.5ex) -- (1cm,.5ex);
  \end{tikzpicture} 和细线
  \begin{tikzpicture}[arrows={#1-#2},thin,baseline]
    \useasboundingbox (0pt,-0.5ex) rectangle (1cm,2ex);
    \draw (0pt,.5ex) -- (1cm,.5ex);
  \end{tikzpicture}
}

\newcommand\sarrowdouble[2]{
  \index{#1@\protect\texttt{#1} arrow tip}%
  \index{Arrow tips!#1@\protect\texttt{#1}}
  \index{#2@\protect\texttt{#2} arrow tip}%
  \index{Arrow tips!#2@\protect\texttt{#2}}
  \texttt{#1-#2}& 产生粗线
  \begin{tikzpicture}[arrows={#1-#2},thick,baseline]
    \useasboundingbox (0pt,-0.5ex) rectangle (1cm,2ex);
    \draw (0pt,.5ex) -- (1cm,.5ex);
  \end{tikzpicture} 和细线
  \begin{tikzpicture}[arrows={#1-#2},thin,baseline]
    \useasboundingbox (0pt,-0.5ex) rectangle (1cm,2ex);
    \draw (0pt,.5ex) -- (1cm,.5ex);
  \end{tikzpicture}，双线
  \begin{tikzpicture}[arrows={#1-#2},thick,baseline]
    \useasboundingbox (0pt,-0.5ex) rectangle (1cm,2ex);
    \draw[double,double equal sign distance] (0pt,.5ex) -- (1cm,.5ex);
  \end{tikzpicture} 和
  \begin{tikzpicture}[arrows={#1-#2},thin,baseline]
    \useasboundingbox (0pt,-0.5ex) rectangle (1cm,2ex);
    \draw[double,double equal sign distance] (0pt,.5ex) -- (1cm,.5ex);
  \end{tikzpicture}
}

\newcommand\carrow[1]{
  \index{#1@\protect\texttt{#1} arrow tip}%
  \index{Arrow tips!#1@\protect\texttt{#1}}
  \texttt{#1}& 对线宽 1ex 产生
  \begin{tikzpicture}[arrows={#1-#1},line width=1ex,baseline]
    \useasboundingbox (0pt,-0.5ex) rectangle (1.5cm,2ex);
    \draw (0pt,.5ex) -- (1.5cm,.5ex);
  \end{tikzpicture}
}
\def\myvbar{\char`\|}
\newcommand\plotmarkentry[1]{%
  \index{#1@\protect\texttt{#1} plot mark}%
  \index{Plot marks!#1@\protect\texttt{#1}}
  \texttt{\char`\\pgfuseplotmark\char`\{\declare{#1}\char`\}} &
  \tikz\draw[color=black!25] plot[mark=#1,mark options={fill=examplefill,draw=black}] coordinates{(0,0) (.5,0.2) (1,0) (1.5,0.2)};\\
}
\newcommand\plotmarkentrytikz[1]{%
  \index{#1@\protect\texttt{#1} plot mark}%
  \index{Plot marks!#1@\protect\texttt{#1}}
  \texttt{mark=\declare{#1}} & \tikz\draw[color=black!25]
  plot[mark=#1,mark options={fill=examplefill,draw=black}]
    coordinates {(0,0) (.5,0.2) (1,0) (1.5,0.2)};\\
}



\ifx\scantokens\@undefined
  \PackageError{pgfmanual-macros}{You need to use extended latex
    (elatex) or (pdfelatex) to process this document}{}
\fi

% Define \find@example such that it doesn't destroy catcodes:
\begingroup
\catcode`|=0
\catcode`[= 1
\catcode`]=2
\catcode`\{=12
\catcode `\}=12
\catcode`\\=12 |gdef|find@example#1\end{codeexample}[|endofcodeexample[#1]]
|endgroup

% define \returntospace.
%
% It should define NEWLINE as {}, spaces and tabs as \space.
\begingroup
\catcode`\^=7
\catcode`\^^M=13
\catcode`\^^I=13
\catcode`\ =13%
\gdef\returntospace{\catcode`\ =13\def {\space}\catcode`\^^I=13\def^^I{\space}\catcode`\^^M=13\def^^M{\pgfkeysvalueof{/codeexample/newline}}}%
\endgroup

\begingroup
\catcode`\%=13
\catcode`\^^M=13
\gdef\commenthandler{\catcode`\%=13\def%{\@gobble@till@return}}
\gdef\@gobble@till@return#1^^M{}
\gdef\@gobble@till@return@ignore#1^^M{\ignorespaces}
\gdef\typesetcomment{\catcode`\%=13\def%{\@typeset@till@return}}
\gdef\@typeset@till@return#1^^M{{\def%{\char`\%}\pgfkeysvalueof{/codeexample/comment font}{\char`\%#1}}\par}
\endgroup

% Define tab-implementation functions
%   \codeexample@tabinit@replacementchars@
% and
%   \codeexample@tabinit@catcode@
%
% They should ONLY be used in case that tab replacement is active.
%
% This here is merely a preparation step.
%
% Idea:
% \codeexample@tabinit@catcode@ will make TAB active
% and
% \codeexample@tabinit@replacementchars@ will insert as many spaces as
% /codeexample/tabsize contains.
{
\catcode`\^^I=13
% ATTENTION: do NOT use tabs in these definitions!!
\gdef\codeexample@tabinit@replacementchars@{%
 \begingroup
 \count0=\pgfkeysvalueof{/codeexample/tabsize}\relax
 \toks0={}%
 \loop
 \ifnum\count0>0
  \advance\count0 by-1
  \toks0=\expandafter{\the\toks0\ }%
 \repeat
 \xdef\codeexample@tabinit@replacementchars@@{\the\toks0}%
 \endgroup
 \let^^I=\codeexample@tabinit@replacementchars@@
}%
\gdef\codeexample@tabinit@catcode@{\catcode`\^^I=13}%
}%

% Called after any options have been set. It assigns
%   \codeexample@tabinit@catcode
% and
%   \codeexample@tabinit@replacementchars
% which are used inside of
%\begin{codeexample}
% ...
%\end{codeexample}
%
% \codeexample@tabinit@catcode  is either \relax or it makes tab
% active.
%
% \codeexample@tabinit@replacementchars is either \relax or it inserts
% a proper replacement sequence for tabs (as many spaces as
% configured)
\def\codeexample@tabinit{%
    \ifnum\pgfkeysvalueof{/codeexample/tabsize}=0\relax
        \let\codeexample@tabinit@replacementchars=\relax
        \let\codeexample@tabinit@catcode=\relax
    \else
        \let\codeexample@tabinit@catcode=\codeexample@tabinit@catcode@
        \let\codeexample@tabinit@replacementchars=\codeexample@tabinit@replacementchars@
    \fi
}

\pgfqkeys{/codeexample}{%
    width/.code=    {\setlength\codeexamplewidth{#1}},
    graphic/.code=    {\colorlet{graphicbackground}{#1}},
    code/.code=    {\colorlet{codebackground}{#1}},
    execute code/.is if=code@execute,
    code only/.code=    {\code@executefalse},
    pre/.code=    {\def\code@pre{#1}},
    post/.code=    {\def\code@post{#1}},
    % #1 is the *complete* environment contents as it shall be
    % typeset. In particular, the catcodes are NOT the normal ones.
    typeset listing/.code=    {#1},
    vbox/.code=    {\def\code@pre{\vbox\bgroup\setlength{\hsize}{\linewidth-6pt}}\def\code@post{\egroup}},
    ignorespaces/.code=    {\let\@gobble@till@return=\@gobble@till@return@ignore},
    leave comments/.code=    {\def\code@catcode@hook{\catcode`\%=12}\let\commenthandler=\relax\let\typesetcomment=\relax},
    comment font/.initial=\textsl,% FIXME : doesn't work!?
    tabsize/.initial=0,
    newline/.initial=^^J,% use 'newline=^^J' to insert a *real* newline
    every codeexample/.style={width=4cm+7pt},
}

\def\code@pre{}
\def\code@post{}
\def\code@catcode@hook{}

\newdimen\codeexamplewidth
\newif\ifcode@execute
\newbox\codeexamplebox
\def\codeexample[#1]{%
  \begingroup%
  \code@executetrue
  \pgfqkeys{/codeexample}{every codeexample,#1}%
  \codeexample@tabinit% assigns \codeexample@tabinit@[catcode,replacementchars]
  \parindent0pt
  \begingroup%
  \par%
  \medskip%
  \let\do\@makeother%
  \dospecials%
  \obeylines%
  \@vobeyspaces%
  \catcode`\%=13%
  \catcode`\^^M=13%
  \code@catcode@hook%
  \codeexample@tabinit@catcode
  \relax%
  \find@example}
\def\endofcodeexample#1{%
  \endgroup%
  \ifcode@execute%
    \setbox\codeexamplebox=\hbox{%
      {%
        {%
          \returntospace%
          \commenthandler%
          \xdef\code@temp{#1}% removes returns and comments
        }%
%\message{Nach returntospace: \meaning\code@temp^^J^^J^^J}%
        \expandafter\pgfmanual@codeexample@generate@graphic\expandafter{\code@temp}%
      }%
    }%
    \ifdim\wd\codeexamplebox>\codeexamplewidth%
      \def\code@start{\par}%
      \def\code@flushstart{}\def\code@flushend{}%
      \def\code@mid{\parskip2pt\par\noindent}%
      \def\code@width{\linewidth-6pt}%
      \def\code@end{}%
    \else%
      \def\code@start{%
        \linewidth=\textwidth%
        \parshape \@ne 0pt \linewidth
        \leavevmode%
        \hbox\bgroup}%
      \def\code@flushstart{\hfill}%
      \def\code@flushend{\hbox{}}%
      \def\code@mid{\hskip6pt}%
      \def\code@width{\linewidth-12pt-\codeexamplewidth}%
      \def\code@end{\egroup}%
    \fi%
    \code@start%
    \noindent%
    \begin{minipage}[t]{\codeexamplewidth}\raggedright
      \hrule width0pt%
      \footnotesize\vskip-1em%
      \code@flushstart\box\codeexamplebox\code@flushend%
      \vskip-1ex
      \leavevmode%
    \end{minipage}%
  \else%
    \def\code@mid{\par}
    \def\code@width{\linewidth-6pt}
    \def\code@end{}
  \fi%
  \code@mid%
  \colorbox{codebackground}{%
    \begin{minipage}[t]{\code@width}%
      {%
        \let\do\@makeother
        \dospecials
        \frenchspacing\@vobeyspaces
        \normalfont\ttfamily\footnotesize
        \typesetcomment%
        \codeexample@tabinit@replacementchars
        \@tempswafalse
        \def\par{%
          \if@tempswa
          \leavevmode \null \@@par\penalty\interlinepenalty
          \else
          \@tempswatrue
          \ifhmode\@@par\penalty\interlinepenalty\fi
          \fi}%
        \obeylines
        \everypar \expandafter{\the\everypar \unpenalty}%
        \pgfmanual@@typeset@listing{#1}%
    }
    \end{minipage}}%
  \code@end%
  \par%
  \medskip
  \end{codeexample}
}

% the purpose of these extra-macros is to plug-in special tex4ht
% instructions for the html version:
\long\def\pgfmanual@codeexample@generate@graphic#1{%
    \colorbox{graphicbackground}{\pgfmanual@codeexample@generate@graphic@{#1}}%
}%
\long\def\pgfmanual@codeexample@generate@graphic@#1{%
    \color{black}\ignorespaces%
    \code@pre%
    \pgfmanual@scantokens{#1}\ignorespaces%
    \code@post\ignorespaces%
}%

\ifluatex
    % LUA has a couple of bugs around \scantokens, see
    % http://tex.stackexchange.com/questions/197634/workaround-for-open-lualatex-bug-in-scantokens
    \def\pgfmanual@scantokens#1{
        \directlua{%
            local s = "\luaescapestring{\code@temp}"
            local splitted = pgfplots.stringsplit(s, "\string\n")
            for i =1,\pgfplotsHASH splitted do
                local line=splitted[i]
              tex.print(line)
            end
        }%
        %\def\pgfmanual@tempfile{\jobname.code.tmp}%
        %\immediate\openout\w@pgf@writea=\pgfmanual@tempfile\relax
        %\immediate\write\w@pgf@writea{#1}%
        %\immediate\closeout\w@pgf@writea
        %\input{\pgfmanual@tempfile}
    }%
\else
    \let\pgfmanual@scantokens=\scantokens
\fi

\long\def\pgfmanual@@typeset@listing#1{%
    \pgfkeysvalueof{/codeexample/typeset listing/.@cmd}{#1}\pgfeov
}%%

\def\endcodeexample{\endgroup}


\makeatother


%%% Local Variables:
%%% mode: latex
%%% TeX-master: "beameruserguide"
%%% End:



\def\pgfplotsifdocpackageuptodate#1#2{%
    \pgfkeysifdefined{/codeexample/prettyprint/word/.@cmd}{#1}{#2}
}%

\pgfplotsiffileexists{pgfmanual.sty}{%
    \RequirePackage{pgfmanual}
    \pgfplotsifdocpackageuptodate{}{%
        \makeatletter
        \input{pgfplotsoldpgfsupp_pgfmanual.code.tex}
        \makeatother
    }%
}{%
    \makeatletter
    \input{pgfplotsoldpgfsupp_pgfmanual.code.tex}
    \makeatother
}%

\makeatletter
\def\pgfplotsmakefilelinkifuseful#1#2{%
    \protect\pgfplotsmakefilelinkifuseful@{#1}{#2}%
}%
\def\pgfplotsmakefilelinkifuseful@#1#2{%
    \edef\temp{#1}%
    \edef\tempb{\jobname}%
    \edef\temp{\meaning\temp}% \meaning normalizes the catcodes.
    \edef\tempb{\meaning\tempb}%
    \ifx\temp\tempb
        % we are processing '#1'. Don't make a link.
        #2%
    \else
        \href{file:#1.pdf}{#2}%
    \fi
}%
\makeatother


\pgfkeys{
    /codeexample/prettyprint/cs arguments/pgfplotscreateplotcyclelist/.initial=2,
    /codeexample/prettyprint/cs/pgfplotscreateplotcyclelist/.code args={#1#2#3}{\pgfmanualpdfref{#1}{#1}\{#2\}\{\pgfmanualprettyprintpgfkeys{#3}\pgfmanualclosebrace},
    /codeexample/prettyprint/cs arguments/tikzset/.initial=1,
    /codeexample/prettyprint/cs/tikzset/.code 2 args={\pgfmanualpdfref{#1}{#1}\{\pgfmanualprettyprintpgfkeys{#2}\pgfmanualclosebrace},
    /codeexample/prettyprint/cs arguments/pgfplotsset/.initial=1,
    /codeexample/prettyprint/cs/pgfplotsset/.code 2 args={\pgfmanualpdfref{#1}{#1}\{\pgfmanualprettyprintpgfkeys{#2}\pgfmanualclosebrace},
    /codeexample/prettyprint/cs arguments/pgfplotstableset/.initial=1,
    /codeexample/prettyprint/cs/pgfplotstableset/.code 2 args={\pgfmanualpdfref{#1}{#1}\{\pgfmanualprettyprintpgfkeys{#2}\pgfmanualclosebrace},
    /codeexample/prettyprint/cs arguments/usepgfplotslibrary/.initial=1,
    /codeexample/prettyprint/cs/usepgfplotslibrary/.code 2 args={\pgfmanualpdfref{#1}{#1}\{\pgfmanualpdfref{#2}{#2}\pgfmanualclosebrace},
    %
    %
    %/codeexample/prettyprint/key value/cycle list/.code 2 args={\pgfmanualprettyprintpgfkeys{#2}},
    /codeexample/prettyprint/key value/xticklabel/.code 2 args={\pgfmanualprettyprintcode{#2}},
    /codeexample/prettyprint/key value/yticklabel/.code 2 args={\pgfmanualprettyprintcode{#2}},
    /codeexample/prettyprint/key value/zticklabel/.code 2 args={\pgfmanualprettyprintcode{#2}},
    /codeexample/prettyprint/key value/nodes near coords/.code 2 args={\pgfmanualprettyprintcode{#2}},
    /codeexample/prettyprint/key value/nodes near coords*/.code 2 args={\pgfmanualprettyprintcode{#2}},
    %
    /codeexample/prettyprint/key value/includegraphics/.code 2 args={\pgfmanualprettyprintpgfkeys{#2}},
    /codeexample/prettyprint/key value/contour gnuplot/.code 2 args={\pgfmanualprettyprintpgfkeys{#2}},
    /codeexample/prettyprint/key value/contour lua/.code 2 args={\pgfmanualprettyprintpgfkeys{#2}},
    /codeexample/prettyprint/key value/contour prepared/.code 2 args={\pgfmanualprettyprintpgfkeys{#2}},
    /codeexample/prettyprint/key value/contour filled/.code 2 args={\pgfmanualprettyprintpgfkeys{#2}},
    /codeexample/prettyprint/key value/hist/.code 2 args={\pgfmanualprettyprintpgfkeys{#2}},
    /codeexample/prettyprint/key value/tieline/.code 2 args={\pgfmanualprettyprintpgfkeys{#2}},
    /codeexample/prettyprint/key value/quiver/.code 2 args={\pgfmanualprettyprintpgfkeys{#2}},
    /codeexample/prettyprint/key value/boxplot/.code 2 args={\pgfmanualprettyprintpgfkeys{#2}},
    /codeexample/prettyprint/key value/boxplot prepared/.code 2 args={\pgfmanualprettyprintpgfkeys{#2}},
    %
    /codeexample/prettyprint/key value/colormap/.code 2 args={\pgfmanualprettyprintpgfkeys{#2}},
    /codeexample/prettyprint/key value/of colormap/.code 2 args={\pgfmanualprettyprintpgfkeys{#2}},
    /codeexample/prettyprint/key value/samples of colormap/.code 2 args={\pgfmanualprettyprintpgfkeys{#2}},
    /codeexample/prettyprint/key value/levels from colormap/.code 2 args={\pgfmanualprettyprintpgfkeys{#2}},
    /codeexample/prettyprint/key value/colors of colormap/.code 2 args={\pgfmanualprettyprintpgfkeys{#2}},
    %
    %
    % whenever an unqualified key is found, the following key prefix
    % list is tried to find a match.
    /pdflinks/search key prefixes in={/pgfplots/table/,/pgfplots/error bars/,/pgfplots/,/pgfplots/plot file/,/tikz/,/pgf/},
    %
    % the link prefix written to the pdf file:
    /pdflinks/internal link prefix=pgfp,
    %
    /pdflinks/warnings=false,
    /pdflinks/codeexample links=true,
    /pdflinks/show labels=false,
}%


% should be used to show something in red which doesn't need to get a
% hyper ref.
%
% Examples are descriptions of key labels.
\def\declaretext#1{\texttt{\declare{#1}}}

% To be used whenever something NEW has been declared.
% In this case, a \pgfmanualpdflabel will be generated using '#1'.
%
% Use '\declaretext' if you only describe something local (for example
% the documentation of key values).
\def\declarelabel#1{%
    \texttt{\declare{#1}}%
    \pgfmanualpdflabel{#1}{}%
}

\def\pgfmanualbar{\char`\|}
\makeatletter

%%%%%%%%%%%%%%%%%%%%%%%%%%%%%%%%%%%%%%%%%%%%%%%%%%%%%%%%%%%%%%%%%%%%%%%%%%%%%%

\usepgfplotslibrary{external}

% use \pgfplotsmanualenableexternalizationofexpensive in the preamble
% to enable externalization of expensive examples.
% This will ONLY externalize expensive examples, i.e. those for which
% \pgfplotsexpensiveexample is written:
%
% \pgfplotsexpensiveexample
% \begin{codeexample}
\def\pgfplotsmanualenableexternalizationofexpensive{%
    \pgfplotsmanual@enable@externalization@for@expensivetrue
    \tikzexternalize[
        prefix=figures/expensiveexamples/,% the 'X' suffix is to avoid confusion in git: previous versions contained the pdfs in git
        figure name={},
        export=false, % needs to be activated for single pictures (i.e. expensive ones)
        mode=list and make,
        verbose=false,
        %xport=true,% FASTER FOR DEBUGGING
    ]
    \tikzifexternalizing{%
        \nofiles
        \pgfkeys{/pdflinks/codeexample links=false}%
    }{}%
}%
%\pgfkeys{/pgf/images/include external/.code={\href{file:#1}{\pgfimage{#1}}}} % FIXME : NOT FOR THE FINAL VERSION

\newif\ifpgfplotsmanual@enable@externalization@for@expensive
\newif\ifpgfplots@example@is@expensive

\pgfkeys{
    /codeexample/every codeexample/.append code={%
        \ifpgfplots@example@is@expensive
            \pgfkeys{/tikz/external/export=true}%
            \global\pgfplots@example@is@expensivefalse
        \fi
    }
}

% Write this macro directly in front of \begin{codeexample} (without arguments):
\def\pgfplotsexpensiveexample{%
    \ifpgfplotsmanual@enable@externalization@for@expensive
        \pgfplots@example@is@expensivetrue
    \else
        \message{[NOTE: I am now about to typeset an expensive example. You will need to ENLARGE YOUR TeX MEMORY CAPACITIES if this fails.]}%
    \fi
}%

%%%%%%%%%%%%%%%%%%%%%%%%%%%%%%%%%%%%%%%%%%%%%%%%%%%%%%%%%%%%%%%%%%%%%%%%%%%%%%

\newif\ifpgfplotsmanualhtmlmode

% MACROS FOR HTML OUTPUT:
% HTML will be compiled in a separate folder. See the Makefile
\tikzset{
    external/html export mode/.style={
        system call={%
        % FIXME : which converter to png ? Why did I use pdftoppm |
        % pnmtopng in the first draft?
            %pdflatex \tikzexternalcheckshellescape -halt-on-error -interaction=batchmode -jobname "\image" "\\def\\pgfplotsmanualhtlatexmode{1}\texsource" && pdftoppm "\image.pdf" | pnmtopng > "\image.png"%
            pdflatex \tikzexternalcheckshellescape -halt-on-error -interaction=batchmode -jobname "\image" "\\def\\pgfplotsmanualhtlatexmode{1}\texsource" && pdftoppm -png "\image.pdf" "\image"%
            %pdflatex \tikzexternalcheckshellescape -halt-on-error -interaction=batchmode -jobname "\image" "\\def\\pgfplotsmanualhtlatexmode{1}\texsource" && convert -density 150 "\image.pdf" "\image.png"%
        },
        /pgf/images/external info,
        /pgf/images/include external/.code={%
            \includegraphics
               [width=\pgfexternalwidth,height=\pgfexternalheight]
               {##1.png}%
        },
        prefix=figures/generated/,
        figure name={manual},
        mode=list and make,
        verbose=false,
    },
}

\def\pgfplotsmanual@configure@for@htlatexmode{%
    % disable this; we will externalize everything now:
    \let\pgfplotsmanualenableexternalizationofexpensive=\relax
    %
    \message{^^Jpgfplots manual: initializing external lib for HTML output (png)^^J}%
    \tikzexternalize[
        html export mode,
    ]
    %
    \@ifpackageloaded{tex4ht}{%
        %\usepackage[html,png,3]{tex4ht}
        \tikzset{external/mode=only graphics}%
        \pgfplotsmanual@configure@for@htlatexmode@
        %\patches@for@htlatex
    }{%
    }%
    %
    %
    \pgfplotsmanualhtmlmodetrue
    %
    \gdef\pgfsys@imagesuffixlist{.png}
    %
    \tikzifexternalizing{%
        %\nofiles
        \pgfkeys{/pdflinks/codeexample links=false}%
    }{}%
}

% explanation of how to customize tex4ht:
% http://www.cvr.cc/tex4ht-low-level-commands/#more-482
%
% some reference:
% http://www.cvr.cc/tex4ht-options/
%
% some tex4ht mailing list discussion:
% http://tug.org/pipermail/tex4ht/2010q4/000246.html
\def\patches@for@htlatex{%
 \let\savecolor\color
 \NewConfigure{color}[2]{\def\a at color{##1}\def\b at color{##2}}
 \def\@@tmp##1{\a at color##1\b at color\savecolor{##1}\aftergroup\endspan}
 \let\color\@@tmp
 \def\endspan{\Tg</span>}
 \Configure{color}{\HCode{<span style="color:}}{\HCode{;">}}
}%

\def\pgfplotsmanual@configure@for@htlatexmode@{%
    \NewConfigure{endofcodeexample}[2]{%
        \def\endofcodeexample@pre{##1}%
        \def\endofcodeexample@post{##2}%
    }%
    \let\endofcodeexample@orig=\endofcodeexample
    \def\endofcodeexample##1{%
        \endofcodeexample@pre
        \endofcodeexample@orig{##1}%
        \endofcodeexample@post
    }%
    \Configure{endofcodeexample}
        {\Tg<div class="img">}
        {\Tg</div>}%
    %
    \NewConfigure{pgfmanual@@typeset@listing}[2]{%
        \def\pgfmanual@@typeset@listing@pre{##1}%
        \def\pgfmanual@@typeset@listing@post{##2}%
    }%
    \let\pgfmanual@@typeset@listing@orig=\pgfmanual@@typeset@listing
    \def\pgfmanual@@typeset@listing##1{%
        \pgfmanual@@typeset@listing@pre
        \pgfmanual@@typeset@listing@orig{##1}%
        \pgfmanual@@typeset@listing@post
    }%
    \Configure{pgfmanual@@typeset@listing}
        {\Tg<div class="texsrc">}
        {\Tg</div>}%
    %
    %
    \NewConfigure{pgfmanual@codeexample@generate@graphic}[2]{%
        \def\pgfmanual@codeexample@generate@graphic@pre{##1}%
        \def\pgfmanual@codeexample@generate@graphic@post{##2}%
    }%
    \def\pgfmanual@codeexample@generate@graphic##1{%
        \pgfmanual@codeexample@generate@graphic@pre
        % the '@' variant skips the colorbox:
        \pgfmanual@codeexample@generate@graphic@{##1}%
        \pgfmanual@codeexample@generate@graphic@post
    }%
    \Configure{pgfmanual@codeexample@generate@graphic}
        {\Tg<span class="graphics">}
        {\Tg</span>}%
}

\@ifpackageloaded{tex4ht}{
    % always active here:
    \def\pgfplotsmanualhtlatexmode{1}%
}{
}

% the macro \pgfplotsmanualhtlatexmode should be set from command line
% (to any value). If it is known, the manual will be translated to
% HTML.
\pgfutil@ifundefined{pgfplotsmanualhtlatexmode}{%
}{%
    \pgfplotsmanual@configure@for@htlatexmode
}%


%%%%%%%%%%%%%%%%%%%%%%%%%%%%%%%%%%%%%%%%%%%%%%%%%%%%%%%%%%%%%%%%%%%%%%%%%%%%%%



\newenvironment{addplotoperation}[3][]{
  \begin{pgfmanualentry}
      {%
    \let\ltxdoc@marg=\marg
    \let\ltxdoc@oarg=\oarg
    \let\ltxdoc@parg=\parg
    \let\ltxdoc@meta=\meta
    \def\marg##1{{\normalfont\ltxdoc@marg{##1}}}%
    \def\oarg##1{{\normalfont\ltxdoc@oarg{##1}}}%
    \def\parg##1{{\normalfont\ltxdoc@parg{##1}}}%
    \def\meta##1{{%
    \ifmmode
        \mathit{##1}%
    \else
        \normalfont\ltxdoc@meta{##1}%
    \fi
}}%
    \pgfmanualentryheadline{\textcolor{gray}{{\ttfamily\char`\\addplot\ }}%
      \declare{\texttt{#2}} \texttt{#3;}}%
      \unskip
     \nobreak
    \pgfmanualentryheadline{\textcolor{gray}{\texttt{\char`\\addplot}\oarg{options} }%
      \declare{\texttt{#2}} \texttt{#3} \textcolor{gray}{\meta{trailing path commands}}\texttt{;}}%
      \unskip
     \nobreak
    \pgfmanualentryheadline{\textcolor{gray}{{\ttfamily\char`\\addplot3}} $\dotsc$}%
    \def\pgfmanualtest{#1}%
    \ifx\pgfmanualtest\@empty%
      \index{#2@\protect\textcolor{gray}{\protect\texttt{plot}}\protect\texttt{ #2}}%
      \index{Plot operations!plot #2@\protect\texttt{plot #2}}%
    \fi%
    \pgfmanualpdflabel{\textbackslash addplot #2}{}%
    \pgfmanualpdflabel{plot #2}{}%
    \pgfmanualpdflabel{#2}{}%
    }%
    \pgfmanualbody
}
{
  \end{pgfmanualentry}
}

\newenvironment{addplot+}{
  \begin{pgfmanualentry}
      {%
    \let\ltxdoc@marg=\marg
    \let\ltxdoc@oarg=\oarg
    \let\ltxdoc@parg=\parg
    \let\ltxdoc@meta=\meta
    \def\marg##1{{\normalfont\ltxdoc@marg{##1}}}%
    \def\oarg##1{{\normalfont\ltxdoc@oarg{##1}}}%
    \def\parg##1{{\normalfont\ltxdoc@parg{##1}}}%
    \def\meta##1{{%
    \ifmmode
        \mathit{##1}%
    \else
        \normalfont\ltxdoc@meta{##1}%
    \fi
}}%
    \pgfmanualentryheadline{{\ttfamily\declare{\char`\\addplot+}\oarg{options} \textcolor{gray}{\dots};}}%
    \index{addplot+@\protect\texttt{\protect\textbackslash addplot+}}%
    \pgfmanualpdflabel{\textbackslash addplot+}{}%
    }%
    \pgfmanualbody
}
{
  \end{pgfmanualentry}
}
\newenvironment{addplot3generic}{
  \begin{pgfmanualentry}
      {%
    \let\ltxdoc@marg=\marg
    \let\ltxdoc@oarg=\oarg
    \let\ltxdoc@parg=\parg
    \let\ltxdoc@meta=\meta
    \def\marg##1{{\normalfont\ltxdoc@marg{##1}}}%
    \def\oarg##1{{\normalfont\ltxdoc@oarg{##1}}}%
    \def\parg##1{{\normalfont\ltxdoc@parg{##1}}}%
    \def\meta##1{{\normalfont\ltxdoc@meta{##1}}}%
    \pgfmanualentryheadline{{\ttfamily\declare{\char`\\addplot3}\oarg{options} \meta{input data} \meta{trailing path commands};}}%
    \index{addplot3@\protect\texttt{\protect\textbackslash addplot3}}%
    \pgfmanualpdflabel{\textbackslash addplot3}{}%
    }%
    \pgfmanualbody
}
{
  \end{pgfmanualentry}
}
\newenvironment{addplot3operation}[3][]{
  \begin{pgfmanualentry}
      {%
    \let\ltxdoc@marg=\marg
    \let\ltxdoc@oarg=\oarg
    \let\ltxdoc@parg=\parg
    \let\ltxdoc@meta=\meta
    \def\marg##1{{\normalfont\ltxdoc@marg{##1}}}%
    \def\oarg##1{{\normalfont\ltxdoc@oarg{##1}}}%
    \def\parg##1{{\normalfont\ltxdoc@parg{##1}}}%
    \def\meta##1{{\normalfont\ltxdoc@meta{##1}}}%
    \pgfmanualentryheadline{\textcolor{gray}{{\ttfamily\char`\\addplot3\ }}%
      \declare{\texttt{#2}} \texttt{#3;}}%
      \unskip
     \nobreak
    \pgfmanualentryheadline{\textcolor{gray}{\texttt{\char`\\addplot3}\oarg{options} }%
      \declare{\texttt{#2}} \texttt{#3} \textcolor{gray}{\meta{trailing path commands}}\texttt{;}}%
    \def\pgfmanualtest{#1}%
    \ifx\pgfmanualtest\@empty%
      \index{#2@\protect\texttt{#2}}%
      \index{Plot operations!addplot3 #2@\protect\texttt{#2}}%
    \fi%
    \pgfmanualpdflabel{\textbackslash addplot3 #2}{}%
    \pgfmanualpdflabel{plot3 #2}{}%
    }%
    \pgfmanualbody
}
{
  \end{pgfmanualentry}
}

\newenvironment{codekey}[1]{%
  \begin{pgfmanualentry}
     \pgfmanualentryheadline{{\ttfamily\declarekey{#1}\textcolor{gray}{/\pgfmanualpdfref{/handlers/.code}{.code}}=\marg{...}}\hfill}%
    \def\mykey{#1}%
    \def\mypath{}%
    \def\myname{}%
    \firsttimetrue%
    \decompose#1/\nil%
    \pgfmanualbody
}
{
  \end{pgfmanualentry}
}
\newenvironment{codeargskey}[2]{%
  \begin{pgfmanualentry}
      {\toks0={#2}%
      \xdef\argpattern{\the\toks0 }%
    }%
      \pgfmanual@command@to@string\argpattern\argpattern
     \pgfmanualentryheadline{{\ttfamily\declarekey{#1}\textcolor{gray}{/\pgfmanualpdfref{/handlers/.code}{.code args}}=\texttt{\{\argpattern\}}\marg{...}}\hfill}%
    \def\mykey{#1}%
    \def\mypath{}%
    \def\myname{}%
    \firsttimetrue%
    \decompose#1/\nil%
    \pgfmanualbody
}
{
  \end{pgfmanualentry}
}
\def\pgfmanual@command@to@string#1#2{%
    \expandafter\pgfmanual@command@to@string@@\meaning#1\pgfmanual@EOI{#2}%
}%
\xdef\pgfmanual@glob@TMPa{\meaning\pgfutil@empty}%
\expandafter\def\expandafter\pgfmanual@command@to@string@@\pgfmanual@glob@TMPa#1\pgfmanual@EOI#2{%
    \def#2{#1}%
}%

\newenvironment{pgfplotscodekey}[1]{%
    \begin{codekey}{/pgfplots/#1}%
}
{
  \end{codekey}
}
\newenvironment{pgfplotscodetwokey}[1]{%
  \begin{pgfmanualentry}
     \pgfmanualentryheadline{{\ttfamily\declarekey{/pgfplots/#1}\textcolor{gray}{/\pgfmanualpdfref{/handlers/.code 2 args}{.code 2 args}}=\marg{...}}\hfill}%
    \def\mykey{/pgfplots/#1}%
    \def\mypath{}%
    \def\myname{}%
    \firsttimetrue%
    \decompose/pgfplots/#1/\nil%
    \pgfmanualbody
}
{
  \end{pgfmanualentry}
}

\newenvironment{pgfplotsxycodekeylist}[1]{%
    \begingroup
    \let\oldpgfmanualentryheadline=\pgfmanualentryheadline
    \def\pgfmanualentryheadline##1{%
        \pgfmanualentryheadline@##1\pgfplots@EOI
    }%
    \def\pgfmanualentryheadline@##1\hfill##2\pgfplots@EOI{%
        \oldpgfmanualentryheadline{{\ttfamily\declarekey{##1}\textcolor{gray}{/\pgfmanualpdfref{/handlers/.code}{.code}}=\marg{...}}\hfill}%
    }
    \begin{pgfplotsxykeylist}{#1}%
}
{
    \end{pgfplotsxykeylist}
    \endgroup
}

\newenvironment{pgfplotsxyexpressionkeylist}[1]{%
    \begingroup
    \let\oldpgfmanualentryheadline=\pgfmanualentryheadline
    \def\pgfmanualentryheadline##1{%
        \pgfmanualentryheadline@##1\pgfplots@EOI
    }%
    \def\pgfmanualentryheadline@##1\hfill##2\pgfplots@EOI{%
        \oldpgfmanualentryheadline{{\ttfamily\declarekey{##1}\textcolor{gray}{/.expression}=\marg{math expression}}\hfill}%
    }
    \begin{pgfplotsxykeylist}{#1}%
}
{
    \end{pgfplotsxykeylist}
    \endgroup
}
\newenvironment{pgfplotsxyappendexpressionkeylist}[1]{%
    \begingroup
    \let\oldpgfmanualentryheadline=\pgfmanualentryheadline
    \def\pgfmanualentryheadline##1{%
        \pgfmanualentryheadline@##1\pgfplots@EOI
    }%
    \def\pgfmanualentryheadline@##1\hfill##2\pgfplots@EOI{%
        \oldpgfmanualentryheadline{{\ttfamily\declarekey{##1}\textcolor{gray}{/.append expression}=\marg{math expression}}\hfill}%
    }
    \begin{pgfplotsxykeylist}{#1}%
}
{
    \end{pgfplotsxykeylist}
    \endgroup
}

\newenvironment{pgfplotskey}[1]{%
  \begin{key}{/pgfplots/#1}%
}
{
  \end{key}
}

\newenvironment{tikzkey}[1]{%
  \begin{key}{/tikz/#1}%
}
{
  \end{key}
}

\def\choicesep{$\vert$}%
\def\choicearg#1{\texttt{#1}}

\newif\iffirstchoice
\newcommand\mchoice[1]{%
    \begingroup
    \let\margold=\marg
    \def\marg##1{{\normalfont\margold{##1}}}%
    \firstchoicetrue
    \foreach \mchoice@ in {#1} {%
        \iffirstchoice
            \global\firstchoicefalse
        \else
            \choicesep
        \fi
        \choicearg{\mchoice@}%
    }%
    \endgroup
}%




% \begin{xykey}{/path/\x label=value}
% \end{xykey}
%
% has same features with 'default', 'initially' etc as key environment
\newenvironment{xykey}[2][]{%
    \begin{pgfmanualentry}
    \def\extrakeytext{}
    \insertpathifneeded{#2}{#1}%
    \expandafter\pgfutil@in@\expandafter=\expandafter{\mykey}%
    \ifpgfutil@in@%
        \expandafter\xykey@eq\mykey\@nil
    \else
        \expandafter\xykey@noeq\mykey\@nil
    \fi
    \pgfmanualbody
}{%
    \end{pgfmanualentry}
}%

% \begin{xystylekey}{/path/\x label=value}
% \end{xystylekey}
%
% has same features with 'default', 'initially' etc as key environment
\newenvironment{xystylekey}[2][]{%
    \begin{pgfmanualentry}
    \def\extrakeytext{style, }
    \insertpathifneeded{#2}{#1}%
    \expandafter\pgfutil@in@\expandafter=\expandafter{\mykey}%
    \ifpgfutil@in@%
        \expandafter\xykey@eq\mykey\@nil
    \else
        \expandafter\xykey@noeq\mykey\@nil
    \fi
    \pgfmanualbody
}{%
    \end{pgfmanualentry}
}%

% \insertpathifneeded{a key}{/pgfplots} -> assign mykey={/pgfplots/a key}
% \insertpathifneeded{/tikz/a key}{/pgfplots} -> assign mykey={/tikz/a key}
%
% #1: the key
% #2: a default path (or empty)
\def\insertpathifneeded#1#2{%
    \def\insertpathifneeded@@{#2}%
    \ifx\insertpathifneeded@@\empty
        \def\mykey{#1}%
    \else
        \insertpathifneeded@#1\@nil
        \ifpgfutil@in@
            \def\mykey{#1}%
        \else
            \def\mykey{#2/#1}%
        \fi
    \fi
}%
\def\insertpathifneeded@#1#2\@nil{%
    \def\insertpathifneeded@@{#1}%
    \def\insertpathifneeded@@@{/}%
    \ifx\insertpathifneeded@@\insertpathifneeded@@@
        \pgfutil@in@true
    \else
        \pgfutil@in@false
    \fi
}%

% \begin{keylist}[default path]
%     {/path/option 1=value,/path/option 2=value2}
% \end{keylist}
\newenvironment{keylist}[2][]{%
    \begin{pgfmanualentry}
    \def\extrakeytext{}%
    \foreach \xx in {#2} {%
        \expandafter\insertpathifneeded\expandafter{\xx}{#1}%
        \expandafter\extractkey\mykey\@nil%
    }%
    \pgfmanualbody
}{%
  \end{pgfmanualentry}
}%

\newenvironment{pgfplotskeylist}[1]{%
    \begin{keylist}[/pgfplots]{#1}%
}{%
    \end{keylist}%
}

\newenvironment{tikzkeylist}[1]{%
    \begin{keylist}[/tikz]{#1}%
}{%
    \end{keylist}%
}

\newenvironment{anchorlist}[1]{
  \begin{pgfmanualentry}
      \foreach \xx in {#1} {%
        \pgfmanualentryheadline{锚点 {\ttfamily\declare{\xx}}}%
        \index{\xx @\protect\texttt{\xx} anchor}%
        \index{Anchors!\xx @\protect\texttt{\xx}}
        \expandafter\pgfmanualpdflabel\expandafter{\xx}{}
    }%
    \pgfmanualbody
}
{
  \end{pgfmanualentry}
}

\newenvironment{coordinatesystemlist}[1]{
  \begin{pgfmanualentry}
      \foreach \xx in {#1} {%
        \pgfmanualentryheadline{坐标系 {\ttfamily\declare{\xx}}}%
        \index{\xx @\protect\texttt{\xx} coordinate system}%
        \index{Coordinate systems!\xx @\protect\texttt{\xx}}
        \expandafter\pgfmanualpdflabel\expandafter{\xx}{}
    }%
    \pgfmanualbody
}
{
  \end{pgfmanualentry}
}
\renewenvironment{coordinatesystem}[1]{
  \begin{pgfmanualentry}
    \pgfmanualentryheadline{坐标系 {\ttfamily\declare{#1}}}%
    \index{#1@\protect\texttt{#1} coordinate system}%
    \index{Coordinate systems!#1@\protect\texttt{#1}}
    \pgfmanualpdflabel{#1}{}
    \pgfmanualbody
}
{
  \end{pgfmanualentry}
}

% \begin{xykeylist}[default path]
%     {/path/option \x1=value,/path/option \x2=value2,/path/option \x3=value}
% \end{xykeylist}
\newenvironment{xykeylist}[2][]{%
    \begin{pgfmanualentry}
    \def\extrakeytext{}
    \foreach \xx in {#2} {%
        \expandafter\insertpathifneeded\expandafter{\xx}{#1}%
        \expandafter\pgfutil@in@\expandafter=\expandafter{\mykey}%
        \ifpgfutil@in@%
            \expandafter\xykey@eq\mykey\@nil
        \else
            \expandafter\xykey@noeq\mykey\@nil
        \fi
    }%
    \pgfmanualbody
}{%
  \end{pgfmanualentry}
}%

\makeatother % FIXME this is almost surely a bug in pgfmanual-en-macros
% \begin{commandlist}
%     {\command1{arg1},\command2{\arg2}}
% \end{commandlist}
\newenvironment{commandlist}[1]{%
    \begin{pgfmanualentry}
    \foreach \xx in {#1} {%
        \expandafter\extractcommand\xx\@@%
    }%
    \pgfmanualbody
}{%
  \end{pgfmanualentry}
}%

\newenvironment{texif}[1]{%
    \begin{pgfmanualentry}
    \pgfmanualentryheadline{\declare{\texttt{\textbackslash if#1}}\meta{true code}\texttt{\textbackslash else}\meta{else code}\texttt{\textbackslash fi}}%
    \index{if#1}%
    \pgfmanualpdflabel{\\if#1}{}%
    \pgfmanualbody
}{%
  \end{pgfmanualentry}
}%
\makeatletter

\newif\ifxykeyfound

\def\pgfmanual@xykey@install@replacements{%
    \def\ { }%
    \def\space{ }%
}%

\def\xykey@eq#1=#2\@nil{%
    \begingroup
    \pgfmanual@xykey@install@replacements
    \def\x{x}%
    \xdef\mykey{#1}%
    \def\xykey@@{#1}%
    \ifx\xykey@@\mykey
        \xykeyfoundfalse
    \else
        \xykeyfoundtrue
    \fi
    \expandafter\extractkey\mykey=#2\@nil%
    \ifxykeyfound
        \def\x{y}%
        \xdef\mykey{#1}%
        \expandafter\extractkey\mykey=#2\@nil%
        \def\x{z}%
        \xdef\mykey{#1}%
        \expandafter\extractkey\mykey=#2\@nil%
    \fi
    \endgroup
}
\def\xykey@noeq#1\@nil{%
    \begingroup
    \pgfmanual@xykey@install@replacements
    \def\x{x}%
    \xdef\mykey{#1}%
    \def\xykey@@{#1}%
    \ifx\xykey@@\mykey
        \xykeyfoundfalse
    \else
        \xykeyfoundtrue
    \fi
    \expandafter\extractkey\mykey\@nil%
    \ifxykeyfound
        \def\x{y}%
        \xdef\mykey{#1}%
        \expandafter\extractkey\mykey\@nil%
        \def\x{z}%
        \xdef\mykey{#1}%
        \expandafter\extractkey\mykey\@nil%
    \fi
    \endgroup
}

% \begin{pgfplotsxykey}{\x label=value}
% \end{pgfplotsxykey}
%
% It introduces the path /pgfplots/ automatically.
%
% has same features with 'default', 'initially' etc as key environment
\newenvironment{pgfplotsxykey}[1]{%
    \begin{xykey}[/pgfplots]{#1}%
}{%
    \end{xykey}%
}


\newenvironment{pgfplotsxykeylist}[1]{%
    \begin{xykeylist}[/pgfplots]{#1}%
}{%
    \end{xykeylist}%
}


% the first, optional argument is the default key path to insert.
\newenvironment{plottype}[2][/tikz]{%
    \begin{keylist}[#1]{#2}%
    \end{keylist}
  \begin{pgfmanualentry}
    \pgfmanualentryheadline{\textcolor{gray}{{\ttfamily\char`\\addplot+[\declare{#2}]}}}%
    \pgfmanualbody
}
{
  \end{pgfmanualentry}
}

\def\index@prologue{\section*{索引}\addcontentsline{toc}{section}{索引}
}

\newenvironment{pgfplotstablecolumnkey}{%
  \begin{pgfmanualentry}
     \pgfmanualentryheadline{{\ttfamily\textcolor{gray}{/pgfplots/table/}\declare{columns/\meta{column name}}\textcolor{gray}{/.style}=\marg{key-value-list}}\hfill}%
    \pgfplotsmanualkeyindex{/pgfplots/table/columns}%
    \pgfmanualbody
}
{
  \end{pgfmanualentry}
}
\newenvironment{pgfplotstabledisplaycolumnkey}{%
  \begin{pgfmanualentry}
     \pgfmanualentryheadline{{\ttfamily\textcolor{gray}{/pgfplots/table/}\declare{display columns/\meta{index}}\textcolor{gray}{/.style}=\marg{key-value-list}}\hfill}%
    \pgfplotsmanualkeyindex{/pgfplots/table/display columns}%
    \pgfmanualbody
}
{
  \end{pgfmanualentry}
}
\newenvironment{pgfplotstablealiaskey}{%
  \begin{pgfmanualentry}
     \pgfmanualentryheadline{{\ttfamily\textcolor{gray}{/pgfplots/table/}\declare{alias/\meta{col name}}\textcolor{gray}{/.initial}=\marg{real col name}}\hfill}%
    \pgfplotsmanualkeyindex{/pgfplots/table/alias}%
    \pgfmanualbody
}
{
  \end{pgfmanualentry}
}


\def\pgfplotsmanualkeyindex#1{%
    \def\mypath{#1}%
    \def\myname{}%
    \firsttimetrue%
    \decompose#1/\nil%
}
\newenvironment{pgfplotstablecreateonusekey}{%
  \begin{pgfmanualentry}
     \pgfmanualentryheadline{{\ttfamily\textcolor{gray}{/pgfplots/table/}\declare{create on use/\meta{col name}}\textcolor{gray}{/.style}=\marg{create options}}\hfill}%
    \def\mykey{/pgfplots/table/create on use}%
    \pgfmanualbody
    \pgfplotsmanualkeyindex{/pgfplots/table/create on use}%
}
{
  \end{pgfmanualentry}
}

\def\pgfplotsassertcmdkeyexists#1{%
    \pgfkeysifdefined{/pgfplots/#1/.@cmd}\relax{%
        \pgfplots@error{DOCUMENTATION ERROR: command key /pgfplots/#1 does not exist!}%
    }%
}%

{
\catcode`\ =12%
\gdef\makespaceexpandable{\def\ { }}}%

\def\pgfplotsassertXYcmdkeyexists#1{%
    {\makespaceexpandable\def\x{x}\edef\pgfplotsassertXYcmdkeyexists@tmp{#1}%
    \pgfkeysifdefined{/pgfplots/\pgfplotsassertXYcmdkeyexists@tmp/.@cmd}\relax{%
        \pgfplots@error{DOCUMENTATION ERROR: command key /pgfplots/#1 does not exist!}%
    }}%
    {\makespaceexpandable\def\x{y}\edef\pgfplotsassertXYcmdkeyexists@tmp{#1}%
    \pgfkeysifdefined{/pgfplots/\pgfplotsassertXYcmdkeyexists@tmp/.@cmd}\relax{%
        \pgfplots@error{DOCUMENTATION ERROR: command key /pgfplots/#1 does not exist!}%
    }}%
}%

\def\pgfplotsshortstylekey #1=#2\pgfeov{%
    \pgfplotsassertcmdkeyexists{#1}%
    \pgfplotsassertcmdkeyexists{#2}%
    \begin{pgfplotskey}{#1=\marg{key-value-list}}
        是 \texttt{\pgfmanualpdfref{#2}{#2}/\pgfmanualpdfref{/handlers/.append style}{.append style}=}\marg{key-value-list} 的简写。

        它把选项追加到已存在的样式 \texttt{\pgfmanualpdfref{#2}{#2}}。
    \end{pgfplotskey}
}
\def\pgfplotsshortxystylekey #1=#2\pgfeov{%
    \pgfplotsassertXYcmdkeyexists{#1}%
    \pgfplotsassertXYcmdkeyexists{#2}%
    \begin{pgfplotsxykey}{#1=\marg{key-value-list}}
        是 {\def\x{x}\texttt{\pgfmanualpdfref{#2}{#2}/\pgfmanualpdfref{/handlers/.append style}{.append style}=}}\marg{key-value-list}
        的简写（或 $y$ 轴对应的样式
            {\def\x{y}\texttt{\pgfmanualpdfref{#2}{#2}/\pgfmanualpdfref{/handlers/.append style}{.append style}=}}\marg{key-{}value-{}list}，
        以及 $z$ 轴
            {\def\x{z}\texttt{\pgfmanualpdfref{#2}{#2}/\pgfmanualpdfref{/handlers/.append style}{.append style}=}}\marg{key-{}value-{}list}）。

        它把选项追加到已存在的样式 {\def\x{x}\texttt{\pgfmanualpdfref{#2}{#2}}}。
    \end{pgfplotsxykey}
}
\def\pgfplotsshortstylekeys #1,#2=#3\pgfeov{%
    \pgfplotsassertcmdkeyexists{#1}%
    \pgfplotsassertcmdkeyexists{#2}%
    \pgfplotsassertcmdkeyexists{#3}%
    \begin{pgfplotskeylist}{%
        #1=\marg{key-value-list},
        #2=\marg{key-value-list}}
        是 \texttt{\pgfmanualpdfref{#3}{#3}/\pgfmanualpdfref{/handlers/.append style}{.append style}=}\marg{key-value-list} 的不同简写。
    \end{pgfplotskeylist}
}
\def\pgfplotsshortxystylekeys #1,#2=#3\pgfeov{%
    \pgfplotsassertXYcmdkeyexists{#1}%
    \pgfplotsassertXYcmdkeyexists{#2}%
    \pgfplotsassertXYcmdkeyexists{#3}%
    \begin{pgfplotsxykeylist}{%
        #1=\marg{key-value-list},
        #2=\marg{key-value-list}}
        是 {\def\x{x}\texttt{\pgfmanualpdfref{#3}{#3}/\pgfmanualpdfref{/handlers/.append style}{.append style}=}}\marg{key-value-list}
        的不同简写（或 $y$ 轴对应的样式
            {\def\x{y}\texttt{\pgfmanualpdfref{#3}{#3}/\pgfmanualpdfref{/handlers/.append style}{.append style}=}\marg{key-{}value-{}list}}，以及 $z$ 轴
            {\def\x{z}\texttt{\pgfmanualpdfref{#3}{#3}/\pgfmanualpdfref{/handlers/.append style}{.append style}=}\marg{key-{}value-{}list}}%
            ）。
    \end{pgfplotsxykeylist}
}


%
% For using the correct form of including libraries in the manual.
%
\newenvironment{pgfplotslibrary}[1]{%
  \begin{pgfmanualentry}
    \pgfmanualentryheadline{{\ttfamily\char`\\usepgfplotslibrary\char`\{\declare{#1}\char`\}\space\space \char`\%\space\space  \LaTeX\space and plain \TeX}}%
    \index{#1@\protect\texttt{#1} library}%
    \index{Libraries!#1@\protect\texttt{#1}}%
    \pgfmanualentryheadline{{\ttfamily\char`\\usepgfplotslibrary[\declare{#1}]\space \char`\%\space\space Con\TeX t}}%
    \pgfmanualentryheadline{{\ttfamily\char`\\usetikzlibrary\char`\{\declare{pgfplots.#1}\char`\}\space\space \char`\%\space\space \LaTeX\space and plain \TeX}}%
    \pgfmanualentryheadline{{\ttfamily\char`\\usetikzlibrary[\declare{pgfplots.#1}]\space \char`\%\space\space Con\TeX t}}%
    \pgfmanualpdflabel{#1}{}%
    \pgfmanualbody
}
{
  \end{pgfmanualentry}
}


%
% Creates and shows a colormap with specification '#1'.
\def\pgfplotsshowcolormapexample#1{%
    \pgfplotscreatecolormap{tempcolormap}{#1}%
    \pgfplotsshowcolormap{tempcolormap}%
}

% Shows the colormap named '#1'.
\def\pgfplotsshowcolormap#1{%
    \pgfplotscolormapifdefined{#1}{\relax}{%
        \pgfplotsset{colormap/#1}%
    }%
    \pgfplotscolormaptoshadingspec{#1}{8cm}\result
    \def\tempb{\pgfdeclarehorizontalshading{tempshading}{1cm}}%
    \expandafter\tempb\expandafter{\result}%
    \pgfuseshading{tempshading}%
}

\makeatother

\def\decompose/#1/#2\nil{%
  \def\test{#2}%
  \ifx\test\empty%
    % aha.
    \index{#1@\protect\texttt{#1} key}%
    \ifx\mypath\empty
    \else
        \index{\mypath#1@\protect\texttt{#1}}%
    \fi
    \def\myname{#1}%
    %\pgfmanualpdflabel{#1}{}% No, its better to use fully qualified keys and search if necessary!
  \else%
    \iffirsttime
        \begingroup
            % also make a pdf link anchor with full key path.
            \def\hyperlabelwithoutslash##1/\nil{%
                \pgfmanualpdflabel{##1}{}%
            }%
            \hyperlabelwithoutslash/#1/#2\nil
        \endgroup
%    CF : disabled for /pgfplots/ prefix.
%        \def\mypath{#1@\protect\texttt{/#1/}!}%
%        \firsttimefalse
        \def\pgfplotslocTMPa{pgfplots}%
        \edef\pgfplotslocTMPb{#1}%
        \ifx\pgfplotslocTMPb\pgfplotslocTMPa
            \def\mypath{}%
        \else
            \def\mypath{#1@\protect\texttt{/#1/}!}%
        \fi
        \firsttimefalse
    \else
      \expandafter\def\expandafter\mypath\expandafter{\mypath#1@\protect\texttt{#1/}!}%
    \fi
    \def\firsttime{}
    \decompose/#2\nil%
  \fi%
}
\def\extracthandler#1#2\@nil{%
  \pgfmanualentryheadline{键处理器 \meta{key}{\ttfamily/\declare{#1}}#2}%
  \index{\gobble#1@\protect\texttt{#1} handler}%
  \index{Key handlers!#1@\protect\texttt{#1}}
  \pgfmanualpdflabel{/handlers/#1}{}%
}
\def\extractcommand#1#2\@@{%
  \pgfmanualentryheadline{\declare{\texttt{\string#1}}#2}%
  \removeats{#1}%
  \index{\strippedat @\protect\myprintocmmand{\strippedat}}%
  \pgfmanualpdflabel{\textbackslash\strippedat}{}%
}
\def\extractenvironement#1#2\@@{%
  \pgfmanualentryheadline{{\ttfamily\char`\\begin\char`\{\declare{#1}\char`\}}#2}%
  \pgfmanualentryheadline{{\ttfamily\ \ }\meta{environment contents}}%
  \pgfmanualentryheadline{{\ttfamily\char`\\end\char`\{\declare{#1}\char`\}}}%
  \index{#1@\protect\texttt{#1} environment}%
  \index{Environments!#1@\protect\texttt{#1}}%
  \pgfmanualpdflabel{#1}{}%
}
\renewenvironment{predefinednode}[1]{
  \begin{pgfmanualentry}
    \pgfmanualentryheadline{预定义节点 {\ttfamily\declare{#1}}}%
    \index{#1@\protect\texttt{#1} node}%
    \index{Predefined node!#1@\protect\texttt{#1}}
    \pgfmanualpdflabel{#1}{}%
    \pgfmanualbody
}
{
  \end{pgfmanualentry}
}


\makeatletter
\@ifpackageloaded{tex4ht}{
}{%
    \IfFileExists{ocg.sty}{%
        \usepackage{ocg}%
    }{%
        \usepackage{pgfplots_ocg_copy}%
    }
}%
\makeatother

\usepackage{nicefrac}

\graphicspath{{figures/}}

\edef\preambleconfig{width=7cm,compat=\pgfkeysvalueof{/pgfplots/compat/mostrecent}}

% generate the following file such that it does not need to be updated
% manually for each release:
\makeatletter
\immediate\openout\w@pgf@writea pgfplots.preliminaries.compatcurrent.tmp
\immediate\write\w@pgf@writea{%
\string\begin{codeexample}[code only]^^J%
\string\usepackage{pgfplots}^^J%
\string\pgfplotsset{compat=\pgfkeysvalueof{/pgfplots/compat/mostrecent}}^^J%
\string\end{codeexample}%
}%
\immediate\closeout\w@pgf@writea


\expandafter\pgfplotsset\expandafter{\preambleconfig}

\pgfplotsset{lua debug=off and silent}

\def\pgfplotsmanualcurlibrary{}

% And now, invoke
%     /codeexample/typeset listing/.add={% Preamble:\pgfplotsset{\preambleconfig}}{}}
% since listings are VERBATIM, I need to do some low-level things
% here to get the correct \catcodes:
\pgfkeys{/codeexample/typeset listing/.add code={%
        \ifcode@execute
            \begingroup
            \def\b@show@header{0}%
            \pgfutil@in@{axis}{#1}%
            \ifpgfutil@in@
                \def\b@show@header{1}%
            \fi
            %
            \pgfutil@in@{addplot}{#1}%
            \ifpgfutil@in@
                \def\b@show@header{1}%
            \fi
            %
            \pgfutil@in@{smithchart}{#1}%
            \ifpgfutil@in@
                \def\b@show@header{1}%
            \fi
            \if1\b@show@header
                {\tiny
                    \% Preamble: \pgfmanualpdfref{\textbackslash pgfplotsset}{\pgfmanual@pretty@backslash pgfplotsset}%
                        \pgfmanual@pretty@lbrace \expandafter\pgfmanualprettyprintpgfkeys\expandafter{\preambleconfig}\pgfmanual@pretty@rbrace
                        \ifx\pgfplotsmanualcurlibrary\pgfutil@empty
                        \else
                            \pgfmanual@pretty@backslash usepgfplotslibrary\pgfmanual@pretty@lbrace \pgfplotsmanualcurlibrary\pgfmanual@pretty@rbrace
                        \fi
                }%
            \fi
            \endgroup
        \fi
    }{},%
    %/codeexample/typeset listing/.show code,
}%
\makeatother

\pgfplotsset{
    %every axis/.append style={width=7cm},
    filter discard warning=false,
}

\pgfqkeys{/codeexample}{%
    every codeexample/.append style={
        width=8cm,
        /pgfplots/legend style={fill=graphicbackground},
        /pgfplots/contour/every contour label/.append style={
            every node/.append style={fill=graphicbackground}
        },
    },
    tabsize=4,
}

\usetikzlibrary{backgrounds,patterns}
% Global styles:
\tikzset{
  shape example/.style={
    color=black!30,
    draw,
    fill=yellow!30,
    line width=.5cm,
    inner xsep=2.5cm,
    inner ysep=0.5cm}
}

\newcommand{\FIXME}[1]{\textcolor{red}{(FIXME: #1)}}

% fuer environment 'sidewaysfigure' bspw
% \usepackage{rotating}

\newcommand\Tikz{Ti\textit kZ}
\newcommand\PGF{\textsc{pgf}}
\newcommand\PGFPlots{\pgfplotsmakefilelinkifuseful{pgfplots}{\textsc{pgfplots}}}
\newcommand\PGFPlotstable{\pgfplotsmakefilelinkifuseful{pgfplotstable}{\textsc{PgfplotsTable}}}

\makeindex

% Fix overfull hboxes automatically:
\tolerance=2000
\emergencystretch=10em

\tikzset{prefix=gnuplot/pgfplots_} % prefix for 'plot function'

\author{%
    Dr.\ Christian Feuersänger\\
    {\footnotesize\texttt{cfeuersaenger@users.sourceforge.net}}}%


\usepackage{todonotes}
\newcommand\todosp[2][]{%        % Stefan Pinnow
    \todo[
        disable,
        color=blue!40,
        #1]{#2}
}

% ---------------------------------------------------------------
% 修复：doc 宏包 v2 的 \ltxdoc@meta 在 TeX Live 2026 中未定义，
% 导致 \marg/\oarg/\parg（内部调用 \meta）写 aux/idx 时报
% "Undefined control sequence"。
% 对策（双保险）：
%  1) 兜底提供 \ltxdoc@meta（doc 的 \meta 依赖它）；
%  2) 在 \AtBeginDocument 末尾覆盖 \meta 为鲁棒命令
%     （doc v2 会在 begin document 时重定义 \meta，必须最后覆盖）。
\makeatletter
\providecommand\ltxdoc@meta[1]{%
    {\normalfont\itshape\textlangle #1\textrangle}%
}
\AtBeginDocument{%
    \DeclareRobustCommand\meta[1]{%
        {\normalfont\itshape\textlangle #1\textrangle}%
    }%
}
\makeatother


\usepackage{array}
\usepackage{colortbl}
\usepackage{booktabs}
\usepackage{eurosym}
\usepackage{multirow}

\pgfqkeys{/codeexample}{%
    every codeexample/.style={
        width=4cm,
        /pgfplots/every axis/.append style={legend style={fill=graphicbackground}}
    },
    narrow/.style={width=7cm},
    tabsize=4,
}

\pgfkeys{/pgfmanual/gray key prefixes/.add={/pgfplots/table/,/pgf/number format/,}{}}

\makeatletter
\pgfkeys{%
    /codeexample/prettyprint/key name with handler/.code 2 args={%
        \gdef\pgfplotstablemanualautocheck{}%
        \foreach \special in {columns,create on use,alias,display columns}{%
            \expandafter\pgfutil@in@\expandafter{\special/}{#1}%
            \ifpgfutil@in@
                \xdef\pgfplotstablemanualautocheck{\special}%
                \breakforeach
            \fi
        }%
        \ifx\pgfplotstablemanualautocheck\pgfutil@empty%
            \pgfmanualpdfref{#1}{#1}/\pgfmanualpdfref{/handlers/#2}{#2}%
        \else
            \expandafter\def\expandafter\temp\expandafter##\expandafter1\pgfplotstablemanualautocheck/##2\relax{%
                \pgfmanualpdfref
                    {/pgfplots/table/\pgfplotstablemanualautocheck}%
                    {##1\pgfplotstablemanualautocheck/}%
                ##2/%
                \pgfmanualpdfref{/handlers/#2}{#2}%
            }%
            \temp#1\relax
        \fi
    },
    /pdflinks/search key prefixes in/.add={/pgf/number format/,}{,/pgfplots/table/create col/},
    /pdflinks/show labels=false,
}
\makeatother

\pgfplotstableset{
    %debug=true,
    begin table=\begin{tabular}[b],
}

\title{%
    \PGFPlotstable{} 宏包手册\\
    {\small\PGFPlots{} 的组成部分，版本 \pgfplotsversion}\\
    {\small\url{https://github.com/pgf-tikz/pgfplots}}}

\begin{document}
\maketitle
\begin{abstract}
    本宏包从输入读取以制表符分隔的数值表格，并生成排版精美的
    \LaTeX{} tabular 代码。它会按所需精度四舍五入，并以不同的数字格式样式打印。
\end{abstract}

\tableofcontents


\section{引言}

\PGFPlotstable{} 是 \PGFPlots{} 的一个轻量子宏包，它借用 \PGFPlots{} 的表格输入
方法和数字格式化技术，把以制表符分隔的表格转换为 tabular。

它的输入是一个文本文件，包含以空白分隔的行，可能以列名开头。它的输出是一个
\LaTeX{} tabular\footnote{关于纯 \TeX{} 和 Con\TeX{}t，请参见
节~\ref{sec:pgfplotstable:context} 中的说明。}，其中包含文本表格中被选中的列，
按所需精度四舍五入，并以所需的数字格式（定点、整数、科学计数法等）打印。
输出是 \LaTeX{} 代码，该代码最终由 \LaTeX{} 排版。

换句话说，\PGFPlotstable{} 不过是一个多少有点聪明的代码生成器，吐出类似
|\begin{tabular}...\end{tabular}| 的东西。如果你想定制依赖于行或列的样式，
或者你有数值数据、希望自动格式化其内容，就可以使用它。

它通过下面的方式使用
%
% NO GALLERY
\begin{codeexample}[code only]
\usepackage{pgfplotstable}
% recommended:
%\usepackage{booktabs}
%\usepackage{array}
%\usepackage{colortbl}
\end{codeexample}
%
\noindent 并且要求已安装 \PGFPlots{} 和 \PGF{} $ \ge 2.00$。

\begin{command}{\pgfplotstableset\marg{key-value-options}}
    本宏包的用户界面基于键值选项。它们决定显示什么、如何排版以及计算什么。

    键值对可以通过两种方式设置：
    %
    \begin{enumerate}
        \item 作为整个文档（或文档某一部分）的默认设置，使用
            |\pgfplotstableset|\marg{options}。例如，文档的导言区可以包含
            %
% NO GALLERY
\begin{codeexample}[code only]
\pgfplotstableset{fixed zerofill,precision=3}
\end{codeexample}
            %
            以配置小数点后保留 $3$ 位精度，并补零使所有定点数都恰好有 $3$ 位。
        \item 作为只影响单个表格的选项。它作为可选参数提供给相应的表格排版命令，
            例如 |\pgfplotstabletypeset|\oarg{options}\marg{file}。
    \end{enumerate}
    %
    下面的例子会展示这两种方式。

    了解 |pgfkeys| 有助于更深入地理解本宏包，因为 |/.style|、|/.append style|
    等都是 |pgfkeys| 特有的。要进一步了解 |pgfkeys|，请参阅
    \PGF{} 手册~\cite[Section ``pgfkeys'']{tikz} 或更简短的介绍~\cite{keyvalintro}。
    否则，直接跳到本文档提供的例子即可。

    你会看到键前缀 |/pgfplots/table/| 和 |/pgf/number format/|。
    在 \PGFPlotstable{} 中使用时，这些前缀可以省略；它们属于 |pgfkeys| 的
    ``默认键路径''。
\end{command}


\section{加载并显示数据}

\subsection{文本表格输入格式}

\PGFPlotstable{} 处理纯文本文件表格，其中的条目（``单元格''）由一个分隔字符隔开。
初始的分隔字符是 ``空白''，意思是 ``至少一个空格或制表符''（参见下面的选项
|col sep|）。这类表格可以有一行表头，其中包含列名，而其余大多数列通常包含数值数据。

\noindent 下面的列表展示了 |pgfplotstable.example1.dat|，本文档中会经常用到它。
%
\begin{codeexample}[code only]
# Convergence results
# fictional source, generated 2008
level    dof      error1            error2   info     grad(log(dof),log(error2)) quot(error1)
1        4        2.50000000e-01    7.57858283e-01    48       0                 0
2        16       6.25000000e-02    5.00000000e-01    25       -3.00000000e-01   4
3        64       1.56250000e-02    2.87174589e-01    41       -3.99999999e-01   4
4        256      3.90625000e-03    1.43587294e-01    8        -5.00000003e-01   4
5        1024     9.76562500e-04    4.41941738e-02    22       -8.49999999e-01   4
6        4096     2.44140625e-04    1.69802322e-02    46       -6.90000001e-01   4
7        16384    6.10351562e-05    8.20091159e-03    40       -5.24999999e-01   4
8        65536    1.52587891e-05    3.90625000e-03    48       -5.35000000e-01   3.99999999e+00
9        262144   3.81469727e-06    1.95312500e-03    33       -5.00000000e-01   4.00000001e+00
10       1048576  9.53674316e-07    9.76562500e-04    2        -5.00000000e-01   4.00000001e+00
\end{codeexample}
%
以 `|%|' 或 `|#|' 开头的行被视为注释行，会被忽略。

\begin{command}{\pgfplotstabletypeset\oarg{optional arguments}\marg{file name {\normalfont or} \textbackslash macro {\normalfont or} inline table}}
    加载（或获取）一个表格，并使用当前的数字格式和表格选项配置对其排版。

    如果第一个参数是文件名，则从磁盘加载表格。如果它是一个已加载的表格
    （参见 |\pgfplotstableread| 或 |\pgfplotstablenew|），就直接使用它。
    否则，如果它是内联表格数据，这些数据会像在文件中找到的那样被解析
    （参见 |\pgfplotstableread|）。
    %
\begin{codeexample}[]
\pgfplotstabletypeset[sci zerofill]{
    a b
    5000 1.234e5
    6000 1.631e5
    7000 2.1013e5
    9000 1000000
}
\end{codeexample}

\begin{codeexample}[]
\pgfplotstabletypeset{pgfplotstable.example1.dat}
\end{codeexample}

    \noindent 可以用 \meta{optional arguments} 定制配置。配置可以针对整个表格，
    也可以针对特定的列（或行）。

\begin{codeexample}[]
\pgfplotstableset{% global config, for example in the preamble
    % these columns/<colname>/.style={<options>} things define a style
    % which applies to <colname> only.
    columns/dof/.style={int detect,column type=r,column name=\textsc{Dof}},
    columns/error1/.style={
        sci,sci zerofill,sci sep align,precision=1,sci superscript,
        column name=$e_1$,
    },
    columns/error2/.style={
        sci,sci zerofill,sci sep align,precision=2,sci 10e,
        column name=$e_2$,
    },
    columns/{grad(log(dof),log(error2))}/.style={
        string replace={0}{}, % erase '0'
        column name={$\nabla e_2$},
        dec sep align,
    },
    columns/{quot(error1)}/.style={
        string replace={0}{}, % erase '0'
        column name={$\frac{e_1^{(n)}}{e_1^{(n-1)}}$}
    },
    empty cells with={--}, % replace empty cells with '--'
    every head row/.style={before row=\toprule,after row=\midrule},
    every last row/.style={after row=\bottomrule}
}
\pgfplotstabletypeset[ % local config, applies only for this table
    1000 sep={\,},
    columns/info/.style={
        fixed,fixed zerofill,precision=1,showpos,
        column type=r,
    }
]
{pgfplotstable.example1.dat}
\end{codeexample}

    \noindent 所有这些选项都会在后面各节中详细解释。

    你也可以使用类似于 tabular 环境的输入格式：
    %
\begin{codeexample}[width=8cm]
\pgfplotstabletypeset
    [col sep=&,row sep=\\,sci zerofill]
{
    level &  dof &    error \\
    1 &      4 &      2.50000000e-01 \\
    2 &      16 &     6.25000000e-02 \\
    3 &      64 &     1.56250000e-02 \\
    4 &      256 &    3.90625000e-03 \\
    5 &      1024 &   9.76562500e-04 \\
    6 &      4096 &   2.44140625e-04 \\
    7 &      16384 &  6.10351562e-05 \\
    8 &      65536 &  1.52587891e-05 \\
    9 &      262144 & 3.81469727e-06 \\
    10 &     1048576 &9.53674316e-07 \\
}
\end{codeexample}

    技术说明：每个被打开的文件都会记录到你的日志文件中。
\end{command}

\begin{command}{\pgfplotstabletypesetfile\oarg{optional arguments}\marg{file name}}
    加载表格 \meta{file name} 并对其排版。从 \PGFPlotstable{} 1.2 起，
    这个命令是 |\pgfplotstabletypeset| 的别名，也就是说第一个参数既可以是文件名，
    也可以是一个已加载的表格。
\end{command}

\begin{commandlist}{%
    \pgfplotstableread\marg{file name}\marg{\textbackslash macro},
    \pgfplotstableread\marg{inline table}\marg{\textbackslash macro}%
}
    把一个表格加载到 \TeX{} 宏 \meta{\textbackslash macro} 中。该宏会把表格存为内部
    结构，可多次使用。
    %
\begin{codeexample}[]
\pgfplotstableread{pgfplotstable.example1.dat}\loadedtable
\pgfplotstabletypeset[columns={dof,error1}]\loadedtable
\hspace{2cm}
\pgfplotstabletypeset[columns={dof,error2}]\loadedtable
\end{codeexample}
    %
    第一个参数可以像本例这样是一个 \meta{file name}。也可以直接提供表格数据：
    %
\begin{codeexample}[code only]
% Alternative: inline table data:
\pgfplotstableread{
level   dof     error1  error2  info    grad(log(dof),log(error2))      quot(error1)
1       4       2.50000000e-01  7.57858283e-01  48      0               0
2       16      6.25000000e-02  5.00000000e-01  25      -3.00000000e-01 4
3       64      1.56250000e-02  2.87174589e-01  41      -3.99999999e-01 4
4       256     3.90625000e-03  1.43587294e-01  8       -5.00000003e-01 4
5       1024    9.76562500e-04  4.41941738e-02  22      -8.49999999e-01 4
6       4096    2.44140625e-04  1.69802322e-02  46      -6.90000001e-01 4
7       16384   6.10351562e-05  8.20091159e-03  40      -5.24999999e-01 4
8       65536   1.52587891e-05  3.90625000e-03  48      -5.35000000e-01 3.99999999e+00
9       262144  3.81469727e-06  1.95312500e-03  33      -5.00000000e-01 4.00000001e+00
10      1048576 9.53674316e-07  9.76562500e-04  2       -5.00000000e-01 4.00000001e+00
}\loadedtable
% can be used as above:
\pgfplotstabletypeset[columns={dof,error1}]\loadedtable
\hspace{2cm}
\pgfplotstabletypeset[columns={dof,error2}]\loadedtable
\end{codeexample}
    %
    它会自动检查第一个参数包含的是内联数据还是文件名。

    判断第一个参数是内联数据还是文件名的检查如下进行：如果 |format=auto|，
    则第一个参数被视为文件名，除非它包含 |row sep| 字符（参见 |row sep|）。
    如果 |format=inline|，它总是被视为内联数据。如果 |format=file|，它就是文件名。


    \paragraph{特殊情况与更多细节：}

    \begin{itemize}
        \item 内联数据格式是 ``脆弱的''。如果你遇到问题，请用 `|\\|'
            配合 |row sep=\\| 来结束表格（|row sep| 的文档中包含其他方式和更多解释）。
        \item 这个命令还有一些变体，它们并不真正建立一个结构体，而是把每一行
            报告给一个 ``监听器''。还有一个结构体可以避免被 \TeX{} 作用域保护。
            如果你需要这类东西，请阅读源代码注释。
        \item 技术说明：每个被打开的文件都会记录到你的日志文件中。
        \item 注意：避免把 `|\table|' 用作名字，它会与 \LaTeX{} 的
            |\begin{table}| 冲突。
    \end{itemize}
\end{commandlist}

\begin{key}{/pgfplots/table/col sep=\mchoice{space,tab,comma,semicolon,colon,braces,\&,ampersand} (initially space)}
    指定读取表格时所用的列分隔字符。初始选项 |space| 表示 ``至少一个空白''。
    空白指制表位或空格（换行符总是用于分隔行）。

    例如，文件 |pgfplotstable.example1.csv| 使用逗号作为分隔字符。
    %
\begin{codeexample}[code only]
# Convergence results
# fictional source generated 2008
level,dof,error1,error2,info,{grad(log(dof),log(error2))},quot(error1)
1,9,2.50000000e-01,7.57858283e-01,48,0,0
2,25,6.25000000e-02,5.00000000e-01,25,-1.35691545e+00,4
3,81,1.56250000e-02,2.87174589e-01,41,-1.17924958e+00,4
4,289,3.90625000e-03,1.43587294e-01,8,-1.08987331e+00,4
5,1089,9.76562500e-04,4.41941738e-02,22,-1.04500712e+00,4
6,4225,2.44140625e-04,1.69802322e-02,46,-1.02252239e+00,4
7,16641,6.10351562e-05,8.20091159e-03,40,-1.01126607e+00,4
8,66049,1.52587891e-05,3.90625000e-03,48,-1.00563427e+00,3.99999999e+00
9,263169,3.81469727e-06,1.95312500e-03,33,-1.00281745e+00,4.00000001e+00
10,1050625,9.53674316e-07,9.76562500e-04,2,-1.00140880e+00,4.00000001e+00
\end{codeexample}
    %
    因此，我们读取它时需要指定 |col sep=comma|。
    %
\begin{codeexample}[]
\pgfplotstabletypeset[col sep=comma]{pgfplotstable.example1.csv}
\end{codeexample}
    %
    如果你的所有表格都用逗号作为列分隔符，你可以在导言区调用一次
    |\pgfplotstableset{col sep=comma}|。

    请注意，如果单元格条目（例如列名）包含分隔字符，你需要用\emph{花括号}
    把该列条目括起来：|{grad(log(dof),log(error2)}|。如果你想使用不配对的花括号，
    需要在花括号前写一个反斜杠。例如，名字 `|column{withbrace|' 需要写成
    `|column\{withbrace|'。

    当 |col sep|$\neq$|space| 时，空格会被视为参数的一部分（不做裁剪）。
    不过（和 \TeX{} 中通常一样），连续多个空格和制表符会被替换为一个空白。
    当然，如果 |col sep=tab|，制表符就是列分隔符，会被特殊处理。

    此外，如果在 |col sep=space| 的情况下需要空单元格，你必须用 |{}| 来界定
    这样一个单元格，因为 |col sep=space| 使用\emph{至少}一个空白（会吞掉其后所有的空白）。

    取值 |col sep=braces| 比较特殊，因为它实际上使用两个分隔字符。
    对于这个选项，每个单元格条目都由一对开花括号和闭花括号界定，即 \meta{entry}。
    此外，在 |braces| 情形下，单元格条目之间的任何空白（空格和制表符）都会被\emph{跳过}，
    直到找到下一个 \meta{entry}。

    |col sep=braces| 的另一个特别之处是它支持\emph{多行}单元格：配对花括号内的一切
    都被视为某个单元格的一部分，这包括换行：\footnote{配对花括号内换行的这种处理，
    实际上同样适用于其他所有列分隔符（这是 \TeX{} 的读行特性）。换句话说：只要你把
    单元格用配对花括号括起来，那么对每一种列分隔符你\emph{都可以}使用多行单元格。
    不过，\texttt{col sep=braces} 有一个特殊处理：行尾算作空白字符；而对于其他所有
    \texttt{col sep} 取值，这种空白都会被抑制，以消除多余空格。}

    |col sep=&| 的情形（通常配合 |row sep=\\|）允许像你平时在 \LaTeX{} 中那样
    输入表格来读取。这会自动启用 |trim cells|。
\end{key}

\begin{key}{/pgfplots/table/trim cells=\mchoice{true,false} (initially false)}
    如果启用，读取表格时会去除前导和尾随空白。

    当 |col sep|$\neq$|space| 而你的单元格中含有空格时，这可能是必要的。
    对于 |col sep=&|，它会自动启用。
\end{key}

\begin{key}{/pgfplots/table/header=\mchoice{true,false,has colnames} (initially true)}
    配置在输入过程中是否应自动识别列名。

    第一个非注释行\emph{可以}是包含列名的表头。|header| 键配置如何检测是否真的如此。

    选项 \declaretext{true} 启用列名的自动检测：如果第一个非注释行至少包含一个
    非数值条目（例如 `|a name|'），则该行中的每个条目都被认为是列名。如果第一个
    非注释行只包含数值数据，它就被用作数据行。在这种情况下，列索引将被指定为列 ``名''。

    选项 \declaretext{false} 与上一种情形相同，也就是说即使第一行包含字符串，
    它们也不会被识别为列名。

    最后，选项 \declaretext{has colnames} 与 |false| 相反：它假定第一个非注释行
    \emph{包含}列名。
    换句话说：即使第一行只含有数字，它们也会被视为列\emph{名}。

    注意，这个键只配置\emph{输入}表格的表头。要配置\emph{输出}表头，
    可参看 |every head row|。
\end{key}

\begin{pgfplotskey}{table/format=\mchoice{auto,inline,file} (initially auto)}
    配置 |\pgfplotstableread|\marg{input} 的第一个参数所期待的格式。

    选项 \declaretext{inline} 直接把表格数据作为参数，其中行由 |row sep| 分隔。
    内联数据是 ``脆弱的''，因为 \TeX{} 可能吞掉行尾字符（或 |col sep| 字符）。
    细节见 |row sep|。

    选项 \declaretext{file} 期待一个文件名。

    选项 \declaretext{auto} 在提供给 |\pgfplotstableread| 的第一个参数中查找 |row sep|。
    如果找到 |row sep|，就是内联数据，否则就是文件名。
\end{pgfplotskey}

\begin{pgfplotskey}{table/row sep=\mchoice{newline,\string\\} (initially newline)}
    配置用于分隔内联表格数据格式（见 |format=inline|）中各行的字符。

    选项 \declaretext{newline} 使用表格数据中（即输入文件或任何内联表格数据）
    所出现的行尾。

    选项 \declaretext{\string\\} 用 `|\\|' 来表示一行的结束。

    注意，对内联表格数据使用 \declaretext{newline} 是 ``脆弱的''：你不能在
    \TeX{} 宏内部提供这样的数据（这不适用于输入文件）。每当遇到问题时，
    请按如下方式处理：
    %
    \begin{enumerate}
        \item 第一种办法：在任何宏声明\emph{之外}调用
            |\pgfplotstableread|\marg{data}|\yourmacro|。
        \item 使用 |row sep=\\|。
    \end{enumerate}
    %
    如果你对内联数据和特殊的 |col sep| 选项（如 |col sep=tab|）遇到问题，同样适用。

    这类问题的原因在于：\TeX{} 会扫描宏体并把换行替换为空白。它还会做其他这类替换，
    而这些替换无法撤销（甚至可能无法被发现）。
\end{pgfplotskey}

\begin{key}{/pgfplots/table/ignore chars=\marg{comma-separated-list} (initially empty)}
    允许为单个字符定义一个 ``忽略列表''。在输入文件中找到、同时也出现在
    \meta{comma-separated-list} 中的任何字符，都会被悄悄丢弃。其处理完全等同于
    你一开始根本就没有写出它们。

    例如，假设我们有 |pgfplotstable.example5.dat|，其内容为
    %
        \lstinputlisting[basicstyle=\ttfamily\footnotesize]{pgfplotstable.example5.dat}

    \noindent 那么，我们可以通过下面的写法忽略其中若干字符
    %
\begin{codeexample}[]
\pgfplotstabletypeset
    [col sep=comma,ignore chars={(,),\ ,\#}]
    {pgfplotstable.example5.dat}
\end{codeexample}

    \meta{comma-separated-list} 的每个列表元素应恰好包含一个字符，这些单个字符之间
    用逗号分隔。某些特殊字符，如逗号、空白、井号、百分号或反斜杠，需要在前面加一个
    反斜杠来转义。

    除了普通字符之外，它还支持从输入文件中消除任何二进制代码。例如，假设你的文件中
    有代码为 |0x01|（十六进制记法）的二进制字符。那么，使用
    %
\begin{codeexample}[code only]
\pgfplotstableset{ignore chars={\^^01}}
\end{codeexample}
    \noindent 来悄悄消除它们。|^^|\meta{digit}\meta{digit} 记法是 \TeX{} 的一个特性，
    用于以十六进制编码提供字符，其中 \meta{digit} 是 |0123456789abcdef| 之一。
    我不知道 |\^^01| 中的反斜杠是否总是必要，你可以试试看。还有一种基于字符的记法，
    其中 |\^^M| 是 \meta{newline}，|\^^I| 是 \meta{tab}。更多细节参见~\cite{texbook}。

    注意，剥离所有这些字符之后，输入表格必须仍然有效——它仍应包含列分隔符且各列平衡。

    这个设置同样适用于 \PGFPlots{} 的 |\addplot table| 和 |\addplot file|。

    注意，当 |ignore chars| 应用于 |format=inline| 或 |format=auto| 时是 ``脆弱的''。
    如果遇到问题，请考虑使用 |format=file|。\footnote{关于处理脆弱的内联表格格式的
    更多信息，另见 |row\ sep|。}
\end{key}

\begin{key}{/pgfplots/table/white space chars=\marg{comma-separated-list} (initially empty)}
    允许定义一个单个字符的列表，这些字符实际上被当作空白处理（除制表符和空格之外）。
    它在想要使用多个列分隔字符时可能有用。

    |white space chars| 列表的用法与 |ignore chars| 完全相同，上面的说明同样适用。
\end{key}

\begin{key}{/pgfplots/table/text special chars=\marg{comma-separated-list} (initially empty)}
    允许定义一个单个字符的列表，这些字符实际上被当作普通文本处理
    （就像任何类别码为~$12$ 的字符一样）。

    |text special chars| 列表的用法与 |ignore chars| 完全相同，上面的说明同样适用。

    注意，某些（特定的）特殊字符是由 |verb string type| 安装的。
\begin{codeexample}[]
\pgfplotstabletypeset[
    font=\ttfamily,
    verb string type,
]{
URI
http://example.org/a#b
}
\end{codeexample}
\end{key}

\begin{key}{/pgfplots/table/comment chars=\marg{comma-separated-list} (initially empty)}
    允许添加一个或多个\emph{额外的}注释字符。这些字符中的每一个都具有与 |#| 字符
    类似的效果，也就是说该输入行中其后的所有字符都会被跳过。
    %
\begin{codeexample}[]
\pgfplotstabletypeset[comment chars=!]{
! Some comments
1 0
2 -10
! another comment line
3 0
}
\end{codeexample}
    %
    上面的例子使用 `|!|' 作为额外的注释字符（这样就能解析 Touchstone 文件）。
\end{key}

\begin{key}{/pgfplots/table/percent is letter=\mchoice{true,false} (initially false)}
    允许控制如何处理输入文件中的百分号字符。\PGFPlotstable{} 过去会保留 \TeX{} 中
    百分号作为注释字符的特殊含义，也就是说百分号之后的所有字符都会被忽略。

    取值 \declaretext{true} 表示百分号只是一个普通字母，没有特殊含义。
    取值 \declaretext{false} 表示百分号仍像过去那样是特殊的。
\end{key}

\begin{key}{/pgfplots/table/text indicator=\marg{char} (initially empty)}
    允许定义一个文本指示字符。如果定义了它，就可以用它在单元格正文中括起分隔符：
    %
\begin{codeexample}[]
\pgfplotstabletypeset[text indicator=",string type]{
A B
"A long cell" normalcell
}
\end{codeexample}

    如果 |text indicator| 也出现在单元格内部，就必须把它双写：
    %
\begin{codeexample}[]
\pgfplotstabletypeset[text indicator=",string type]{
A B
"A long ""cell""" normalcell
}
\end{codeexample}

    不以文本指示符开头的单元格按通常方式处理。注意，|text indicator| 不支持多行单元格。
\end{key}

\begin{key}{/pgfplots/table/skip first n=\marg{integer} (initially 0)}
    允许跳过输入文件的前 \meta{integer} 行。这些行不会被处理。
    %
\begin{codeexample}[]
\pgfplotstabletypeset[skip first n=4]{%<- this '%' is important. Otherwise, the
                                      %newline here would delimit an (empty) row.
    XYZ Format,
    Version 1.234
    Date 2010-09-01
    @author Mustermann
    A B C
    1 2 3
    4 5 6
}
\end{codeexample}
\end{key}

% FIXME: duplicate code in pgfplotstable.tex and pgfplots.reference.axis-addplot.tex
\begin{key}{/pgfplots/table/search path=\marg{comma-separated-list} (initially .)}
    允许为输入表格提供搜索路径。每当 \PGFPlots{} 试图读取一个数据文件时，都会计算
    这个变量。这包括 |\pgfplotstableread| 和 |\addplot table|；它的值是一个逗号分隔的
    路径名列表。所请求的文件将从该列表中第一个匹配的位置读取。

    使用 `|.|' 表示用 \TeX{} 正常的文件搜索过程来搜索（即通常在当前工作目录中）。

    \meta{comma-separated-list} 中的条目可以相对于当前工作目录，也就是说像
    |search path={.,../datafiles}| 这样是可以接受的。
\end{key}

% FIXME: duplicate code in pgfplotstable.tex and pgfplots.reference.axis-addplot.tex
\begin{key}{/pgfplots/table/search path/implicit .=\mchoice{true,false} (initially true)}
    \PGFPlotstable{} 允许隐式地把 `|.|' 添加到 |search path| 的值中，因为在搜索路径的
    许多应用中通常都这样假定。

    初始配置 |search path/implicit .=true| 将确保：如果用户的值不包含 `|.|'，
    就在 |search path| 前面加上 `|.|'。

    取值 |search path/implicit .=false| 则不添加 `|.|'。

    请记住，`|.|' 表示 ``让 \TeX{} 自己去搜索文件''。这通常会在当前工作目录中找到文件，
    但也会包括对 |TEXINPUTS| 的处理。
\end{key}


\subsection{选择列及其外观样式}

\begin{key}{/pgfplots/table/columns=\marg{comma-separated-list}}
    选择表格中特定的列。如果此选项为空（未提供），则选择所有可用的列。

    在 \meta{comma-separated-list} 内部，应填写出现在表格表头中的列名。如果没有表头，
    就直接用列索引。如果有列名，则可以用特殊语法 |[index]|\meta{integer} 按索引选择列。
    第一列的索引为~$0$。
    %
\begin{codeexample}[]
\pgfplotstabletypeset[columns={dof,level,[index]4}]{pgfplotstable.example1.dat}
\end{codeexample}

    |pgfkeys| 的特殊功能 |\pgfplotstableset{columns/.add={}{,a further col}}|
    允许把一个值（这里是 `|,a further col|'）\emph{追加}到当前值之后。
    细节参见 |/.add|。
\end{key}

\begin{pgfplotstablealiaskey}
    为 \meta{real col name} 所表示的列指定新名字 \meta{col name}。此后，
    访问 \meta{col name} 将使用与列 \meta{real col name} 相关联的数据。
    %
\begin{codeexample}[]
% in preamble:
\pgfplotstableset{
    alias/newname/.initial=b,
}

% in document:
\pgfplotstabletypeset[
    columns={a,newname},% access to `newname' is the same as to `b'
]{
    a b
    1 2
    3 4
    5 6
}%
\end{codeexample}
    %
    你可以用 |columns/|\meta{col name}|/.style| 来为这个别名（而不是原始列名）指定样式。

    如果同时存在一个别名和一个同名的列，则优先使用列名。此外，如果还存在一个同名的
    |create on use| 声明，也优先使用它。

    如果 \meta{col name} 包含键设置所需的字符，你需要用花括号把它括起来：
    ``|alias/{name=wi/th,special}/.initial={othername}|''。

    每当查询列时都会用到这个键，它也适用于 \PGFPlots{} 的 |\addplot table| 语句。
\end{pgfplotstablealiaskey}

\begin{pgfplotstablecolumnkey}
    只为 \meta{column name} 设置 \meta{key-value-list} 中的所有选项。
    %
\begin{codeexample}[]
\pgfplotstabletypeset[
    columns/error1/.style={
        column name=$L_2$,
        sci,sci zerofill,sci subscript,
        precision=3},
    columns/error2/.style={
        column name=$A$,
        sci,sci zerofill,sci subscript,
        precision=2},
    columns/dof/.style={
        int detect,
        column name=\textsc{Dof}
    }
]
    {pgfplotstable.example1.dat}
\end{codeexample}
    %
    如果你的列名包含逗号 `|,|'、斜杠 `|/|' 或等号 `|=|'，你需要用花括号把列名括起来。
    %
\begin{codeexample}[narrow]
\pgfplotstabletypeset[
    columns={dof,error1,{grad(log(dof),log(error2))}},
    columns/error1/.style={
        column name=$L_2$,
        sci,sci zerofill,sci subscript,
        precision=3},
    columns/dof/.style={
        int detect,
        column name=\textsc{Dof}},
    columns/{grad(log(dof),log(error2))}/.style={
        column name=slopes $L_2$,
        fixed,fixed zerofill,
        precision=1}
]
    {pgfplotstable.example1.dat}
\end{codeexample}
    %
    如果你的表格没有列名，可以直接用整数索引代替 \meta{column name} 来引用列。
    如果你有列名，就不能用索引来设置列样式。
\end{pgfplotstablecolumnkey}

\begin{pgfplotstabledisplaycolumnkey}
    把 \meta{key-value-list} 中的所有选项只应用于输出表中出现在位置 \meta{index} 的那一列。

    与 |table/columns/|\meta{name} 样式不同，此选项引用的是输出表而不是输入表。
    由于输出表没有唯一的列名，你只能通过索引访问列。

    索引从~$\mathit{index}=0$ 开始。

    显示列样式覆盖输入列样式。
\end{pgfplotstabledisplaycolumnkey}

\begin{stylekey}{/pgfplots/table/every col no \meta{index}}
    一个与 |display columns/|\meta{index} 相同的样式：它只应用于输出表中位置
    \meta{index} 处的列。

    细节见 |display columns/|\meta{index}。
\end{stylekey}

\begin{key}{/pgfplots/table/column type=\marg{tabular column type} (initially c)}
    包含 |tabular| 的列类型。

    如果所有列类型都为空，整个参数就被跳过（假定不生成 |tabular| 环境）。

    使用 |\pgfplotstableset{column type/.add=|\marg{before}\marg{after}|}|
    来\emph{修改}一个值，而不是覆盖它。|/.add| 键处理器对其他选项同样有效。
    %
\begin{codeexample}[narrow]
\pgfplotstabletypeset[
    columns={dof,error1,info},
    column type/.add={|}{}% results in '|c'
]
    {pgfplotstable.example1.dat}
\end{codeexample}
\end{key}

\begin{key}{/pgfplots/table/column name=\marg{\TeX{} display column name} (initially \textbackslash pgfkeysnovalue)}
    在当前上下文中设置列的显示名。

    建议在某个特定列的样式内部提供此选项，也就是使用

    |columns/|\marg{lowlevel colname}|/.style={column name=|\marg{\TeX{} display column name}|}| .

    这里，``lowlevel colname'' 指的是出现在你输入表格中的列名。这个底层列名有一些限制
    （例如它必须可展开——也就是说很多控制序列是被禁止的）。|column name| 的值只在
    |\pgfplotstabletypeset| 内部使用，用于为相关列生成一个\emph{显示}名。

    注意，你可以用 |column name={}| 提供一个\emph{空}显示名。在 |\pgfplotstabletypeset|
    中使用时，这会在该列上方得到一个空字符串。

    初始配置是 |column name=\pgfkeysnovalue|。这构成一个特殊的 ``空'' 值，告诉
    |\pgfplotstabletypeset| 应回退到该列的名字（即底层名字）。
\end{key}

\begin{codekey}{/pgfplots/table/assign column name}
    允许\emph{修改} |column name| 的值。

    参数 |#1| 是当前的列名，也就是在计算完 |column name| 之后的值。在
    |assign column name| 之后，应设置一个新的（可能已修改的）|column name| 值。

    也就是说，你可以用 |column name| 来直接指定名字本身，而用 |assign column name|
    来生成最终的 \TeX{} 代码（例如插入 |\multicolumn{1}{c}{#1}|）。
    %
\begin{codeexample}[]
\pgfplotstabletypeset[
    assign column name/.style={
        /pgfplots/table/column name={\textbf{#1}}
    },
]
{
first second
1 2
3 4
}
\end{codeexample}

    注意，\PGFPlotstable{} 对表头提供的后处理能力有限，特别是 |postproc cell content|
    及其同类只适用于 ``普通'' 单元格。如果你想修改表头单元格的外观，应考虑使用
    |assign column name|。

    默认为空，表示不做改动。
\end{codekey}

\begin{stylekey}{/pgfplots/table/multicolumn names=\marg{tabular column type} (initially c)}
    一个样式，它在当前上下文中用
    |\multicolumn{1}|\marg{tabular column type}\marg{the column name} 语句来排版每个列名。

    这里，\meta{the column name} 照例用 |column name| 来设置。
\end{stylekey}

\begin{stylekey}{/pgfplots/table/dec sep align=\marg{header column type} (initially c)}
    一个样式，它把数值型列按小数点对齐。

    第一个参数决定表头列的对齐方式。

    请注意，使用这个样式需要 |\usepackage{array}|。
    %
\begin{codeexample}[]
% requires \usepackage{array}
\pgfplotstabletypeset[
    columns={dof,error1,error2,info,{grad(log(dof),log(error2))}},
    columns/error1/.style={dec sep align},
    columns/error2/.style={sci,sci subscript,sci zerofill,dec sep align},
    columns/info/.style={fixed,dec sep align},
    columns/{grad(log(dof),log(error2))}/.style={fixed,dec sep align}
]
    {pgfplotstable.example1.dat}
\end{codeexample}

    或者用逗号作为小数点：
    %
\begin{codeexample}[]
% requires \usepackage{array}
\pgfplotstabletypeset[
    use comma,
    columns={dof,error1,error2,info,{grad(log(dof),log(error2))}},
    columns/error1/.style={dec sep align},
    columns/error2/.style={sci,sci subscript,sci zerofill,dec sep align},
    columns/info/.style={fixed,dec sep align},
    columns/{grad(log(dof),log(error2))}/.style={fixed,dec sep align}
]
    {pgfplotstable.example1.dat}
\end{codeexample}
    %
    建议使用 |fixed zerofill| 和/或 |sci zerofill| 来强制小数点后至少有一位数字，
    以改善指数的位置：
    %
\begin{codeexample}[]
% requires \usepackage{array}
\pgfplotstabletypeset[
    use comma,
    columns={dof,error1,error2,info,{grad(log(dof),log(error2))}},
    columns/error1/.style={dec sep align,sci zerofill},
    columns/error2/.style={sci,sci subscript,sci zerofill,dec sep align},
    columns/info/.style={fixed,dec sep align},
    columns/{grad(log(dof),log(error2))}/.style={fixed,dec sep align,fixed zerofill}
]
    {pgfplotstable.example1.dat}
\end{codeexample}

    样式 |dec sep align| 实际上引入了两个新的 |tabular| 列，\footnote{遗憾的是，
    \texttt{dec sep align} 目前在列类型修改方面还不够灵活。特别是，无法把彩色列或
    单元格与 \texttt{dec sep align} 结合使用。\texttt{\textbackslash rowcolor} 命令可以
    正常工作；由 \texttt{colortbl} 引入的颜色残留会被自动调整。}即 |r@{}l|。
    它会相应地为列名引入多列，并处理没有小数点的数字。

    注意，对于定点数，也可以用 |fixed zerofill| 配合 |column type=r| 来获得类似效果，
    作为一种替代方案。

    请注意，这个样式会覆盖 |column type|、|assign cell content| 以及某些数字格式设置。
\end{stylekey}

%--------------------------------------------------
% FIXME : doesn't really work as intended:
% \begin{stylekey}{/pgfplots/table/dec sep align/no unbounded}
%     Changes the internal processing of |dec sep align| such that unbounded values like $NaN$ and $\pm \infty$ will be aligned in exactly the same way as the column name.
%
%     This avoids funny spacing around unbounded coordinates when used with |dec sep align|.

% NO GALLERY
% \begin{codeexample}[]
% \pgfplotstabletypeset{pgfplotstable.example4.dat}
% \vrule
%
% \pgfplotstabletypeset[
%     columns/b/.style={fixed,precision=1},
% ]
%     {pgfplotstable.example4.dat}
% \vrule
%
% \pgfplotstabletypeset[
%     columns/b/.style={fixed,precision=1,dec sep align},
% ]
%     {pgfplotstable.example4.dat}
% \vrule
%
% \pgfplotstabletypeset[
%     columns/b/.style={fixed,precision=1,dec sep align,dec sep align/no unbounded},
% ]
%     {pgfplotstable.example4.dat}
% \end{codeexample}
% \end{stylekey}
%--------------------------------------------------

\begin{stylekey}{/pgfplots/table/sci sep align=\marg{header column type} (initially c)}
    一个样式，它把数值型列按科学计数法中的指数对齐。

    第一个参数决定表头列的对齐方式。

    它的工作方式与 |dec sep align| 类似，即引入两个人为列 |r@{}l| 用于对齐。

    请注意，使用这个样式需要 |\usepackage{array}|。

    请注意，这个样式会覆盖 |column type|、|assign cell content| 以及某些数字格式设置。
\end{stylekey}

\begin{stylekey}{/pgfplots/table/dcolumn=\marg{tabular column type}\marg{type for column name} (initially \{D\{.\}\{.\}\{2\}\}\{c\})}
    一个样式，可与 David Carlisle 的 |dcolumn| 宏包配合使用。它也能实现按小数点对齐。
    不过，小数点必须恰好是一个字符，这与 `|{,}|'（|use comma| 的默认设置）不兼容。
\end{stylekey}

\begin{pgfplotskey}{table/sort=\marg{true,false} (initially false)}
    如果设为 |true|，|\pgfplotstabletypeset| 会在执行其操作之前先对表格排序。

    关于如何配置 |sort key| 和 |sort cmp|，参见 |\pgfplotstablesort| 的说明。
    %
\begin{codeexample}[]
\pgfplotstabletypeset[
    sort,sort key=error2,
    columns={dof,error1,error2},
    columns/error1/.style={sci,sci subscript,sci zerofill,dec sep align},
    columns/error2/.style={sci,sci subscript,sci zerofill,dec sep align},
]
    {pgfplotstable.example1.dat}
\end{codeexample}

    |sort| 机制在真正的排版例程开始之前应用，也就是说它的效果与你在排版表格之前
    手动调用 |\pgfplotstablesort| 相同（不过 |sort| 键的优点是会尊重 |include outfiles|
    缓存机制）。任何 |create on use| 规范都会在调用 |sort key| 之前被解析。

    %The |sort| key applies also to |\addplot table[sort,read completely]|.
\end{pgfplotskey}

\begin{stylekey}{/pgfplots/table/every first column}
    一个只为每个第一列安装的样式。
    %
\begin{codeexample}[narrow]
\pgfplotstabletypeset[
 every head row/.style={before row=\hline,after row=\hline\hline},
 every last row/.style={after row=\hline},
 every first column/.style={
  column type/.add={|}{}
 },
 every last column/.style={
  column type/.add={}{|}
 }]
    {pgfplotstable.example1.dat}
\end{codeexample}
\end{stylekey}

\begin{stylekey}{/pgfplots/table/every last column}
    一个只为每个最后一列安装的样式。
\end{stylekey}

\begin{stylekey}{/pgfplots/table/every even column}
    一个为每个列索引为偶数的列（从~$0$ 开始）安装的样式。
{
\pgfplotstableset{
    columns={dof,error1,{grad(log(dof),log(error2))},info},
    columns/error1/.style={
        column name=$L_2$,
        sci,sci zerofill,sci subscript,
        precision=3},
    columns/dof/.style={
        int detect,
        column name=\textsc{Dof}},
    columns/{grad(log(dof),log(error2))}/.style={
        column name=slopes $L_2$,
        fixed,fixed zerofill,
        precision=1}}

\begin{codeexample}[code only]
\pgfplotstableset{
    columns={dof,error1,{grad(log(dof),log(error2))},info},
    columns/error1/.style={
        column name=$L_2$,
        sci,sci zerofill,sci subscript,
        precision=3},
    columns/dof/.style={
        int detect,
        column name=\textsc{Dof}},
    columns/{grad(log(dof),log(error2))}/.style={
        column name=slopes $L_2$,
        fixed,fixed zerofill,
        precision=1}}
\end{codeexample}

\begin{codeexample}[narrow,graphic=white]
% requires \usepackage{colortbl}
\pgfplotstabletypeset[
 every even column/.style={
  column type/.add={>{\columncolor[gray]{.8}}}{}
}]
    {pgfplotstable.example1.dat}
\end{codeexample}
}
\end{stylekey}

\begin{stylekey}{/pgfplots/table/every column}
    一个应用于每一列的样式。
    %
\begin{codeexample}[]
\pgfplotstabletypeset[
    every column/.code={
        \ifnum\pgfplotstablecol>0
            \pgfkeysalso{column type/.add={|}{}}%
        \fi
    },
]
{
first second third
1 2 3
4 5 6
}
\end{codeexample}
\end{stylekey}

\begin{stylekey}{/pgfplots/table/every odd column}
    一个为每个列索引为奇数的列（从~$0$ 开始）安装的样式。
\end{stylekey}

\begin{command}{\pgfplotstablecol}
    在计算行或列选项期间，这个命令展开为当前列的索引。
\end{command}

\begin{command}{\pgfplotstablecolname}
    在计算列选项期间，这个命令展开为当前列的名字。它在 |\pgfplotstabletypeset|
    处理列样式（包括节~\ref{sec:pgfplotstable:preproc} 中讲解的预处理步骤）、
    准备输出单元格内容以及检查行谓词期间有效。
\end{command}

\begin{command}{\pgfplotstablerow}
\label{pgfplotstable:page:tablerow}
    在计算行或列选项期间，这个命令展开为当前行的索引。

    ``列选项的计算'' 适用于节~\ref{sec:data:processing} 中的任何样式，
    即预处理、排版和后处理。

    宏 |\pgfplotstablerow| 可以取 $0,1,2,\dotsc,n-1$ 中的任意值，
    其中 $n$ 是 |\pgfplotstablerows| 的值。

    ``行选项的计算'' 适用于诸如 |every last row| 之类的内容。

    注意，对于表头行，它会取特殊值 $-1$。
\end{command}

\begin{command}{\pgfplotstablecols}
    在计算行或列选项期间，这个命令展开为输出表中的列总数。
\end{command}

\begin{command}{\pgfplotstablerows}
    在计算\emph{列}期间，这个命令展开为\emph{输入}行的总数。
    你可以在 |row predicate| 内部使用它。

    在计算\emph{行}期间，这个命令展开为\emph{输出}行的总数。
\end{command}

\begin{command}{\pgfplotstablename}
    在 |\pgfplotstabletypeset| 期间，这个宏以顶层展开的形式包含表格的宏名。
    如果你不熟悉 ``顶层展开'' 和 `|\expandafter|'，那你大概永远也用不到这个宏。

    进阶用户可能会从类似

    |\expandafter\pgfplotstabletypeset\pgfplotstablename|.

    的表达式中受益。

    对于从磁盘加载（且没有显式指定宏名）的表格，它会展开为一个临时宏。
\end{command}


\subsection{配置行外观：样式}

下面这些样式允许在\emph{所有单元格内容都已分配之后}配置最终的表格代码。

\begin{key}{/pgfplots/table/before row=\marg{\TeX{} code}}
    包含一段 \TeX{} 代码，它会被安装在一行中第一个单元格之前。

    请记住，\PGFPlotstable{} 并不做什么神奇的事——它只是一个生成 |tabular| 环境的
    代码生成器。因此，你可以在这里添加任何你平时会写进 |tabular| 环境的 \TeX{} 代码。

    一个例子是多列标题，对此 \PGFPlotstable{} 没有自带的解决方案：
    %
\begin{codeexample}[]
% \usepackage{booktabs}
\pgfplotstabletypeset[
    column type=l,
    every head row/.style={
        before row={
            \toprule
            & \multicolumn{2}{c}{Singular} & \multicolumn{2}{c}{Plural}\\
        },
        after row=\midrule,
    },
    every last row/.style={
        after row=\bottomrule},
    columns/person/.style       ={column name=},
    columns/singGaeilge/.style  ={column name=Gaeilge},
    columns/pluralGaeilge/.style={column name=Gaeilge},
    columns/singEnglish/.style  ={column name=English},
    columns/pluralEnglish/.style={column name=English},
    col sep=&,row sep=\\,
    string type,
]{
person & singEnglish  & singGaeilge & pluralEnglish  & pluralGaeilge\\
1st    & at me        & agam        & at us          & againn\\
2st    & at you       & agat        & at you         & agaibh\\
3st    & at him       & aige        & at them        & acu\\
       & at her       & aici        &                &\\
}
\end{codeexample}
%
    \noindent 这个例子声明了很多选项，最后跟着一个简短的内联表格。|every head row|
    样式通过逐字的 |tabular| 代码配合 |booktabs| 横线来配置 |\multicolumn| 标题。
    在试验过程中，考虑使用 |debug| 或 |outfile| 键可能会有帮助。要修改列标题，
    需要用 |columns/...| 样式。

    有时，你会想把 |\multicolumn| 和 |\rowcolor| 结合起来。据我对 \LaTeX{} 的了解，
    这有点复杂：它要求在 |\multicolumn| 内部使用 |\columncolor|。和上面的例子一样，
    有必要对 \PGFPlotstable{} 生成的代码稍作修改。请记住，\PGFPlotstable{} 只是一个
    生成 |tabular| 环境的代码生成器——你想怎么改就怎么改。下面的例子演示了
    |\multicolumn| 与 |\rowcolor| 的结合。它取自一个（必须承认，相当进阶的）应用：
    %
\begin{codeexample}[]
\newcolumntype{C}{>{\centering\arraybackslash}p{6mm}}% a centered fixed-width-column
\pgfplotstabletypeset[
    col sep=&,
    row sep=\\,
    every head row/.style={
    % as in the previous example, this patches the first row:
        before row={
            \hline
            \rowcolor{lightgray}
            \multicolumn{3}{|>{\columncolor{lightgray}}c|}{Quantenzahlen} & Term-\\
            \rowcolor{lightgray}
        },
        after row=\hline,
    },
    every last row/.style={
        after row=\hline},
    % define column-specific styles:
    columns/n/.style={column type=|C,column name=$n$},
    columns/l/.style={column type=|C,column name=$\ell$},
    columns/lambda/.style={column type=|C,column name=$\lambda$},
    columns/text/.style={column type=|c|,column name=bezeichnung,
        string type % <-it contains formatted data
    },
] {
n & l & lambda & text\\
1 & 0 & 0 & $1 s\sigma$ \\
2 & 0 & 0 & $2 s\sigma$ \\
2 & 1 & 0 & $2 p\sigma$ \\
2 & 1 & 1 & $2 p \pi $\\
3 & 2 & 0 & $3 d\sigma$ \\
3 & 2 & 2 & $3 d\delta$ \\
}
\end{codeexample}

    除数字格式（它实际上会调用 |\pgfmathprintnumber|）之外，上面的代码等价于下面的列表
    %
\begin{codeexample}[code only]
\newcolumntype{C}{>{\centering\arraybackslash}p{6mm}}% a centered fixed-width-column
\begin{tabular}{|C|C|C|c|}
\hline
\rowcolor{lightgray} \multicolumn{3}{|>{\columncolor{lightgray}}c|}{Quantenzahlen} & Term-\\
\rowcolor{lightgray} $n$ & $\ell$ & $\lambda$ & bezeichnung\\
\hline
$1$ & $0$ & $0$ & $1 s\sigma$ \\
$2$ & $0$ & $0$ & $2 s\sigma$ \\
$2$ & $1$ & $0$ & $2 p\sigma$ \\
$2$ & $1$ & $1$ & $2 p \pi$ \\
$3$ & $2$ & $0$ & $3 d\sigma$ \\
$3$ & $2$ & $2$ & $3 d\delta$ \\
\hline
\end{tabular}
\end{codeexample}

    显然，定义大量样式所带来的额外开销，只有在你需要数字打印、自动化处理，
    或拥有一大批相似表格时才值得。
\end{key}

\begin{key}{/pgfplots/table/after row=\marg{\TeX{} code}}
    包含一段 \TeX{} 代码，它会被安装在一行中最后一个单元格之后（即在 |\\| 之后）。
\end{key}

\begin{stylekey}{/pgfplots/table/every even row}
    一个为每个行索引为偶数的行安装的样式。第一行被认为是 ``表头'' 行，不计入。
    索引从~$0$ 开始。
{
\pgfplotstableset{
    columns={dof,error1,{grad(log(dof),log(error2))}},
    columns/error1/.style={
        column name=$L_2$,
        sci,sci zerofill,sci subscript,
        precision=3},
    columns/dof/.style={
        int detect,
        column name=\textsc{Dof}},
    columns/{grad(log(dof),log(error2))}/.style={
        column name=slopes $L_2$,
        fixed,fixed zerofill,
        precision=1},
}

\begin{codeexample}[code only]
\pgfplotstableset{
    columns={dof,error1,{grad(log(dof),log(error2))}},
    columns/error1/.style={
        column name=$L_2$,
        sci,sci zerofill,sci subscript,
        precision=3},
    columns/dof/.style={
        int detect,
        column name=\textsc{Dof}},
    columns/{grad(log(dof),log(error2))}/.style={
        column name=slopes $L_2$,
        fixed,fixed zerofill,
        precision=1}}
\end{codeexample}

\begin{codeexample}[narrow,graphic=white]
% requires \usepackage{booktabs}
\pgfplotstabletypeset[
    every head row/.style={
        before row=\toprule,after row=\midrule},
    every last row/.style={
        after row=\bottomrule},
]
    {pgfplotstable.example1.dat}
\end{codeexample}

\begin{codeexample}[narrow,graphic=white]
% requires \usepackage{booktabs,colortbl}
\pgfplotstabletypeset[
    every even row/.style={
        before row={\rowcolor[gray]{0.9}}},
    every head row/.style={
        before row=\toprule,after row=\midrule},
    every last row/.style={
        after row=\bottomrule},
]
    {pgfplotstable.example1.dat}
\end{codeexample}
}
\end{stylekey}

\begin{stylekey}{/pgfplots/table/every odd row}
    一个为每个行索引为奇数的行安装的样式。第一行被认为是 ``表头'' 行，不计入。
    索引从~$0$ 开始。
\end{stylekey}

\begin{stylekey}{/pgfplots/table/every head row}
    一个为 tabular 中每个第一行安装的样式。它可用于调整列名的选项，
    或添加额外的横线/颜色。

\begin{codeexample}[]
\pgfplotstabletypeset[
  % suppress the header row 'col1  col2  col3':
  every head row/.style={output empty row},
  col sep=comma,
  columns/col1/.style={string type,column type=r},
  columns/col2/.style={string type,column type=l},
  columns/col3/.style={string type,column type=l},
  ]
{
    col1,col2,col3
    Col A,B,C
    The first column,E,F
}
\end{codeexample}
\end{stylekey}

\begin{stylekey}{/pgfplots/table/every first row}
    一个为每个第一\emph{数据}行（即在表头行之后）安装的样式。
\end{stylekey}

\begin{stylekey}{/pgfplots/table/every last row}
    一个为每个最后一行安装的样式。
\end{stylekey}

\begin{stylekey}{/pgfplots/table/every row no \meta{index}}
    一个为索引为 \meta{index} 的行安装的样式。
\end{stylekey}

\begin{keylist}{%
    /pgfplots/table/every nth row=\marg{integer}\marg{options},
    /pgfplots/table/every nth row=\marg{integer\oarg{shift}}\marg{options}%
}
    这允许为每第 $n$ 行安装 \meta{options}，其中 $n=$\meta{integer}。
    %
\begin{codeexample}[]
\pgfplotstabletypeset[
    every nth row={3}{before row=\midrule},
    every head row/.style={
        before row=\toprule,after row=\midrule},
    every last row/.style={
        after row=\bottomrule},
]{
    a b
    0 0
    1 1
    2 2
    3 3
    4 4
    5 5
    6 6
    7 7
    8 8
}
\end{codeexample}
    %
    |every nth row| 只考虑数据行；它永远不会应用于列名，数据行按 $i=0,1,2,\dotsc$
    编号（上面的例子把它应用于 $a = 3,6$ 的行）。由于 $i=0$ 的情形可由
    |every first row| 处理，它不纳入 |every nth row| 的考虑。

    第二种语法允许提供一个额外的 \meta{shift}：
\begin{codeexample}[]
\pgfplotstabletypeset[
    every nth row={3[+1]}{before row=\midrule},
]{
    a b
    0 0
    1 1
    2 2
    3 3
    4 4
    5 5
    6 6
    7 7
    8 8
    9 9
    10 10
}
\end{codeexample}
    %
    \noindent 这里，该样式应用于行 $i=1,4,7,10$（数学上，当 $(i \mod n) = $
    \meta{shift} 时应用它）。\meta{shift} 可以为负。

    你可以定义多个 |every nth row| 样式，它们按出现顺序处理（考虑使用
    |before row/.add=|\marg{before existing}\marg{after existing} 来修改已有的值）。

    注意，|every nth row/.style 2 args=...| 与 |every nth row=...| 相同。
\end{keylist}

\begin{stylekey}{/pgfplots/table/output empty row}
    一个抑制当前行输出的样式。

    这个样式会在很晚的时候才计算，即在所有特定列的内容修改都已应用之后。
    它等价于
    %
\begin{codeexample}[code only]
\pgfplotstableset{
    output empty row/.style={
        typeset cell/.style={@cell content={}}
    },
}
\end{codeexample}

    一个应用示例参见 |every head row|。
\end{stylekey}


\subsection{配置单个单元格外观：样式}

除了修改列样式和行样式的各种方式之外，还有一些样式可以修改单个单元格。

\begin{stylekey}{/pgfplots/table/every row \meta{rowindex} column \meta{colindex}}
    一个最多应用于一个单元格的样式：即行索引为 \meta{rowindex}、列索引为
    \meta{colindex} 的那个单元格。这些索引都从~$0$ 开始。

    这个样式在与 |preproc cell content|、|assign cell content| 和
    |postproc cell content| 键相同的上下文中计算，修改这些参数中的任何一个都是
    正当做法。也可以通过给 |@cell content| 赋一个值来替换初始单元格值。

    例如，看下面这个未经修改的表格：
    %
\begin{codeexample}[]
\pgfplotstabletypeset[
    col sep=&,row sep=\\]{
  colA & colB & colC \\
  11   & 12   & 13   \\
  21   & 22   & 23   \\
}
\end{codeexample}

    现在，我们修改其中一个单元格的数字格式，同时修改另一个（单个）单元格的格式：
    %
\begin{codeexample}[]
\pgfplotstabletypeset[
    every row 1 column 2/.style={/pgf/number format/sci},
    every row 0 column 0/.style={postproc cell content/.style={@cell content=\textbf{##1}}},
    col sep=&,row sep=\\]{
  colA & colB & colC \\
  11   & 12   & 13   \\
  21   & 22   & 23   \\
}
\end{codeexample}
    注意，这个样式（只）应用于输入的行/列值。
\end{stylekey}

\begin{stylekey}{/pgfplots/table/every row no \meta{rowindex} column no \meta{colindex}}
    这实际上是一样的——|row no| 和 |row| 都受支持，|column| 和 |column no| 也是如此。
\end{stylekey}

\begin{stylekey}{/pgfplots/table/every row \meta{rowindex} column \meta{colname}}
    一个与上面类似的样式，但它允许用列\emph{名}而不是列索引。列名需要用与其他
    特定列样式相同的方式提供（如果 \meta{colname} 包含特殊字符，还要加上额外的
    花括号）。

    上面的例子因此可以变成：
\begin{codeexample}[]
\pgfplotstabletypeset[
    every row 1 column colB/.style={string replace*={2}{4}},
    every row 0 column colA/.style={preproc/expr={##1*8}},
    col sep=&,row sep=\\]{
  colA & colB & colC \\
  11   & 12   & 13   \\
  21   & 22   & 23   \\
}
\end{codeexample}
    这个例子使用了 |string replace*| 预处理器样式和 |preproc/expr| 样式。
    所有预处理器或后处理器样式都可以使用。

    预定义的选项参见节~\ref{sec:data:processing}。
\end{stylekey}


\subsection{定制并获取 Tabular 代码}

下面这些键允许修改对齐方式（|begin table|）和 |font|，并允许把生成的代码写入
|outfile|（另见 |write to macro|）。此外，还可以对生成的代码做精细调整，
以提供 \LaTeX{} 之外的其他种类表格输出。

\begin{stylekey}{/pgfplots/table/every table=\marg{file name}}
    一个在每个 |\pgfplotstabletypeset| 命令开头安装的样式。\footnote{\texttt{every table}
    样式是在提供给 \texttt{\textbackslash pgfplotstabletypeset} 的选项\emph{之后}
    安装的；它优先级更高。}

    表格文件名作为第一个参数给出。
\end{stylekey}

\begin{key}{/pgfplots/table/font=\marg{font name} (initially empty)}
    指定用于整个表格的字体。
\end{key}

\begin{key}{/pgfplots/table/begin table=\marg{code} (initially \textbackslash begin\{tabular\})}
    包含一段 \marg{code}，它作为表格的开头生成。

    下面的例子用 |longtable| 代替\index{longtable} |tabular|：
    %
\begin{codeexample}[code only]
\pgfplotstableset{
    begin table=\begin{longtable},
    end table=\end{longtable},
}
\end{codeexample}

    注意，|longtable| 允许对多页表格\index{Multi-page} 借助其
    \declareandlabel{\endhead} 宏使用\emph{重复表头}：\index{Replicate headers}
    %
\begin{codeexample}[code only]
% replicate column names on top of every page of a multi-page table:
\pgfplotstableset{
    row sep=\\,
    begin table=\begin{longtable},
    end table=\end{longtable},
    every head row/.append style={after row=\endhead},
}
\end{codeexample}

    如果希望第一页有不同的表头，可以用 |longtable| 宏包提供的
    \declareandlabel{\endfirsthead}：
    %
\begin{codeexample}[code only]
% replicate column names on top of every page of a multi-page table,
% but with different values for first page:
\pgfplotstableset{
    row sep=\\,
    begin table=\begin{longtable},
    end table=\end{longtable},
    every head row/.append style={after row={%
            \caption{The caption}%
            \endfirsthead
            \multicolumn{3}{c}{{\bfseries \tablename\ \thetable{} -- continued from previous page}} \\
            \endhead
        },
    },
}
\end{codeexample}
    %
    \noindent 上一个例子用到了 |longtable| 的宏 |\caption|、|\endfirsthead|、|\thetable|
    和 |\endhead|。此外，它还要求\emph{显式}给出列数（这里是 |{3}|）。

    也可以\emph{修改} |begin table| 的值。例如，
\begin{codeexample}[code only]
\pgfplotstableset{
    begin table/.add={}{[t]},
}
\end{codeexample}
    %
    会在前面 prepend 空字符串 |{}|，并在后面 append 前缀 |[t]|。
    于是 `|\begin{tabular}|' 变成 `|\begin{tabular}[t]|'。
\end{key}

\begin{key}{/pgfplots/table/end table=\marg{code} (initially \textbackslash end\{tabular\})}
    包含一段 \meta{code}，它作为表格的结尾生成。
\end{key}

\begin{codekey}{/pgfplots/table/typeset cell}
    一个代码键，它把 \declareandlabel{/pgfplots/table/@cell content} 赋给当前单元格的
    最终输出。

    第一个参数 |#1| 是该单元格的最终值。在这个宏之后，|@cell content| 的值会被写入输出。

    默认实现是
    %
\begin{codeexample}[code only]
/pgfplots/table/typeset cell/.code={%
    \ifnum\c@pgfplotstable@colindex=\c@pgfplotstable@numcols\relax
        \pgfkeyssetvalue{/pgfplots/table/@cell content}{#1\\}%
    \else
        \pgfkeyssetvalue{/pgfplots/table/@cell content}{#1&}%
    \fi
},
\end{codeexample}

    \paragraph{注意：}

    在这个上下文中，|\pgfplotstablecol| 的值从 $1$ 开始，也就是说它取 $1,\dotsc,n$
    范围内的值，其中 $n=$ |\pgfplotstablecols|。这简化了对是否为最后一列的判断。
\end{codekey}

\begin{key}{/pgfplots/table/outfile=\marg{file name} (initially empty)}
\label{page:outfile}
    把生成的 tabular 代码写入 \meta{file name}。之后就可以用
    |\input|\marg{file name} 来使用它，此时不再需要 \PGFPlotstable{}，因为它包含的是
    一个完全普通的 |tabular|。
    %
    \pgfplotstableset{begin table=\begin{tabular}}
\begin{codeexample}[]
\pgfplotstabletypeset[
    columns={dof,error1},
    outfile=pgfplotstable.example1.out.tex]
    {pgfplotstable.example1.dat}
\end{codeexample}
    %
    而 |pgfplotstable.example1.out.tex| 的内容为
    %
    %\lstdefineformat{inp}{\\\\=\string\newline}%
    \lstinputlisting[basicstyle=\ttfamily\footnotesize]{pgfplotstable.example1.out.tex}

    命令 |\pgfutilensuremath| 会检查数学模式是否处于激活状态，并在必要时切换到数学模式。
    \footnote{请注意，如果你想脱离 \PGF{} 独立使用这个输出文件，需要把
    \lstinline{\\pgfutilensuremath} 替换为 \lstinline{\\ensuremath}。这可以通过
    \lstinline{\\let\\pgfutilensuremath=\\ensuremath} 来完成，从而启用 \LaTeX{} 命令
    \lstinline{\\ensuremath}。}
\end{key}

\begin{key}{/pgfplots/table/include outfiles=\marg{boolean} (initially false)}
    如果启用，任何已存在的 outfile 都会被 |\input| 进来，而不是被覆盖。
    %
\begin{codeexample}[code only]
\pgfplotstableset{include outfiles} % for example in the document's preamble
\end{codeexample}
    %
    这允许你手动把任何修正放进生成的输出文件中，因为 \PGFPlotstable{} 不会自动覆盖
    所得表格。

    这会影响那些设置了 |outfile| 选项的表格。如果你想把它应用到每个表格，可以考虑
    %
\begin{codeexample}[code only]
\pgfplotstableset{every table/.append style={outfile={#1.out}}}
\end{codeexample}
    %
    \noindent 它会为每个表格生成一个 |outfile| 名字。
\end{key}

\begin{key}{/pgfplots/table/force remake=\marg{boolean} (initially false)}
    如果启用，|include outfiles| 的效果就会被禁用。由于所有键设置只持续到下一个
    花括号（或 |\end|\meta{}）为止，这个键可用于重新生成某些输出文件，而其他文件仍被
    包含进来。
\end{key}

\begin{key}{/pgfplots/table/write to macro=\marg{\textbackslash macroname}}
    如果 |write to macro| 的值非空，完整生成的（tabular）代码将被写入宏
    \meta{\textbackslash macroname}。

    如果你\emph{只}需要所得的宏，参见 |typeset=false| 键。
\end{key}

\begin{key}{/pgfplots/table/skip coltypes=\mchoice{true,false} (initially false)}
    允许跳过 |\begin{tabular}|\marg{coltypes} 中的 \meta{coltypes}。
    这可以为那些没有 \LaTeX{} 表格格式的其他表格类型提供简化。
\end{key}

\begin{key}{/pgfplots/table/typeset=\mchoice{true,false} (initially true)}
    一个布尔值，用于禁用最后的排版阶段。如果只关心生成的代码、而不希望 \TeX{} 在
    输出 |pdf| 中产生任何内容，请配合 |write to macro| 使用 |typeset=false|。
\end{key}

\begin{key}{/pgfplots/table/debug=\marg{boolean} (initially false)}
    如果启用，它会把每个最终的 tabular 代码写入你的日志文件。
\end{key}

\begin{key}{/pgfplots/table/TeX comment=\marg{comment sign} (initially \%)}
    插入 outfile 中用于抑制尾随空白的注释符号。
\end{key}

\noindent 作为最后一个例子，我们用 \PGFPlotstable{} 来写一个 |.html| 文件
（还包括数字格式和四舍五入！）：
%
% \usepackage{listings}
\begin{codeexample}[width=8cm]
\pgfplotstabletypeset[
    begin table={<table>}, end table={</table>},
    typeset cell/.style={
      /pgfplots/table/@cell content={<td>#1</td>}
    },
    before row=<tr>,after row=</tr>,
    skip coltypes, typeset=false,
    verbatim,% configures number printer
    TeX comment=,
    columns={level,dof,error1},
    outfile=pgfplotstable.example1.out.html,
]{pgfplotstable.example1.dat}
\lstinputlisting
    [basicstyle=\ttfamily\footnotesize]
    {pgfplotstable.example1.out.html}
\end{codeexample}


\subsection{为 \texttt{tabular} 定义列类型}

除了输入文本文件之外，有时还希望为已有的 \texttt{tabular} 环境定义列类型。

\begin{command}{\newcolumntype\marg{letter}\oarg{number of
    arguments}$>$\marg{before column}\meta{column type}$<$\marg{after column}
}
    命令 |\newcolumntype| 是 |array| 宏包的一部分，它定义一个新的列类型
    \meta{letter}，供 \LaTeX{} tabular 环境使用。
    %
\begin{codeexample}[code only]
\usepackage{array}
\end{codeexample}

\begin{codeexample}[]
\newcolumntype{d}{>{-}c<{+}}
\begin{tabular}{dl}
a & b \\
c & d \\
\end{tabular}
\end{codeexample}

    现在，可以在 \meta{before column} 和 \meta{after column} 中使用环境
    |pgfplotstablecoltype| 来定义数值型列：
    %
% \usepackage{array}
\begin{codeexample}[]
% requires \usepackage{array}
\newcolumntype{L}[1]
    {>{\begin{pgfplotstablecoltype}[#1]}r<{\end{pgfplotstablecoltype}}}

\begin{tabular}{L{int detect}L{sci,sci subscript,sci zerofill}}
9      & 2.50000000e-01\\
25     & 6.25000000e-02\\
81     & 1.56250000e-02\\
289    & 3.90625000e-03\\
1089   & 9.76562500e-04\\
4225   & 2.44140625e-04\\
16641  & 6.10351562e-05\\
66049  & 1.52587891e-05\\
263169 & 3.81469727e-06\\
1050625& 9.53674316e-07\\
\end{tabular}
\end{codeexample}
    %
    \noindent 环境 |pgfplotstablecoltype| 接受一个可选参数，其中可包含任何数字格式选项。
    如果数值型列中含有非数值数据，就会出错，因此列名可能需要用 |\multicolumn| 来处理。
    %
\begin{codeexample}[]
% requires \usepackage{array}
\newcolumntype{L}[1]
    {>{\begin{pgfplotstablecoltype}[#1]}r<{\end{pgfplotstablecoltype}}}

\begin{tabular}{L{int detect}L{sci,sci subscript,sci zerofill}}
\multicolumn{1}{r}{Dof} & \multicolumn{1}{r}{Error}\\
9      & 2.50000000e-01\\
25     & 6.25000000e-02\\
81     & 1.56250000e-02\\
289    & 3.90625000e-03\\
1089   & 9.76562500e-04\\
4225   & 2.44140625e-04\\
16641  & 6.10351562e-05\\
66049  & 1.52587891e-05\\
263169 & 3.81469727e-06\\
1050625& 9.53674316e-07\\
\end{tabular}
\end{codeexample}
\end{command}


\subsection{数字格式选项}
\label{sec:number:printing}

下面这段摘自 \cite{tikz}，解释如何配置数字格式。公共选项前缀
|/pgf/number format| 可以省略；它会被自动识别。

所有这些数字格式选项也可应用于 \PGFPlots{}。

\begin{command}{\pgfmathprintnumber\marg{x}}
    为（实）数 \meta{x} 生成美观排版的输出。输入数 \meta{x} 用
    |\pgfmathfloatparsenumber| 解析，它支持任意精度。

    数字在数学模式下用当前这一组数字打印选项排版，见下文。也可以用
    |\pgfmathprintnumber[|\meta{options}|]|\marg{x} 来提供可选参数。
\end{command}

\begin{command}{\pgfmathprintnumberto\marg{x}\marg{\textbackslash macro}}
    把所得的数返回给 \meta{\textbackslash macro}，而不是直接排版它。
\end{command}

\begin{key}{/pgf/number format/fixed}
    配置 |\pgfmathprintnumber| 把数字四舍五入到小数点后固定位数，
    并丢弃任何尾随零。
    %
\begin{codeexample}[]
\pgfkeys{/pgf/number format/.cd,fixed,precision=2}
\pgfmathprintnumber{4.568}\hspace{1em}
\pgfmathprintnumber{5e-04}\hspace{1em}
\pgfmathprintnumber{0.1}\hspace{1em}
\pgfmathprintnumber{24415.98123}\hspace{1em}
\pgfmathprintnumber{123456.12345}
\end{codeexample}

    如何修改外观参见节~\ref{sec:number:styles}。
\end{key}

\begin{key}{/pgf/number format/fixed zerofill=\marg{boolean}  (default true)}
    启用或禁用对任何以定点格式绘制的数字的补零。
    %
\begin{codeexample}[]
\pgfkeys{/pgf/number format/.cd,fixed,fixed zerofill,precision=2}
\pgfmathprintnumber{4.568}\hspace{1em}
\pgfmathprintnumber{5e-04}\hspace{1em}
\pgfmathprintnumber{0.1}\hspace{1em}
\pgfmathprintnumber{24415.98123}\hspace{1em}
\pgfmathprintnumber{123456.12345}
\end{codeexample}
    %
    这个键影响用 |fixed| 或 |std| 样式绘制的数字（后者仅在未选择科学格式时）。
    %
\begin{codeexample}[]
\pgfkeys{/pgf/number format/.cd,std,fixed zerofill,precision=2}
\pgfmathprintnumber{4.568}\hspace{1em}
\pgfmathprintnumber{5e-05}\hspace{1em}
\pgfmathprintnumber{1}\hspace{1em}
\pgfmathprintnumber{123456.12345}
\end{codeexample}

    如何修改外观参见节~\ref{sec:number:styles}。
\end{key}

\begin{key}{/pgf/number format/sci}
    配置 |\pgfmathprintnumber| 以科学格式显示数字，即符号、尾数和指数（基数~$10$）。
    尾数会被四舍五入到所需的 |precision|（或 |sci precision|，见下文）。
    %
\begin{codeexample}[]
\pgfkeys{/pgf/number format/.cd,sci,precision=2}
\pgfmathprintnumber{4.568}\hspace{1em}
\pgfmathprintnumber{5e-04}\hspace{1em}
\pgfmathprintnumber{0.1}\hspace{1em}
\pgfmathprintnumber{24415.98123}\hspace{1em}
\pgfmathprintnumber{123456.12345}
\end{codeexample}

    关于如何改变指数显示样式，见节~\ref{sec:number:styles}。
\end{key}

\begin{key}{/pgf/number format/sci zerofill=\marg{boolean}  (default true)}
    启用或禁用对任何以科学格式绘制的数字的补零。
    %
\begin{codeexample}[]
\pgfkeys{/pgf/number format/.cd,sci,sci zerofill,precision=2}
\pgfmathprintnumber{4.568}\hspace{1em}
\pgfmathprintnumber{5e-04}\hspace{1em}
\pgfmathprintnumber{0.1}\hspace{1em}
\pgfmathprintnumber{24415.98123}\hspace{1em}
\pgfmathprintnumber{123456.12345}
\end{codeexample}
    %
    和 |fixed zerofill| 一样，这个选项只影响以 |sci| 格式绘制的数字
    （或在选择了科学格式时的 |std|）。

    如何修改指数显示样式参见节~\ref{sec:number:styles}。
\end{key}

\begin{stylekey}{/pgf/number format/zerofill=\marg{boolean} (default true)}
    同时设置 |fixed zerofill| 和 |sci zerofill|。
\end{stylekey}

\begin{keylist}{/pgf/number format/std,%
    /pgf/number format/std=\meta{lower e},
    /pgf/number format/std=\meta{lower e}:\meta{upper e}%
}
    把 |\pgfmathprintnumber| 配置为一种标准算法。它根据数量级选择 |fixed| 或 |sci|。
    设 $n=s \cdot m \cdot 10^e$ 为输入数，$p$ 为当前精度。如果
    $-p/2 \le e \le 4$，就用 |fixed| 格式显示该数；否则用 |sci| 格式显示。

\begin{codeexample}[]
\pgfkeys{/pgf/number format/.cd,std,precision=2}
\pgfmathprintnumber{4.568}\hspace{1em}
\pgfmathprintnumber{5e-04}\hspace{1em}
\pgfmathprintnumber{0.1}\hspace{1em}
\pgfmathprintnumber{24415.98123}\hspace{1em}
\pgfmathprintnumber{123456.12345}
\end{codeexample}
    %
    可以用可选的整数参数定制这些参数：如果 $\text{\meta{lower e}} \le e \le
    \text{\meta{upper e}}$，就以 |fixed| 格式显示该数，否则以 |sci| 格式显示。
    注意，为得到有用的结果，\meta{lower e} 应为负。科学格式所用的精度必要时可用
    |sci precision| 调整。

\end{keylist}

\begin{keylist}{/pgf/number format/relative*=\meta{exponent base 10}}
    配置 |\pgfmathprintnumber| 相对于某个数量级 $10^r$ 来格式化数字，
    其中 $r$ 是一个整数。

    这个键针对不同的使用场景。


    \paragraph{第一个使用场景：}

    为一\emph{序列}数字提供统一的格式。考虑下面这个测试：
    %
\begin{codeexample}[]
\pgfkeys{/pgf/number format/relative*={1}}
\pgfmathprintnumber{6.42e-16}\hspace{1em}
\pgfmathprintnumber{1.2}\hspace{1em}
\pgfmathprintnumber{6}\hspace{1em}
\pgfmathprintnumber{20.6}\hspace{1em}
\pgfmathprintnumber{87}
\end{codeexample}
    %
    \noindent 用任何其他样式，|6.42e-16| 都会被作为一个孤立的数字来格式化。在这里，
    它被四舍五入为 |0|，因为相对于 $10^1$（指数 $1$ 是传给 |relative| 的参数）来看，
    它没有有效数字。
    %
\begin{codeexample}[]
\pgfkeys{/pgf/number format/relative*={2}}
\pgfmathprintnumber{123.345}\hspace{1em}
\pgfmathprintnumber{0.0012}\hspace{1em}
\pgfmathprintnumber{0.0014}\hspace{1em}
\end{codeexample}
    %
    \noindent 上面的例子把初始的 |precision=2| 应用到 |123.345|——相对于 $100$。
    |123.345| 相对于 $100$ 的两位有效数字是 |123|。注意，``|123.345| 的 $2$ 位
    有效数字'' 等同于 ``把 |1.2345| 四舍五入到 $2$ 位''，结果会是 |1.2300|。
    类似地，另外两个数相对于 $100$、用给定的 |precision| 来看都是 |0|。
    %
\begin{codeexample}[]
\pgfkeys{/pgf/number format/relative*={-3}}
\pgfmathprintnumber{123.345}\hspace{1em}
\pgfmathprintnumber{0.0012}\hspace{1em}
\pgfmathprintnumber{0.0014}\hspace{1em}
\end{codeexample}

    \paragraph{第二个使用场景：}

    在存在\emph{不精确}数字时改善四舍五入。假设某种有限精度的运算得到了结果
    |123456999|（就像 \pgfname 的 |fpu| 那样）。你知道它的精度大约是五到六位
    有效数字。而你想输出定点数。在这种情况下，尾随的几位 |....999| 是有限精度
    造成的数值假象。用 |relative*=3,precision=0| 来消除这些假象：
    %
\begin{codeexample}[]
\pgfkeys{/pgf/number format/.cd,relative*={3},precision=0}
\pgfmathprintnumber{123456999}\hspace{1em}
\pgfmathprintnumber{123456999.12}
\end{codeexample}
    %
    \noindent 这里，|precision=0| 表示我们考察 |123456.999| 并把它四舍五入到 $0$ 位。
    最后再把小数点移回初始位置。加上 |relative style=fixed| 就得到定点输出格式：
    %
\begin{codeexample}[]
\pgfkeys{/pgf/number format/.cd,relative*={3},precision=0,relative style=fixed}
\pgfmathprintnumber{123456999}\hspace{1em}
\pgfmathprintnumber{123456999.12}
\end{codeexample}
    %
    \noindent 注意，这个使用场景还有另一种替代方案，稍后会讨论：即 |fixed relative| 样式。
    %
\begin{codeexample}[]
\pgfkeys{/pgf/number format/.cd,fixed relative,precision=6}
\pgfmathprintnumber{123456999}\hspace{1em}
\pgfmathprintnumber{123456999.12}
\end{codeexample}

    你也许会问，为什么这个键的名字里有个星号。简短回答是：还有一个
    \declareandlabel{/pgf/number format/relative} 数字打印器，它会做出一些出乎意料的
    行为。键 |relative*| 修正了这一点。已有代码仍会沿用旧的行为。

    从技术上讲，这个键的工作方式如下：如前所述，作用于 |123456999.12| 的
    |relative*=3| 键把小数点移动三位，并考察 |123456.99912|。用数学语言说，
    给定一个数 $x = \pm m \cdot 10^e$，我们尝试应用 |relative*=|$r$。该方法随后把
    $x / 10^r$ 四舍五入到 |precision| 位。之后再把结果乘以 $10^r$ 并排版。
\end{keylist}

\begin{stylekey}{/pgf/number format/every relative}
    一个样式，它配置 |relative| 方法最终如何显示其结果。

    初始配置是
    %
\begin{codeexample}[code only]
\pgfkeys{/pgf/number format/every relative/.style=std}
\end{codeexample}

    注意，在计算所得样式时舍入是关闭的（因为 |relative| 已经对该数做过舍入）。

    尽管受支持，我不建议在此情形下使用 |fixed zerofill| 或 |sci zerofill|——
    它可能造成精度比实际所用更高的错觉（因为 |fixed zerofill| 可能只是简单地补上
    |.00|，而在 |relative| 对结果舍入之前原本还有别的信息）。
\end{stylekey}

\begin{key}{/pgf/number format/relative style=\marg{options}}
    等同于 |every relative/.append style=|\marg{options}。
\end{key}

\begin{keylist}{/pgf/number format/fixed relative}
    配置 |\pgfmathprintnumber| 以与 |fixed| 样式类似的方式格式化数字，
    但 |precision| 是相对于该数的指数来解释的。

    其动机是获得与 |sci| 相同的四舍五入效果，但以 |fixed| 样式显示该数：
    %
\begin{codeexample}[]
\pgfkeys{/pgf/number format/.cd,fixed relative,precision=3}
\pgfmathprintnumber{1000.0123}\hspace{1em}
\pgfmathprintnumber{100.0567}\hspace{1em}
\pgfmathprintnumber{0.000010003452}\hspace{1em}
\pgfmathprintnumber{0.010073452}\hspace{1em}
\pgfmathprintnumber{1.23567}\hspace{1em}
\pgfmathprintnumber{1003.75}\hspace{1em}
\pgfmathprintnumber{1006.75}\hspace{1em}
\end{codeexample}

    |fixed relative| 的效果是：无论该数有多少个前导零，它都把该数四舍五入到
    \emph{恰好}前 \meta{precision} 位非零数字。

    如果你想要 |fixed|、并且知道只有前 $n$ 位数字是正确的，就用 |fixed relative|。
    如果你需要科学显示样式、且只有前 $n$ 位数字是正确的，就用 |sci|。

    注意，|fixed relative| 会忽略 |fixed zerofill| 标志。

    另见 |relative*| 键。注意，|relative=|\marg{exponent} 键在尝试对数字四舍五入之前，
    会先把小数点显式地移到某个指定位置。之后它 ``从右往左舍入''，也就是舍入到那个
    显式选定的数字位置。与之相反，|fixed relative| ``从左往右舍入''：它取\emph{第一个}
    非零数字，临时把小数点放在这个数字之后，然后对该数四舍五入。舍入样式 |fixed|
    则保持小数点不动，只对该数字之后的一切进行舍入。|sci| 样式与 |fixed relative| 类似。
\end{keylist}

\begin{key}{/pgf/number format/int detect}
    配置 |\pgfmathprintnumber| 自动检测整数。如果输入数是整数，则完全不显示小数点。
    否则选择科学格式。

\begin{codeexample}[]
\pgfkeys{/pgf/number format/.cd,int detect,precision=2}
\pgfmathprintnumber{15}\hspace{1em}
\pgfmathprintnumber{20}\hspace{1em}
\pgfmathprintnumber{20.4}\hspace{1em}
\pgfmathprintnumber{0.01}\hspace{1em}
\pgfmathprintnumber{0}
\end{codeexample}
\end{key}

\begin{command}{\pgfmathifisint\marg{number constant}\marg{true code}\marg{false code}}
    一个命令，它做与 |int detect| 相同的检查，但如果 \meta{number constant} 确实是
    整数，就执行 \meta{true code}，否则执行 \meta{false code}。

    作为副作用，|\pgfretval| 会包含解析后的数字，要么是整数格式，要么是解析出的
    浮点数。

    参数 \meta{number constant} 会用 |\pgfmathfloatparsenumber| 解析。
\begin{codeexample}[]
15 \pgfmathifisint{15}{is an int: \pgfretval.}{is no int}\hspace{1em}
15.5 \pgfmathifisint{15.5}{is an int: \pgfretval.}{is no int}
\end{codeexample}
\end{command}

\begin{key}{/pgf/number format/int trunc}
    把每个数字截断为整数（丢弃小数点后的任何数字）。
    %
\begin{codeexample}[]
\pgfkeys{/pgf/number format/.cd,int trunc}
\pgfmathprintnumber{4.568}\hspace{1em}
\pgfmathprintnumber{5e-04}\hspace{1em}
\pgfmathprintnumber{0.1}\hspace{1em}
\pgfmathprintnumber{24415.98123}\hspace{1em}
\pgfmathprintnumber{123456.12345}
\end{codeexample}
\end{key}

\begin{key}{/pgf/number format/frac}
    把数字显示为分数。
    %
\begin{codeexample}[width=3cm]
\pgfkeys{/pgf/number format/frac}
\pgfmathprintnumber{0.333333333333333}\hspace{1em}
\pgfmathprintnumber{0.5}\hspace{1em}
\pgfmathprintnumber{2.133333333333325e-01}\hspace{1em}
\pgfmathprintnumber{0.12}\hspace{1em}
\pgfmathprintnumber{2.666666666666646e-02}\hspace{1em}
\pgfmathprintnumber{-1.333333333333334e-02}\hspace{1em}
\pgfmathprintnumber{7.200000000000000e-01}\hspace{1em}
\pgfmathprintnumber{6.666666666666667e-02}\hspace{1em}
\pgfmathprintnumber{1.333333333333333e-01}\hspace{1em}
\pgfmathprintnumber{-1.333333333333333e-02}\hspace{1em}
\pgfmathprintnumber{3.3333333}\hspace{1em}
\pgfmathprintnumber{1.2345}\hspace{1em}
\pgfmathprintnumber{1}\hspace{1em}
\pgfmathprintnumber{-6}
\end{codeexample}

    \begin{key}{/pgf/number format/frac TeX=\marg{\textbackslash macro} (initially \texttt{\textbackslash frac})}
        允许在 |frac| 显示类型内部使用不同的 |\frac| 实现。
    \end{key}

    \begin{key}{/pgf/number format/frac denom=\meta{int} (initially empty)}
        允许为 |frac| 提供自定义分母。
        %
\begin{codeexample}[width=3cm]
\pgfkeys{/pgf/number format/.cd,frac, frac denom=10}
\pgfmathprintnumber{0.1}\hspace{1em}
\pgfmathprintnumber{0.5}\hspace{1em}
\pgfmathprintnumber{1.2}\hspace{1em}
\pgfmathprintnumber{-0.6}\hspace{1em}
\pgfmathprintnumber{-1.4}\hspace{1em}
\end{codeexample}
    \end{key}

    \begin{key}{/pgf/number format/frac whole=\mchoice{true,false} (initially true)}
        配置是否应把完整的整数部分放在分数部分之前。这样分数部分会小于 $1$。
        用 |frac whole=false| 来避免整数部分。
        %
\begin{codeexample}[width=3cm]
\pgfkeys{/pgf/number format/.cd,frac, frac whole=false}
\pgfmathprintnumber{20.1}\hspace{1em}
\pgfmathprintnumber{5.5}\hspace{1em}
\pgfmathprintnumber{1.2}\hspace{1em}
\pgfmathprintnumber{-5.6}\hspace{1em}
\pgfmathprintnumber{-1.4}\hspace{1em}
\end{codeexample}
    \end{key}

    \begin{key}{/pgf/number format/frac shift=\marg{integer} (initially 4)}
        如果你遇到稳定性问题，可以尝试换一个 |frac shift|。较大的 shift 值 $k$
        会使结果对不精确的数据或不精确的运算更敏感。

        从技术上讲，过程如下。如果 $r < 1$ 是尾数的小数部分，则计算一个尺度
        $i = 1/r \cdot 10^k$，其中 $k$ 是 shift；$i$ 的小数部分被忽略。$1/r$ 是在
        内部计算的，其误差会被放大。

        如果你仍然遇到稳定性问题，请在导言区中启用 |fp|。这样 |frac| 样式就会在
        计算 $1/r$ 时自动采用 |fp| 更高的绝对精度。
    \end{key}
\end{key}

\begin{key}{/pgf/number format/precision=\marg{number}}
    为任何显示操作设置期望的舍入精度。对于科学记数格式，这会影响尾数。
\end{key}

\begin{key}{/pgf/number format/sci precision=\meta{number or empty} (initially empty)}
    仅为 |sci| 样式设置期望的舍入精度。

    用 |sci precision={}| 可恢复初始配置（即对所有数字样式都使用提供给 |precision|
    的参数）。
\end{key}

\begin{key}{/pgf/number format/read comma as period=\mchoice{true,false} (initially false)}
    这是少数几个允许定制数字解析器的键之一。如果打开这个开关，逗号会被当作小数点来读。
    %
\begin{codeexample}[]
\pgfkeys{/pgf/number format/read comma as period}
\pgfmathprintnumber{1234,56}
\end{codeexample}
    %
    这通常是不希望的，因为它可能与数学解析指令产生副作用。不过，它支持用于格式化
    输入数字或输入表格。考虑用 |use comma| 让结果也以逗号排版。
    %
\begin{codeexample}[]
\pgfkeys{/pgf/number format/.cd,
    read comma as period,
    use comma}
\pgfmathprintnumber{1234,56}
\end{codeexample}

    \paragraph{注意：}

    这个键需要比~\pgfname~3.0.0 更新的 \pgfname\ 版本。
\end{key}


\subsubsection{修改数字格式的显示样式}%
\label{sec:number:styles}

你可以修改数字的显示方式。例如，如果你使用 `\texttt{fixed}' 样式，输入数会被
四舍五入到所需精度，并用当前的定点显示样式来排版该数。其他任何格式也同样如此：
先用四舍五入例程得到正确的数字，再由一个显示样式生成合适的 \TeX{} 代码。

\begin{key}{/pgf/number format/set decimal separator=\marg{text}}
    把 \meta{text} 指定为任何定点数（包括 sci 格式中的尾数）的小数点。
\end{key}

\begin{key}{/pgf/number format/dec sep=\marg{text}}
    只是 |set decimal separator| 的另一个名字。
\end{key}

\begin{key}{/pgf/number format/set thousands separator=\marg{text}}
    把 \meta{text} 指定为任何定点数（包括 sci 格式中的尾数）的千位分隔符。
    %
\begin{codeexample}[]
\pgfkeys{/pgf/number format/.cd,
    fixed,
    fixed zerofill,
    precision=2,
    set thousands separator={}}
\pgfmathprintnumber{1234.56}
\end{codeexample}

\begin{codeexample}[]
\pgfkeys{/pgf/number format/.cd,
    fixed,
    fixed zerofill,
    precision=2,
    set thousands separator={}}
\pgfmathprintnumber{1234567890}
\end{codeexample}

\begin{codeexample}[]
\pgfkeys{/pgf/number format/.cd,
    fixed,
    fixed zerofill,
    precision=2,
    set thousands separator={.}}
\pgfmathprintnumber{1234567890}
\end{codeexample}

\begin{codeexample}[]
\pgfkeys{/pgf/number format/.cd,
    fixed,
    fixed zerofill,
    precision=2,
    set thousands separator={,}}
\pgfmathprintnumber{1234567890}
\end{codeexample}

\begin{codeexample}[]
\pgfkeys{/pgf/number format/.cd,
    fixed,
    fixed zerofill,
    precision=2,
    set thousands separator={{{{,}}}}}
\pgfmathprintnumber{1234567890}
\end{codeexample}
    %
    最后一个例子使用逗号，并禁用了默认的逗号间距。
\end{key}

\begin{key}{/pgf/number format/1000 sep=\marg{text}}
    只是 |set thousands separator| 的另一个名字。
\end{key}

\begin{key}{/pgf/number format/1000 sep in fractionals=\marg{boolean} (initially false)}
    配置小数部分是否也应按三位一组分组。

    取值 |true| 会对整数部分和小数部分都启用 |1000 sep|。取值 |false| 则只对整数部分
    启用 |1000 sep|。
    %
\begin{codeexample}[]
\pgfkeys{/pgf/number format/.cd,
    fixed,
    precision=999,
    set thousands separator={\,},
    1000 sep in fractionals,
}
\pgfmathprintnumber{1234.1234567}
\end{codeexample}

\begin{codeexample}[]
\pgfkeys{/pgf/number format/.cd,
    fixed,fixed zerofill,
    precision=9,
    set thousands separator={\,},
    1000 sep in fractionals,
}
\pgfmathprintnumber{1234.1234567}
\end{codeexample}
\end{key}

\begin{key}{/pgf/number format/min exponent for 1000 sep=\marg{number} (initially 0)}
    定义绘制千位分隔符所需的科学记数法中的最小指数。指数等于位数减一，因此
    $\mathit{number}=4$ 会从 $1e4 = 10000$ 起使用千位分隔符。
    %
\begin{codeexample}[]
\pgfkeys{/pgf/number format/.cd,
    int detect,
    1000 sep={\,},
    min exponent for 1000 sep=0}
\pgfmathprintnumber{5000}; \pgfmathprintnumber{1000000}
\end{codeexample}

\begin{codeexample}[]
\pgfkeys{/pgf/number format/.cd,
    int detect,
    1000 sep={\,},
    min exponent for 1000 sep=4}
\pgfmathprintnumber{1000}; \pgfmathprintnumber{5000}
\end{codeexample}

\begin{codeexample}[]
\pgfkeys{/pgf/number format/.cd,
    int detect,
    1000 sep={\,},
    min exponent for 1000 sep=4}
\pgfmathprintnumber{10000}; \pgfmathprintnumber{1000000}
\end{codeexample}
    %
    \noindent 取值 |0| 会禁用这个特性（负值会被忽略）。
\end{key}


\begin{key}{/pgf/number format/use period}
    一个预定义样式，它把句号 `\texttt{.}' 装为小数点、把逗号 `\texttt{,}' 装为
    千位分隔符。这个样式是默认的。
    %
\begin{codeexample}[]
\pgfkeys{/pgf/number format/.cd,fixed,precision=2,use period}
\pgfmathprintnumber{12.3456}
\end{codeexample}

\begin{codeexample}[]
\pgfkeys{/pgf/number format/.cd,fixed,precision=2,use period}
\pgfmathprintnumber{1234.56}
\end{codeexample}
\end{key}

\begin{key}{/pgf/number format/use comma}
    一个预定义样式，它把逗号 `\texttt{,}' 装为小数点、把句号 `\texttt{.}' 装为
    千位分隔符。
    %
\begin{codeexample}[]
\pgfkeys{/pgf/number format/.cd,fixed,precision=2,use comma}
\pgfmathprintnumber{12.3456}
\end{codeexample}

\begin{codeexample}[]
\pgfkeys{/pgf/number format/.cd,fixed,precision=2,use comma}
\pgfmathprintnumber{1234.56}
\end{codeexample}
\end{key}

\begin{key}{/pgf/number format/skip 0.=\marg{boolean} (initially false)}
    配置像 $0.1$ 这样的数字是否应排版为 $.1$。
    %
\begin{codeexample}[]
\pgfkeys{/pgf/number format/.cd,
    fixed,
    fixed zerofill,precision=2,
    skip 0.}
\pgfmathprintnumber{0.56}
\end{codeexample}

\begin{codeexample}[]
\pgfkeys{/pgf/number format/.cd,
    fixed,
    fixed zerofill,precision=2,
    skip 0.=false}
\pgfmathprintnumber{0.56}
\end{codeexample}
\end{key}

\begin{key}{/pgf/number format/showpos=\marg{boolean} (initially false)}
    启用或禁用对非负数显示加号。
    %
\begin{codeexample}[]
\pgfkeys{/pgf/number format/showpos}
\pgfmathprintnumber{12.345}
\end{codeexample}

\begin{codeexample}[]
\pgfkeys{/pgf/number format/showpos=false}
\pgfmathprintnumber{12.345}
\end{codeexample}

\begin{codeexample}[]
\pgfkeys{/pgf/number format/.cd,showpos,sci}
\pgfmathprintnumber{12.345}
\end{codeexample}
\end{key}

\begin{stylekey}{/pgf/number format/print sign=\marg{boolean}}
    一个样式，它就是 |showpos=|\meta{boolean} 的别名。
\end{stylekey}

\begin{key}{/pgf/number format/sci 10e}
    对任何以科学格式显示的数字使用 $m \cdot 10^e$。

\begin{codeexample}[]
\pgfkeys{/pgf/number format/.cd,sci,sci 10e}
\pgfmathprintnumber{12.345}
\end{codeexample}
\end{key}

\begin{key}{/pgf/number format/sci 10\textasciicircum e}
    与 `|sci 10e|' 相同。
\end{key}

\begin{key}{/pgf/number format/sci e}
    对任何以科学格式显示的数字，使用常见科学工具所生成的 `$1e{+}0$' 格式。

\begin{codeexample}[]
\pgfkeys{/pgf/number format/.cd,sci,sci e}
\pgfmathprintnumber{12.345}
\end{codeexample}
\end{key}

\begin{key}{/pgf/number format/sci E}
    与上相同，但用大写 `\texttt{E}'。
    %
\begin{codeexample}[]
\pgfkeys{/pgf/number format/.cd,sci,sci E}
\pgfmathprintnumber{12.345}
\end{codeexample}
\end{key}

\begin{key}{/pgf/number format/sci subscript}
    对任何以科学格式显示的数字，把指数排版为下标。这个样式占用的空间非常小。
    %
\begin{codeexample}[]
\pgfkeys{/pgf/number format/.cd,sci,sci subscript}
\pgfmathprintnumber{12.345}
\end{codeexample}
\end{key}

\begin{key}{/pgf/number format/sci superscript}
    对任何以科学格式显示的数字，把指数排版为上标。这个样式占用的空间非常小。
    %
\begin{codeexample}[]
\pgfkeys{/pgf/number format/.cd,sci,sci superscript}
\pgfmathprintnumber{12.345}
\end{codeexample}
\end{key}

\begin{key}{/pgf/number format/sci generic=\marg{keys}}
    允许为科学格式定义自定义数字样式。这里，\meta{keys} 可以是下列选项之一
    （省略长键前缀）：

    \begin{key}{/pgf/number format/sci generic/mantissa sep=\marg{text} (initially empty)}
        提供尾数与指数之间的分隔符。例如它可以是 |\cdot|，
    \end{key}

    \begin{key}{/pgf/number format/sci generic/exponent=\marg{text} (initially empty)}
        提供用于格式化指数的文本。实际的指数可作为参数 |#1| 获取（见下文）。
    \end{key}

\begin{codeexample}[]
\pgfkeys{
    /pgf/number format/.cd,
    sci,
    sci generic={mantissa sep=\times,exponent={10^{#1}}}}
\pgfmathprintnumber{12.345};
\pgfmathprintnumber{0.00012345}
\end{codeexample}
    %
    这些 \meta{keys} 可以依赖三个参数：|#1| 是指数，|#2| 包含浮点数的标志位，
    |#3| 是（未经处理、未格式化的）尾数。

    注意，|sci generic|\emph{并不}适合用来修改定点数的外观，也不能用来格式化尾数
    （尾数是像定点数那样排版的）。要定制尾数，请用 |dec sep|、|1000 sep| 和
    |print sign|。
\end{key}

\begin{key}{/pgf/number format/retain unit mantissa=\mchoice{true,false} (initially true)}
    允许省略单位尾数。
    %
\begin{codeexample}[]
\pgfkeys{
    /pgf/number format/.cd,
    sci, retain unit mantissa=false}
\pgfmathprintnumber{10.5};
\pgfmathprintnumber{10};
\pgfmathprintnumber{1010};
\pgfmathprintnumber[precision=1]{-1010};
\end{codeexample}
    %
    这个特性在四舍五入到所需精度之后应用：如果余下的尾数等于~$1$，就把它省略。
    它适用于所有涉及科学格式的样式（包括 |std|）。
\end{key}

\begin{key}{/pgf/number format/@dec sep mark=\marg{text}}
    会被放在小数点应在位置的紧前面。不过，即使没有小数点，也会插入 \meta{text}。
    它旨在作为辅助例程查找对齐位置的占位符。

    这个键绝不应用来修改小数点！请改用 |dec sep|。
\end{key}

\begin{key}{/pgf/number format/@sci exponent mark=\marg{text}}
    会被放在科学记数法中指数的紧前面。它旨在作为辅助例程查找对齐位置的占位符。

    这个键绝不应用来修改指数！
\end{key}

\begin{key}{/pgf/number format/assume math mode=\marg{boolean} (default true)}
    如果你不希望做任何数学模式检查，就把它设为 |true|。

    初始设置会在每个最终数字周围装一个 |\pgfutilensuremath|，以便在必要时切换到
    数学模式。如果你知道数学模式处于激活状态、且不想要 |\pgfutilensuremath|，
    就用 |assume math mode=true|。
\end{key}

\begin{stylekey}{/pgf/number format/verbatim}
    一个样式，它配置数字打印器产生逐字文本输出，也就是不包含 \TeX{} 宏。
    %
\begin{codeexample}[]
\pgfkeys{
    /pgf/fpu,
    /pgf/number format/.cd,
    sci,
    verbatim}
\pgfmathprintnumber{12.345};
\pgfmathprintnumber{0.00012345};
\pgfmathparse{exp(15)}
\pgfmathprintnumber{\pgfmathresult}
\end{codeexample}
    %
    这个样式会重置 |1000 sep|、|dec sep|、|print sign|、|skip 0.| 并设置
    |assume math mode|。此外，它还为科学数字的逐字输出装一个 |sci generic| 格式。

    不过，它仍会尊重 |precision|、|fixed zerofill|、|sci zerofill| 以及总体样式
    |fixed|、|sci|、|int detect|（及其变体）。如果你想写输出文件，它可能有用。
\end{stylekey}


\section{从输入数据到输出表格：数据处理}
\label{sec:data:processing}

从一个未经处理的输入表格转换到最终排版好的 |tabular| 代码，对每个单元格都经过
四个阶段，
%
\begin{enumerate}
    \item 加载表格，
    \item 预处理，
    \item 排版，
    \item 后处理。
\end{enumerate}
%
其主要思想是选定一个排版算法（例如 ``用配置好的数字样式格式化我的数字''）。
这个算法通常不需要改动。然后可以用零个、一个或多个预处理器和后处理器来做精细调整。
预处理可以指只选择特定的行，或在排版算法看到内容之前施加某种操作。后处理则指对所得的
\TeX{} 输出做精细调整——例如处理空单元格、插入单位后缀，或修改单个单元格的字体。

注意，这类操作是逐列处理的。因此，预处理、排版和后处理的键适用于列上的样式。
它们对 |every last row| 之类的行样式不可用。如果你需要在特定行上施加这类基于内容的
操作，就得写类似 ``如果 |\pgfplotstablerow| $=$ 就以不同方式生成内容'' 的代码。


\subsection{加载表格}

排版表格的这第一步涉及显而易见的输入操作。此外，节~\ref{pgfplotstable:createcol}
中讲解的 ``新建列'' 操作也在此时处理。表格数据如本手册前面所述那样被读取（或获取）。
然后，如果缺少列，列别名和 |create on use| 规范会作为加载过程的一部分被处理。
关于新建列的细节参见节~\ref{pgfplotstable:createcol}。


\subsection{排版单元格内容}

排版单元格就是取它们的值并 ``做点什么''。在很多情况下，这会涉及数字格式例程。
例如，``原始'' 输入数据 |12.56| 可能变成 |1.26| |\cdot| |10^1|。这一阶段的结果不再
适用于基于内容的计算。排版步骤在预处理步骤之后进行。

排版单元格内容是逐列应用的，也就是说它也能为行样式配置。如何只在特定行上应用这类
样式，细节见节~\ref{sec:data:processing} 引言中的说明。

\begin{codekey}{/pgfplots/table/assign cell content}
    允许重新定义分配单元格内容的算法。参数 |#1| 是输入表的（未格式化的）内容。

    所得的输出需要写入 |/pgfplots/table/@cell content|。
    %
\pgfplotstableset{
    %debug=true,
    begin table=\begin{tabular},% eliminate the '[b]' from the preamble here
}
% \usepackage{booktabs}
% \usepackage{multirow}
\begin{codeexample}[]
% An example how to use
% \usepackage{multirow} and
% \usepackage{booktabs}:
\pgfplotstabletypeset[
    columns/Z/.style={
        column name={},
        assign cell content/.code={% use \multirow for Z column:
            \ifnum\pgfplotstablerow=0
                \pgfkeyssetvalue{/pgfplots/table/@cell content}
                    {\multirow{4}{*}{##1}}%
            \else
                \pgfkeyssetvalue{/pgfplots/table/@cell content}{}%
            \fi
        },
    },
    % use \booktabs as well (compare examples above):
    every head row/.style={before row=\toprule,after row=\midrule},
    every last row/.style={after row=\bottomrule},
    row sep=\\,col sep=&,
    outfile=pgfplotstable.multirow.out,% write it to file
]{% here: inline data in tabular format:
    Z    & a & b \\
    data & 1 & 2 \\
         & 3 & 4 \\
         & 5 & 6 \\
         & 7 & 8 \\
}
% ... and show the generated file:
\lstinputlisting[basicstyle=\footnotesize\ttfamily]{pgfplotstable.multirow.out}
\end{codeexample}
        \index{multirow}%
    %
    \noindent 上面的例子用 |\usepackage{multirow}| 来格式化列 |Z|。更准确地说，
    它对列 |Z| 的第 \#0 行使用 |\multirow{4}{*}{data}|，而对列 |Z| 的其他行使用
    空字符串。

    请注意，你可能需要特别留意 |#1=|\marg{}，也就是空字符串。如果某一列的行数比
    第一列少，就会出现这种情况。在这种情况下，\PGFPlotstable{} 会自动平衡各列，
    插入足够多的空单元格，使其与第一列的行数相匹配。

    还要注意，如果某一列的条目比第一列多，这些多余的条目会被跳过，并在日志文件中
    发出一条警告信息。

    这个键是在一个局部 \TeX{} 分组内部计算的，因此之后任何局部宏赋值都会被清除。
\end{codekey}

\begin{stylekey}{/pgfplots/table/numeric type}
    一个样式，它把 |assign cell content|（重）定义回其原始值，即假定数据为数值型。

    它调用 |\pgfmathprintnumberto|，并把结果写入 |@cell content|。
\end{stylekey}

\begin{stylekey}{/pgfplots/table/string type}
    一个样式，它把 |assign cell content| 重定义为直接返回 ``原始'' 输入数据，
    也就是作为文本列。这假定输入表格含有合法的 \LaTeX{} 内容（不支持逐字打印）。
\end{stylekey}

\begin{stylekey}{/pgfplots/table/verb string type}
    一个样式，它把 |assign cell content| 重定义为按原样返回 ``原始'' 内容。
    因此它与 |string type| 相当类似——但它会返回控制序列以及（许多，而非全部）
    特殊字符，而不展开它们。

    你可能需要把 |verb string type| 与 |special chars| 结合使用。
\end{stylekey}

\begin{stylekey}{/pgfplots/table/numeric as string type}
    一个样式，它把 |assign cell content| 重定义为假定输入数据为数值型。它返回一个
    字符串字面量，以整数或科学（指数）记数法描述输入数。与 |numeric type| 不同，
    它不应用数字格式。
\end{stylekey}

\begin{stylekey}{/pgfplots/table/date type=\marg{date format}}% (initially \year-\month-\day)}
    一个样式，它在每个单元格中期望形如 |YYYY-MM-DD| 的 ISO 日期，并在输出时产生
    美观排版的字符串。输出格式由 \meta{date format} 给出。在 \meta{date format}
    内部，可以使用下面讲解的若干宏。
    %
% \usepackage{pgfcalendar}
\begin{codeexample}[]
% Requires
% \usepackage{pgfcalendar}
\pgfplotstableset{columns={date,account1}}

% plotdata/accounts.dat contains:
%
% date              account1  account2  account3
% 2008-01-03        60        1200      400
% 2008-02-06        120       1600      410
% 2008-03-15        -10       1600      410
% 2008-04-01        1800      500       410
% 2008-05-20        2300      500       410
% 2008-06-15        800       1920      410

% Show the contents in `string type':
\pgfplotstabletypeset[
    columns/date/.style={string type}
]{plotdata/accounts.dat}
\hspace{1cm}
% Show the contents in `date type':
\pgfplotstabletypeset[
    columns/date/.style={date type={\monthname\ \year}}
]{plotdata/accounts.dat}
\end{codeexample}
    %
    这个样式 \textbf{需要} 加载 \PGF{} \textbf{日历宏包}：
    %
\begin{codeexample}[code only]
\usepackage{pgfcalendar}
\end{codeexample}

    \begin{command}{\year}
        在 \meta{date format} 内部，这个宏展开为以数字表示的年份（如 |2008|）。
    \end{command}

    \begin{command}{\month}
        在 \meta{date format} 内部，这个宏展开为以数字表示的月份，从~$1$ 开始
        （如 |1|）。
    \end{command}

    \begin{command}{\monthname}
        在 \meta{date format} 内部，这个宏展开为当前语言所设定的月份名
        （如 |January|）。如何修改语言见下文。
    \end{command}

    \begin{command}{\monthshortname}
        在 \meta{date format} 内部，这个宏展开为当前语言所设定的月份缩写名
        （如 |Jan|）。如何修改语言见下文。
    \end{command}

    \begin{command}{\day}
        在 \meta{date format} 内部，这个宏展开为以数字表示的日（如 |31|）。
    \end{command}

    \begin{command}{\weekday}
        在 \meta{date format} 内部，这个宏展开为星期几的编号（周一为 $0$，
        周二为 $1$，依此类推）。
    \end{command}

    \begin{command}{\weekdayname}
        在 \meta{date format} 内部，这个宏展开为当前语言中的星期名
        （如 |Wednesday|）。如何修改语言见下文。
    \end{command}

    \begin{command}{\weekdayshortname}
        在 \meta{date format} 内部，这个宏展开为当前语言中的星期缩写名
        （如 |Wed|）。如何修改语言见下文。
    \end{command}


    \subsubsection*{修改日期所用的语言}

    日期功能是用 \PGF{} 日历模块实现的。这个模块会使用 |translator| 宏包
    （如果它已被加载）。抱歉，我还没有更多细节。更多信息请参见 \cite{tikz}。
\end{stylekey}


\subsection{预处理单元格内容}
\label{sec:pgfplotstable:preproc}

预处理步骤允许在应用任何排版例程（如数字格式）\emph{之前}修改单元格内容。
因此，如果表格含有数值数据，就可以在这一阶段施加数学运算。此外，还可以根据
单元格的数值来清除单元格。预处理步骤在数据获取步骤（``加载步骤''）之后进行。
这尤其意味着你可以创建（或复制）列并应用
operations on them.

预处理按逐列方式应用，也就是说它可以为行样式配置。关于如何只在特定行上应用这类样式，
细节见节~\ref{sec:data:processing} 引言中的说明。

\begin{codekey}{/pgfplots/table/preproc cell content}
    允许在调用 |assign cell content| \emph{之前}\emph{修改}单元格内容。

    其语义如下：在预处理程序之前，|@cell content| 含有原始输入数据（或者，
    可能是另一次预处理调用的结果）。预处理之后，|@cell content| 被填入一个——
    可能被修改过的——值。所得值随后用作 |assign cell content| 的输入。

    在默认设置下，|assign cell content| 期待数值型输入。因此，
    预处理程序应产生数值型输出。

    可以用 |/.append code| 或 |/.append style| 键处理器提供多个预处理指令。

    如果你不想要（或不需要）可叠加的预处理程序，也可以用 `|#1|' 来取得文件中
    所找到的原始输入数据。此外，键 |@unprocessed cell content| 也会包含该原始
    输入数据。
\end{codekey}

\begin{stylekey}{/pgfplots/table/string replace=\marg{cell match}\marg{cell replacement}}
    向当前的 |preproc cell content| 值追加代码，它把任何与 \meta{cell match}
    精确匹配的单元格替换为 \meta{cell replacement}。这一步不做任何展开；
    \meta{cell match} 必须逐字匹配。
    %
\begin{codeexample}[]
\pgfplotstabletypeset[columns={level,dof}]
    {pgfplotstable.example1.dat}

\pgfplotstabletypeset[
    columns={level,dof},
    columns/level/.style={string replace={A}{B}}, % does nothing because there is no cell 'A'
    columns/dof/.style={string replace={256}{-42}}]  % replace cell '256' with '-42'
    {pgfplotstable.example1.dat}
\end{codeexample}

    子串替换参见 |string replace*| 方法。
\end{stylekey}

\begin{stylekey}{/pgfplots/table/string replace*=\marg{pattern}\marg{replacement}}
    向当前的 |preproc cell content| 值追加代码，它把每一处出现的 \meta{pattern}
    替换为 \meta{replacement}。这一步不做任何展开；\meta{pattern} 必须逐字匹配。
    %
\begin{codeexample}[]
\pgfplotstabletypeset[
    string replace*={2}{6},
    col sep=&,row sep=\\]{
  colA & colB & colC \\
  11   & 12   & 13   \\
  21   & 22   & 23   \\
}
\end{codeexample}
\end{stylekey}

\begin{stylekey}{/pgfplots/table/clear infinite}
    向当前的 |preproc cell content| 值追加代码，它把每个无穷数替换为空字符串。
    这会清除任何含有 $\pm \infty$ 和 NaN 的单元格。
\end{stylekey}

\begin{stylekey}{/pgfplots/table/preproc/expr=\marg{math expression}}
    向当前的 |preproc cell content| 值追加代码，它对每个单元格计算
    \meta{math expression}。运算以浮点数进行。

    在 \meta{math expression} 内部，用下列表达式之一来取得当前单元格的值。
    %
    \begin{itemize}
        \item 字符串 `|##1|' 展开为该单元格在输入文件中所找到的内容，
            忽略前面的预处理程序。

            这通常就足够了。
        \item 命令 |\thisrow|\marg{the currently processed column name} 展开为
            当前单元格的内容。这还会包括前面预处理程序的结果。

            注意，在这种上下文中（预处理程序内部）的 |\thisrow{}| 不如在新建列
            例程上下文中那么强大：其参数必须与当前正在处理的列名精确匹配。
            你也可以用简写

            |\thisrow{\pgfplotstablecolname}|。

        \item 命令 |\pgfkeysvalueof{/pgfplots/table/@cell content}| 与此相同。
    \end{itemize}
    %
\begin{codeexample}[]
\pgfplotstabletypeset[
    columns={level},
    columns/level/.style={
        column name={$2\cdot \text{level}+4$},
        preproc/expr={2*##1 + 4}
    }
]
    {pgfplotstable.example1.dat}
\end{codeexample}

    空单元格不会被处理，因为含有 ``空数'' 的数学表达式会失败。

    注意，还有一个 |create col/expr|，它比 |preproc/expr| 更强大。
\end{stylekey}

\begin{stylekey}{/pgfplots/table/multiply by=\marg{real number}}
    向当前的 |preproc cell content| 值追加代码，它把每个单元格乘以
    \meta{real number}。运算以浮点数进行。
\end{stylekey}

\begin{stylekey}{/pgfplots/table/divide by=\marg{real number}}
    向当前的 |preproc cell content| 值追加代码，它把每个单元格除以
    \meta{real number}。运算以浮点数进行。
\end{stylekey}

\begin{stylekey}{/pgfplots/table/sqrt}
    向当前的 |preproc cell content| 值追加代码，它对每个非空单元格取平方根
    $\sqrt{\cdot}$。运算以浮点数进行。

    下面的例子复制列 |error1|，并对这个副本应用 |sqrt|。
    %
\begin{codeexample}[]
\pgfplotstableset{
    columns={error1,sqrterror1},
    create on use/sqrterror1/.style={create col/copy=error1},
    columns/error1/.style={column name=$\epsilon$},
    columns/sqrterror1/.style={sqrt,column name=$\sqrt \epsilon$},
    sci,sci 10e,precision=3,sci zerofill
}
\pgfplotstabletypeset{pgfplotstable.example1.dat}
\end{codeexample}
    %
    关于 |create on use| 的细节，请参见节~\ref{pgfplotstable:createcol}。
\end{stylekey}

\begin{stylekey}{/pgfplots/table/multiply -1}
    向当前的 |preproc cell content| 值追加代码，它把每个单元格乘以 $-1$。
    这个样式与 |multiply by=-1| 做同样的事，只是更快，因为只改变符号。
    %
\begin{codeexample}[]
\pgfplotstableset{
    columns={dof,error2,slopes2},
    columns/error2/.style={sci,sci zerofill},
    columns/slopes2/.style={dec sep align,empty cells with={\ensuremath{-}}},
    create on use/slopes2/.style=
        {create col/gradient loglog={dof}{error2}}}

\pgfplotstabletypeset{pgfplotstable.example1.dat}

\pgfplotstabletypeset[columns/slopes2/.append style={multiply -1}]
    {pgfplotstable.example1.dat}
\end{codeexample}
\end{stylekey}

\begin{codekey}{/pgfplots/table/row predicate}
    一个布尔谓词，它允许根据当前行的索引来选择输入表格中的特定行。参数 |#1| 含有
    当前行的索引（从~$0$ 开始，不计注释行或列名）。

    返回值被赋给 \TeX{} 条件 \declareandlabel{\ifpgfplotstableuserow}。如果该布尔值
    未被改动，则返回值为真。
    %
\begin{codeexample}[newline=]
% requires \usepackage{booktabs}
\pgfplotstabletypeset[
    every head row/.style={
        before row=\toprule,after row=\midrule},
    every last row/.style={
        after row=\bottomrule},
    row predicate/.code={%
        \ifnum#1>4\relax
            \ifnum#1<8\relax
                \pgfplotstableuserowfalse
            \fi
        \fi}
]
    {pgfplotstable.example1.dat}
\end{codeexample}
    %
    请注意，|row predicate| 是在任何其他影响行（或列）外观的选项\emph{之前}应用的。
    它在 |assign cell content| 之前计算。其后果之一是：偶数/奇数行样式只针对那些
    谓词返回 |true| 的行。事实上，你可以用 |row predicate| 在表格真正被处理之前
    就截断整个表格。

    在 |row predicate| 期间，宏 |\pgfplotstablerows| 包含\emph{输入}行的总数。

    此外，|row predicate| 只适用于排版例程，不适用于读取方法。如果你想只用
    |\addplot table| 绘制选定的表格条目，请用 \PGFPlots{} 的坐标过滤选项。
\end{codekey}

\begin{stylekey}{/pgfplots/table/skip rows between index=\marg{begin}\marg{end}}
    一个样式，它追加一个 |row predicate|，用以丢弃选定的行。选择按索引进行，
    索引从~$0$ 开始。索引满足 $\mathit{begin} \le i < \mathit{end}$ 的每一行都会被跳过。
    %
\begin{codeexample}[narrow]
% requires \usepackage{booktabs}
\pgfplotstabletypeset[
    every head row/.style={
        before row=\toprule,after row=\midrule},
    every last row/.style={
        after row=\bottomrule},
    skip rows between index={2}{4},
    skip rows between index={7}{9}
]
    {pgfplotstable.example1.dat}
\end{codeexample}
\end{stylekey}

\begin{stylekey}{/pgfplots/table/select equal part entry of=\marg{part no}\marg{part count}}
    一个样式，它用一个子集选择谓词覆盖 |row predicate|。其思想是把当前列分成
    \meta{part count} 个等大的部分，只选择 \meta{part no}。

    这可用于模拟多列表格。
    %
\begin{codeexample}[]
% requires \usepackage{booktabs}
\pgfplotstableset{
    every head row/.style={before row=\toprule,after row=\midrule},
    every last row/.style={after row=\bottomrule}}

\pgfplotstabletypeset[string type]{pgfplotstable.example2.dat}%
~
\pgfplotstabletypeset[
    columns={A,B,A,B},
    display columns/0/.style={select equal part entry of={0}{2},string type},% first part of `A'
    display columns/1/.style={select equal part entry of={0}{2},string type},% first part of `B'
    display columns/2/.style={select equal part entry of={1}{2},string type},% second part of `A'
    display columns/3/.style={select equal part entry of={1}{2},string type},% second part of `B'
]
    {pgfplotstable.example2.dat}
\end{codeexample}
    %
    上面的例子左侧原样显示原始文件。右侧显示列 A,B,A,B——但只显示一半元素，
    即按 \#2 中的 \#0 或 \#1 索引选出的那些。这些部分等大，除一个余数外。

    如果可用行数不能被 \meta{part count} 整除，剩余的条目会被均分到前面的几个部分中。
\end{stylekey}

\begin{stylekey}{/pgfplots/table/unique=\marg{column name}}
    一个样式，它追加一个 |row predicate|，用以抑制 \meta{column name} 中相同元素的
    连续出现。例如，如果 \meta{column name} 含有 |1,1,3,5,5,6,5,0|，应用 |unique|
    的结果是 |1,3,5,6,5,0|（最后那个 |5| 被保留——它前面紧挨着的不是另一个 |5|）。

    该算法用字符串记号比较来查找重复出现。\footnote{更准确地说，比较是用
    \texttt{\textbackslash ifx} 完成的，也就是单元格内容不会被展开。只会使用
    输入表格中所见的记号。}

    参数 \meta{column name} 可以是列名、索引、别名或 |create on use| 规范
    （最后一种不得依赖其他 |create on use| 语句）。不必提供一个属于输出的
    \meta{column name}。

    不过，|unique| 谓词必须能从第一列开始对所有列计算。这意味着把 |unique| 放在
    特定列样式的深处是错误的。
\end{stylekey}


\subsection{后处理单元格内容}

后处理步骤在排版阶段之后应用。其主要目的是施加最终的、不基于内容的格式指令。

后处理是逐列应用的，也就是说它也能为行样式配置。如何只在特定行上应用这类样式，
细节见节~\ref{sec:data:processing} 引言中的说明。

\begin{codekey}{/pgfplots/table/postproc cell content}
    允许在已分配的单元格内容被分配\emph{之后}对其\emph{修改}，可以是依赖内容的。
    可能的做法有：把负数画成红色、把单个条目排成粗体，或插入替换文本。

    这个键在 |assign cell content| \emph{之后}计算。其语义是修改一个已有的
    |@cell content| 值。

    如果你用 |/.append code| 或 |/.append style| 来定义它们，可以有不止一个
    |postproc cell content| 命令：
    %
\begin{codeexample}[]
% requires \usepackage{eurosym}
\pgfplotstabletypeset[
    column type=r,
    columns={dof,info},
    columns/info/.style={
        % stupid example for multiple postprocessors:
        postproc cell content/.append style={
            /pgfplots/table/@cell content/.add={$\bf}{$},
        },
        postproc cell content/.append style={
            /pgfplots/table/@cell content/.add={}{\EUR{}},
        }
    }]
    {pgfplotstable.example1.dat}
\end{codeexample}
    %
    上面的代码分两步修改 |@cell content|。其净效果是在前面 prepend ``|$\bf |''，
    在后面 append ``|$ \EUR|''。应当注意，|pgfkeys|（基本上）以相同方式处理
    |/.style| 和 |/.code|——二者都是简单的代码键，都可这样使用。你可以把它们与
    |/.append style| 和 |/.append code| 组合。细节请参见
    ~\cite[section about pgfkeys]{tikz}。

    与 |assign cell content| 中一样，代码可以求值像 |\pgfplotstablerow| 这样的辅助宏，
    以只改动特定条目。此外，后处理程序可以依赖未处理的单元格输入（即输入文件中所
    找到的、或由加载过程产生的输入）和/或预处理后的单元格值。这些值可通过下列方式获得：
    %
    \begin{itemize}
        \item 键 \declareandlabel{@unprocessed cell content}，它存储原始输入，
        \item 键 \declareandlabel{@cell content after rowcol styles}，它存储计算单元格
            特定样式之后 |@cell content| 的值，
        \item 键 \declareandlabel{@preprocessed cell content}，它存储预处理程序的结果，
        \item 键 \declareandlabel{@cell content}，它包含排版例程的结果，
        \item 简写 `|#1|'，它也是在输入表格中所找到的未处理输入参数。
    \end{itemize}
    %
    记住，你可以随时用

    |\pgfkeysvalueof{/pgfplots/table/@preprocessed cell content}|

    来访问键值。

    这允许完全基于上下文的格式选项。请记住，由于列平衡可能会出现空字符串——
    必要时请引入特殊处理。

    有一种特殊情况：当 |@cell content| 本身含有单元格分隔字符 `|&|' 时。
    在这种情况下，|postproc cell content| 会对 & 之前和之后的每一部分\emph{分别}调用，
    然后再插入 &。这使其能与那些在输出中创建人为列的特殊样式兼容（这是允许的，
    见 |dec sep align|）。为便于对每一部分分别处理，你可以用宏
    \declareandlabel{\pgfplotstablepartno}。它只在计算 |postproc cell content| 期间
    定义，并展开为当前部分的索引（从~$0$ 开始）。如果你的文本中没有 &，
    该值将始终为~$0$。

    这个键是在一个局部 \TeX{} 分组内部计算的，因此之后任何局部宏赋值都会被清除。

    下面的例子可用于在斜率列中插入一个短横线 $-$：
    %
\begin{codeexample}[]
\pgfplotstableset{
    create on use/slopes1/.style=
        {create col/gradient loglog={dof}{error1}}}

\pgfplotstabletypeset[
    columns={dof,error1,slopes1},
    columns/error1/.style={sci,sci zerofill},
    columns/slopes1/.style={
        postproc cell content/.append code={%
            \ifnum\pgfplotstablerow=0
                \pgfkeyssetvalue{/pgfplots/table/@cell content}{\ensuremath{-}}%
            \fi
        }%
    }]
    {pgfplotstable.example1.dat}
\end{codeexample}
    %
    由于这在更一般的情形下也可能有用，它已作为 |empty cells with| 样式提供。

    注意，现在表头行也会被这个键修改（如果你想修改表头单元格内容，见
    |assign column name|）。
\end{codekey}

\begin{stylekey}{/pgfplots/table/empty cells with=\marg{replacement}}
    向 |postproc cell content| 追加代码，它把任何空单元格替换为 \meta{replacement}。

    如果 |dec sep align| 处于激活状态，替换值只会插入到小数点之前的那一部分。
\end{stylekey}

\begin{stylekey}{/pgfplots/table/set content=\marg{content}}
    一个样式，它把 |postproc cell content| 重定义为始终返回值 \meta{content}。
\end{stylekey}

\begin{stylekey}{/pgfplots/table/fonts by sign=\marg{\TeX{} code for positive}\marg{\TeX{} code for negative}}
    向 |postproc cell content| 追加代码，它允许为正数和负数设置字体。

    参数 \meta{\TeX{} code for positive} 和 \meta{\TeX{} code for negative}
    会被插入到排版好的单元格内容紧前面。两种修改 \LaTeX{} 字体的方式都可以用：
    |\textbf|\marg{argument} 或 |{\bfseries |\marg{argument}|}| 方式。
    %

% \usepackage{pgfcalendar}
\begin{codeexample}[]
% Requires
% \usepackage{pgfcalendar}

% plotdata/accounts.dat contains:
%
% date              account1  account2  account3
% 2008-01-03        60        1200      400
% 2008-02-06        120       1600      410
% 2008-03-15        -10       1600      410
% 2008-04-01        1800      500       410
% 2008-05-20        2300      500       410
% 2008-06-15        800       1920      410

\pgfplotstabletypeset[
    columns={date,account1},
    column type=r,
    columns/date/.style={date type={\monthname\ \year}},
    columns/account1/.style={fonts by sign={}{\color{red}}}
]
    {plotdata/accounts.dat}
\end{codeexample}
    %
    事实上，这个样式的参数不必是字体改动。样式 |fonts by sign| 会把若干花括号和
    相应的参数插入 |@cell content|。更准确地说，对于负数它会得到
    |{|\meta{\TeX{} code for negative}|{|\meta{cell value}|}}|，而对于其他所有数字得到
    |{|\meta{\TeX{} code for positive}|{|\meta{cell value}|}}|。
\end{stylekey}


\section{在新表格或新列中生成数据}
\label{pgfplotstable:createcol}

可以从零开始创建新表格，也可以在表格从磁盘加载后修改它们。


\subsection{从零创建新表格}

\begin{commandlist}{%
    \pgfplotstablenew\oarg{options}\marg{row count}\marg{\textbackslash table},
    \pgfplotstablenew*\oarg{options}\marg{row count}\marg{\textbackslash table}%
}
    从零创建一个新表格。

    新表格将包含 |columns| 键中列出的所有列。对于 |\pgfplotstablenew|，
    |columns| 键需要在 \oarg{options} 中提供。对于 |\pgfplotstablenew*|，
    则使用 |columns| 的当前值，无论它在何处、何时被设置。

    此外，每个要生成的列都必须有 |create on use| 语句（见下一小节）。\footnote{目前，
    你至少需要提供一个列：对于完全空的表格，实现会出错。如果你不提供任何列名，
    会创建一个哑列。} 各列按其在 |columns| 中出现的顺序独立生成。一旦某一列完成，
    就可以用任何基本层级的访问机制来访问它。因此，你可以构建相互依赖的列。

    该表格将恰好包含 \meta{row count} 行。如果 \meta{row count} 是一条
    |\pgfplotstablegetrowsof| 语句，就会执行该语句并使用所得的行数。否则，
    \meta{row count} 会作为数字计算。
    %
\begin{codeexample}[]
% this key setting could be provided in the document's preamble:
\pgfplotstableset{
    % define how the 'new' column shall be filled:
    create on use/new/.style={create col/set list={4,5,6,7,...,10}}}
% create a new table with 11 rows and column 'new':
\pgfplotstablenew[columns={new}]{11}\loadedtable
% show it:
\pgfplotstabletypeset[empty cells with={---}]\loadedtable
\end{codeexample}

\begin{codeexample}[]
% create a new table with 11 rows and column 'new':
\pgfplotstablenew[
    % define how the 'new' column shall be filled:
    create on use/new/.style={create col/expr={factorial(15+\pgfplotstablerow)}},
    columns={new}]
    {11}
    \loadedtable
% show it:
\pgfplotstabletypeset\loadedtable
\end{codeexample}
\end{commandlist}

\begin{command}{\pgfplotstablevertcat\marg{\textbackslash table1}\marg{\textbackslash table2 or filename}}
\label{table:vertcat}
    把 \meta{\textbackslash table2} 的内容追加到 \meta{\textbackslash table1}
    （``纵向拼接''）。更准确地说，只有 \meta{\textbackslash table1} 中已存在的列
    会被追加，且 \meta{\textbackslash table1} 中存在的每一列都必须在
    \meta{\textbackslash table2} 中存在（或者必须有 |alias| 或 |create on use|
    规范来生成它们）。

    如果第二个参数是文件名，就从磁盘加载该文件。

    如果 \meta{\textbackslash table1} 不存在，就把 \meta{\textbackslash table2}
    复制到 \meta{\textbackslash table1}。
    %
\begin{codeexample}[code only]
\pgfplotstablevertcat{\output}{datafile1} % loads `datafile1' -> `\output'
\pgfplotstablevertcat{\output}{datafile2} % appends rows of datafile2
\pgfplotstablevertcat{\output}{datafile3} % appends rows of datafile3
\end{codeexample}

    \paragraph{说明：}

    输出表格 \meta{\textbackslash table1} 会在当前 \TeX{} 作用域中定义，
    并在之后被清除。当前 \TeX{} 作用域由额外的一对花括号界定。然而，每个
    \LaTeX{} 环境，以及遗憾地还有 \Tikz\ 的 |\foreach| 语句，都会引入 \TeX{} 作用域。

    \PGFPlots{} 有一些循环语句不会引入额外作用域。例如，
    %
\begin{codeexample}[code only]
\pgfplotsforeachungrouped \i in {1,2,...,10} {%
    \pgfplotstablevertcat{\output}{datafile\i} % appends `datafile\i' -> `\output'
}%
\end{codeexample}
    %
    这些循环宏在 \PGFPlots{} 手册的参考章节 ``Miscellaneous Commands'' 中有讲解。
\end{command}

\begin{command}{\pgfplotstableclear\marg{\textbackslash table}}
    清除一个表格。注意，在表格操作周围加上额外的花括号 `|{ ... }|' 要可靠得多——
    这些花括号界定了一个变量（包括表格）的作用域。
\end{command}


\subsection{从已有列创建新列}

\begin{command}{\pgfplotstablecreatecol\oarg{options}\marg{new col name}\marg{\textbackslash table}}
    创建一个名为 \meta{new col name} 的新列，并把它追加到一个已存在的表格
    \meta{\textbackslash table}。

    最终用户大概根本不需要直接使用 |\pgfplotstablecreatecol|——有一个高层框架
    |create on use|，它在内部调用该命令，并可通过简单的键值赋值来使用（见下文）。
    不过，本文档讲解如何利用已有列的值来填充新单元格。

    这个命令提供了一个灵活的框架来生成新列。它被设计为利用已有的值来创建新列——
    例如用逻辑或数值方法组合已有值。它提供对某一行的值、上一行的值和下一行值的
    快速访问。

    下面的文档面向所有想\emph{编写}专用列的人。它并不特别难，只是技术性较强，
    需要一些 |pgfkeys| 知识。如果你不喜欢它，可以改用预定义的列生成样式——
    并在 \meta{options} 中启用这些样式。

    列条目会用命令键 \declaretext{create col/assign} 来创建。它会对表格的每一行调用。
    它应把内容赋给 \declaretext{create col/next content}。在计算期间，宏
    |\thisrow|\marg{col name} 展开为由 \meta{col name} 指定的列在当前行的值。
    此外，|\nextrow|\marg{col name} 展开为指定列在\emph{下}一行的值，
    |\prevrow|\marg{col name} 展开为\emph{上}一行的值。

    因此，其思想就是简单地重定义命令键 |create col/assign|，让它按需要填充新单元格。

    第一行和最后一行有两个特殊的 |assign| 例程：\emph{最后}一行的内容用
    \declaretext{create col/assign last} 计算，其语义相同。\emph{第一}行的内容用
    \declaretext{create col/assign first} 计算，以简化此处的特殊情形。
    这两个首行/末行命令是可选的，默认调用普通的 |assign| 例程。

    |assign| 键的计算在局部 \TeX{} 分组中进行（也就是说之后任何局部定义都会被清除）。

    在单元格赋值期间，下面这些宏很有用：
    %
    \begin{enumerate}
        \item \declareandlabel{\prevrow}\marg{col name} /
            \declareandlabel{\getprevrow}\marg{col name}\marg{\textbackslash
            macro}

            这两个例程返回存储在指定列 \meta{col name} \emph{上}一行的值。
            |get| 例程把它存入 \meta{\textbackslash macro}。

            参数 \meta{col name} 必须表示一个已存在的列名，或一个已存在
            |alias/|\meta{col name} 的列名。
        \item \declareandlabel{\thisrow}\marg{col name} /
            \declareandlabel{\getthisrow}\marg{col name}\marg{\textbackslash
            macro}

            这两个例程返回存储在指定列中的\emph{当前}行的值。|get| 例程把它存入
            \meta{\textbackslash macro}。

            参数 \meta{col name} 必须表示一个已存在的列名，或一个已存在
            |alias/|\meta{col name} 的列名。
        \item \declareandlabel{\nextrow}\marg{col name} /
            \declareandlabel{\getnextrow}\marg{col name}\marg{\textbackslash
            macro}

            这两个例程返回\emph{下}一行的值。

            参数 \meta{col name} 必须表示一个已存在的列名，或一个已存在
            |alias/|\meta{col name} 的列名。
        \item |\pgfplotstablerow| 和 |\pgfplotstablerows|，分别包含当前行的索引和
            行总数。细节见第~\pageref{pgfplotstable:page:tablerow} 页。
        \item |\pgfmathaccuma| 和 |\pgfmathaccumb| 可用于传递中间结果。二者的值会在
            一次列赋值到下一次列赋值之间保持。离开赋值例程后，其他所有局部变量都会被
            删除。二者的初始值都是空字符串，除非它们已被列创建样式初始化。
        \item |\pgfplotstablename|，一个包含当前正在处理的表格名字的宏
            （也就是说它包含 |\pgfplotstablecreatecol| 的第二个参数）。
        \item 在本宏包各部分都有效的命令，例如 |\pgfplotstablerow| 用于取得当前行索引，
            或 |\pgfplotstablerows| 用于取得行总数。
    \end{enumerate}
    %
    \meta{col name} 应为\emph{物理}列名，不允许别名或列索引
    （除非列索引与列名相同）。

    下面的例子取我们熟悉的输入表格，创建 |level| 列的一个副本。此外，它产生大量输出
    以展示可用的宏。最后，它用 |\pgfkeyslet| 把所得 |\entry| 的内容赋给
    |next content|。
    %
\begin{codeexample}[]
\pgfplotstableread{pgfplotstable.example1.dat}\loadedtable
\pgfplotstablecreatecol[
    create col/assign/.code={%
        \getthisrow{level}\entry
        \getnextrow{level}\nextentry
        \edef\entry{thisrow=\entry; nextrow=\nextentry.
            (\#\pgfplotstablerow/\pgfplotstablerows)}%
        \pgfkeyslet{/pgfplots/table/create col/next content}\entry
    }]
    {new}\loadedtable

\pgfplotstabletypeset[
    column type=l,
    columns={level,new},
    columns/new/.style={string type}
]\loadedtable
\end{codeexample}

    还有一个特点：你可以用 |columns=|\marg{column list} 来降低这个命令的运行时复杂度。
    这只在 |columns| 键直接提供在 \meta{options} 中时才有效。此时 |\thisrow| 及其
    变体只为 |columns| 值中列出的那些列定义。


    \paragraph{限制。}

    目前，你一次只能访问一列的三个值：当前行、上一行和下一行。对任意索引的访问
    （尚）不受支持。


    \paragraph{说明：}

    如果你想用这个命令（或相关的 |create on use| 简化形式）从零创建一个表格，
    请参见 |\pgfplotstablenew|。

    |assign| 的默认实现是产生空字符串。|assign last| 的默认实现是调用 |assign|，
    因此如果你从不真正使用下一行的值，就不需要改动 |assign last|。|assign first|
    也是如此。
\end{command}

\begin{pgfplotstablecreateonusekey}
    允许对列 \meta{col name} 进行 ``惰性创建''。每当按名字查询列 \meta{col name} 时
    （例如在 |\pgfplotstabletypeset| 命令中），如果这样一列尚不存在，它就会被即时创建。
    %
\begin{codeexample}[narrow]
% requires \usepackage{array}
\pgfplotstableset{% could be used in preamble
    create on use/quot1/.style=
        {create col/quotient={error1}}}

\pgfplotstabletypeset[
    columns={error1,quot1},
    columns/error1/.style={sci,sci zerofill},
    columns/quot1/.style={dec sep align}]
{pgfplotstable.example1.dat}
\end{codeexample}
    %
    上面的例子查询 |quot1|，而它在输入文件中尚不存在。因此会检查是否存在针对
    |quot1| 的 |create on use| 样式。确实存在，于是就用它来创建这一缺失的列。
    |create col/quotient| 键在下文讨论；它计算列 |error1| 中相邻行的商。

    一条 |create on use| 规范会被转换为

    |\pgfplotstablecreatecol[|\meta{create options}|]|\marg{col name}\marg{the table}，

    或者等价地，转换为

    |\pgfplotstablecreatecol[|create on use/\meta{col name}|]|\marg{col name}\marg{the table}。

    这个特性允许某种惰性，因为你可以省略冗长的表格修改。不过，惰性是有代价的：
    在上面的例子中，生成的列在从 |\pgfplotstabletypeset| 返回后就会\emph{丢失}。

    |create on use| 的优先级高于 |alias|。

    如果 \meta{col name} 包含键设置所需的字符，你需要用花括号把它括起来：
    ``|create on use/{name=wi/th,special}/.style={...}|''。

    关于 |create on use| 的更多例子会在下文讨论可用的列创建样式时给出。

    注意，|create on use| 在 \PGFPlots{} 中也可用，即在 |\addplot table| 与
    |read completely| 键配合使用时。
\end{pgfplotstablecreateonusekey}


\subsection{预定义的列生成方法}

下面这些键既可用于 |\pgfplotstablecreatecol|，也可用于更简单的
|create on use| 框架。


\subsubsection{从某处获取数据}

\begin{stylekey}{/pgfplots/table/create col/set=\marg{value}}
    一个用于列创建上下文的样式，它创建一个新列，并把 \meta{value} 写入每个新单元格。
    该值以字符串形式（逐字）写入。
    %
\begin{codeexample}[]
\pgfplotstableset{
    create on use/my new col/.style={create col/set={--empty--}},
    columns/my new col/.style={string type}
}

\pgfplotstabletypeset[
    columns={level,my new col},
]{pgfplotstable.example1.dat}
\end{codeexample}
\end{stylekey}

\begin{stylekey}{/pgfplots/table/create col/set list=\marg{comma-separated-list}}
    一个用于列创建上下文的样式，它创建一个由 \meta{comma-separated-list} 中各条目
    组成的新列。该值以字符串形式（逐字）写入。

    \meta{comma-separated-list} 通过 \Tikz 的 |\foreach| 命令处理，也就是说你可以用
    |...| 表达式来给出数字（或字符）范围。
    %
\begin{codeexample}[]
\pgfplotstableset{
    create on use/my new col/.style={
        create col/set list={A,B,C,4,50,55,...,100}},
    columns/my new col/.style={string type}
}

\pgfplotstabletypeset[
    columns={level,my new col},
]{pgfplotstable.example1.dat}
\end{codeexample}
    %
    \noindent 新列会被填充或截断到所需的行数。如果列表中的元素不够，就会产生空单元格。
\end{stylekey}

\begin{stylekey}{/pgfplots/table/create col/copy=\marg{column name}}
    一个用于列创建上下文的样式，它直接复制已有的列 \meta{column name}。
    %
\begin{codeexample}[]
\pgfplotstableset{
    create on use/new/.style={create col/copy={level}}
}

\pgfplotstabletypeset[
    columns={level,new},
    columns/new/.style={column name=Copy of level}
]{pgfplotstable.example1.dat}
\end{codeexample}
\end{stylekey}

\begin{stylekey}{/pgfplots/table/create col/copy column from table=\marg{file name or \textbackslash macro}\marg{column name}}
    一个用于列创建上下文的样式，它创建一个由所给表格的 \meta{column name} 中各条目
    组成的新列。该参数既可以是文件名，也可以是一个已加载的表格（即
    |\pgfplotstableread| 所返回的 \meta{\textbackslash macro}）。

    你可以用这个样式（可能配合 |\pgfplotstablenew|）把来自不同表格的某一类公共列
    合并进一个大表格。

    单元格值以字符串形式（逐字）写入。

    \noindent 新列会被填充或截断到所需的行数。如果列表中的元素不够，就会产生空单元格。
\end{stylekey}


\subsubsection{数学运算}

\begin{key}{/pgf/fpu=\mchoice{true,false} (initially true)}
        \index{Precision}
    在开始描述列生成方法之前，先就数学库说一句。其核心始终是由 Mark Wibrow 和
    Till Tantau 编写的 \PGF{} 数学引擎。不过，这个引擎是为生成图形而编写的，
    并不适合科学计算。

    我给 \PGF{} 增加了一个高精度浮点库，它会成为比 \PGF{} $2.00$ 更新的版本的一部分。
    它在 \TeX 中提供 IEEE 双精度计算的完整范围。这个 FPU 也是 \PGFPlotstable{} 的
    一部分，并且默认对 |create col/expr| 和所有其他预定义数学方法激活。

    FPU 不会对新定义的数值型样式激活（尽管它对像 |create col/expr| 这样的预定义
    数学表达式解析样式是激活的）。如果你想添加自己的例程或样式，就需要使用
    %
\begin{codeexample}[code only]
\pgfkeys{/pgf/fpu=true}
\end{codeexample}
    %
    \noindent 以启用扩展精度。标准数学解析器限于 $\pm 16384.00000$ 范围内的定点数。
\end{key}

\begin{stylekey}{/pgfplots/table/create col/expr=\marg{math expression}}
    一个用于 |\pgfplotstablecreatecol| 的样式，它用 \meta{math expression}
    为新列赋值内容。
    %
\begin{codeexample}[]
\pgfplotstableset{
    create on use/new/.style={
        create col/expr={\thisrow{level}*2}}
}

\pgfplotstabletypeset[
    columns={level,new},
    columns/new/.style={column name=$2\cdot $level}
]{pgfplotstable.example1.dat}
\end{codeexample}
    %
    宏 |\thisrow|\marg{col name} 和 |\nextrow|\marg{col name} 可用于使用已有表格的值。

    更多信息请参见 |\pgfplotstablecreatecol|。


    \paragraph{累加列：}

    |expr| 样式在它的第一个列之前把 |\pgfmathaccuma| 初始化为 |0|。每当它计算出一个
    新的列值时，就把 |\pgfmathaccuma| 重定义为该结果。也就是说，你可以在
    \meta{math expression} 内部用 |\pgfmathaccuma| 来累加列。更多细节见
    |create col/expr accum|。


    \paragraph{关于精度和数值范围：}
        \index{Precision}
        \index{Floating Point Unit}

    从 1.2 版起，|expr| 使用一个浮点单元。该 FPU 提供科学计算的完整数据范围，
    相对精度在 $10^{-4}$ 到 $10^{-6}$ 之间。|/pgf/fpu| 键提供更多细节。


    \paragraph{接受的运算：}

    \PGF 的数学解析器配合 FPU 提供下列函数和运算符：

    |+|、|-|、|*|、|/|、|abs|、|round|、|floor|、|mod|、|<|、|>|、|max|、|min|、
    |sin|、|cos|、|tan|、|deg|（弧度转角度）、|rad|（角度转弧度）、|atan|、|asin|、
    |acos|、|cot|、|sec|、|cosec|、|exp|、|ln|、|sqrt|、常数 |pi| 和 |e|、|^|（幂运算）、
    |factorial|\footnote{从 \PGF{} 比 $2.00$ 更新的版本起，你可以用后缀运算符
    \texttt{!} 代替 \texttt{factorial}。}、|rand|（按均匀分布在 $-1$ 到 $1$ 之间取
    随机数）、|rnd|（按均匀分布在 $0$ 到 $1$ 之间取随机数）、数字格式转换 |hex|、
    |Hex|、|oct|、|bin| 等等。数学解析器由 Mark Wibrow 和 Till Tantau 编写~\cite{tikz}，
    FPU 例程是作为 \PGFPlots{} 的一部分开发的。这两部分的文档都可在~\cite{tikz} 中找到。
    \textbf{注意：}三角函数按角度（而非弧度）工作，除非重新配置 |trig format|！
\end{stylekey}

\begin{stylekey}{/pgfplots/table/create col/expr accum=\marg{math expression}\marg{accum initial}}
    |create col/expr| 的一个变体，它还允许定义 |\pgfmathaccuma| 的初始值。
    情形 \meta{accum initial}=|0| \emph{等价于} |expr=|\marg{math expression}。
    %
\begin{codeexample}[]
\pgfplotstableset{
    create on use/new/.style={
        create col/expr={\pgfmathaccuma + \thisrow{level}}},
    create on use/new2/.style={
        create col/expr accum={\pgfmathaccuma * \thisrow{level}}{1}%<- start with `1'
    }
}

\pgfplotstabletypeset[
    columns={level,new,new2},
    columns/new/.style={column name=$\sum$level},
    columns/new2/.style={column name=$\prod$level}
]{pgfplotstable.example1.dat}
\end{codeexample}
    %
    这个例子创建两个列：|new| 列就是 \meta{level} 列中每个值之和（它使用默认的
    |\pgfmathaccuma=0|）。|new2| 列把 |\pgfmathaccuma| 初始化为 |100|，
    然后依次减去 \meta{level} 的值。
\end{stylekey}

\begin{stylekey}{/pgfplots/table/create col/quotient=\marg{column name}}
    一个用于 |\pgfplotstablecreatecol| 的样式，它对由 \meta{column name} 指定的列中
    每个条目 $i = 1,\dotsc, (n-1)$ 计算商 $c_i := m_{i-1} / m_i$。第一个值 $c_0$
    保持为空。
    %
\begin{codeexample}[]
% requires \usepackage{array}
\pgfplotstableset{% configuration, for example, in preamble:
    create on use/quot1/.style={create col/quotient=error1},
    create on use/quot2/.style={create col/quotient=error2},
    columns={error1,error2,quot1,quot2},
    %
    % display styles:
    columns/error1/.style={sci,sci zerofill},
    columns/error2/.style={sci,sci zerofill},
    columns/quot1/.style={dec sep align},
    columns/quot2/.style={dec sep align}
}

\pgfplotstabletypeset{pgfplotstable.example1.dat}
\end{codeexample}
    %
    这个样式使用浮点单元的方法，也就是说它以大约 $10^{-7}$ 的相对精度
    （尾数 $7$ 位有效数字）工作。
\end{stylekey}

\begin{stylekey}{/pgfplots/table/create col/iquotient=\marg{column name}}
    与 |create col/quotient| 类似，但商是倒数。
\end{stylekey}

\begin{stylekey}{/pgfplots/table/create col/dyadic refinement rate=\marg{column name}}
    一个用于 |\pgfplotstablecreatecol| 的样式，它计算列 \meta{column name} 中数据的
    收敛阶 $\alpha$。假定 \meta{column name} 的内容形如 $e_i(h_i) = O(h_i^\alpha)$。
    假设从一行到下一行有二进加密关系 $h_i = h_{i-1}/2$，则有
    $h_{i-1}^\alpha / (h_{i-1}/2)^\alpha = 2^\alpha$，于是用下式求得 $\alpha$：
    %
        \[
            c_i := \log_2\left( \frac{e_{i-1}}{e_i} \right).
        \]
    %
    第一个值 $c_0$ 保持为空。
    %
\begin{codeexample}[]
% requires \usepackage{array}
\pgfplotstabletypeset[% here, configuration options apply only to this single statement:
    create on use/rate1/.style={create col/dyadic refinement rate={error1}},
    create on use/rate2/.style={create col/dyadic refinement rate={error2}},
    columns={error1,error2,rate1,rate2},
    columns/error1/.style={sci,sci zerofill},
    columns/error2/.style={sci,sci zerofill},
    columns/rate1/.style={dec sep align},
    columns/rate2/.style={dec sep align}]
    {pgfplotstable.example1.dat}
\end{codeexample}
    %
    这个样式使用浮点单元的方法，也就是说它以大约 $10^{-6}$ 的相对精度
    （尾数 $6$ 位有效数字）工作。
\end{stylekey}

\begin{stylekey}{/pgfplots/table/create col/idyadic refinement rate=\marg{column name}}
    与 |create col/dyadic refinement rate| 类似，但商是倒数。
\end{stylekey}

\begin{keylist}{%
    /pgfplots/table/create col/gradient=\marg{col x}\marg{col y},
    /pgfplots/table/create col/gradient loglog=\marg{col x}\marg{col y},
    /pgfplots/table/create col/gradient semilogx=\marg{col x}\marg{col y},
    /pgfplots/table/create col/gradient semilogy=\marg{col x}\marg{col y}%
}
    一个用于 |\pgfplotstablecreatecol| 的样式，它对每一行计算分段梯度
    $(y_{i+1} - y_i) / (x_{i+1} - x_i )$。$y$ 值取自列 \meta{col y}，$x$ 值取自
    \meta{col x}。

    对数变体会在开始计算差之前对其参数取自然对数 $\log(\cdot)$。更准确地说，
    |loglog| 变体对 $x$ 和 $y$ 都取对数，|semilogx| 变体只对~$x$ 取对数，
    |semilogy| 变体只对~$y$ 取对数。
    %
\begin{codeexample}[]
% requires \usepackage{array}
\pgfplotstableset{% configuration, for example in preamble:
    create on use/slopes1/.style={create col/gradient loglog={dof}{error1}},
    create on use/slopes2/.style={create col/gradient loglog={dof}{error2}},
    columns={dof,error1,error2,slopes1,slopes2},
    % display styles:
    columns/dof/.style={int detect},
    columns/error1/.style={sci,sci zerofill},
    columns/error2/.style={sci,sci zerofill},
    columns/slopes1/.style={dec sep align},
    columns/slopes2/.style={dec sep align}
}
\pgfplotstabletypeset{pgfplotstable.example1.dat}
\end{codeexample}

\begin{codeexample}[]
% requires \usepackage{array}
\pgfplotstableset{% configuration, for example in preamble:
    create on use/slopes1/.style={create col/gradient semilogy={level}{error1}},
    columns={level,error1,slopes1},
    % display styles:
    columns/level/.style={int detect},
    columns/error1/.style={sci,sci zerofill,sci subscript},
    columns/slopes1/.style={dec sep align}
}
\pgfplotstabletypeset{pgfplotstable.example1.dat}
\end{codeexample}
    %
    这个样式使用浮点单元的方法，也就是说它以大约 $10^{-6}$ 的相对精度
    （尾数 $6$ 位有效数字）工作。
\end{keylist}

\begin{stylekey}{/pgfplots/table/create col/linear regression=\marg{key-value-config}}
    用样本数据 $(x_i,y_i)$ 计算线性（最小二乘）回归 $y(x) = a \cdot x + b$，
    样本数据必须在 \meta{key-value-config} 内部指定。
    %
\begin{codeexample}[pre={\vbox\bgroup\hsize=3cm},post=\egroup]
% load table from somewhere:
\pgfplotstableread{
    x y
    1 1
    2 4
    3 9
    4 16
    5 25
    6 36
}\loadedtbl

% create the `regression' column:
\pgfplotstablecreatecol[linear regression]
    {regression}
    {\loadedtbl}
% store slope
\xdef\slope{\pgfplotstableregressiona}

\pgfplotstabletypeset\loadedtbl\\

The slope is `\slope'.
\end{codeexample}
    %
    \noindent 上面的例子从内联数据加载一个表格，追加一个名为 `|regression|' 的列并排版它。
    由于没有提供 \meta{key-value-config}，将使用 |x=[index]0| 和 |y=[index]1|。
    |\xdef\slope{...}| 命令把回归线的 `$a$' 值存入一个新定义的宏 `|\slope|'。\footnote{\texttt{\textbackslash xdef} 的意思是 ``全局展开定义''：它把参数展开到不能再展开为止，
    并给结果赋一个（全局）名字。细节见任何 \TeX{} 书籍。}

    由于这个特性与绘图关系密切，其完整文档已移到 \PGFPlots{} 中。请参考随本宏包附带的
    \PGFPlots{} 手册。
\end{stylekey}

\begin{stylekey}{/pgfplots/table/create col/function graph cut y=\marg{cut value}\marg{common options}\marg{one key--value set for each plot}}
    一个用于 |create on use| 语句的专用样式，它计算（一个或多个）离散曲线
    $y(x_1), \dotsc, y(x_N)$ 与固定 \meta{cut value} 的交点。$x_i$ 被写入表格的单元格。

    在成本—精度图中，这个特性允许提取固定精度下的成本。与之对偶的 |cut x|
    特性则允许计算固定成本下的精度。
    %
    \pgfplotsset{anchor=center,/tikz/baseline}
\begin{codeexample}[]
\pgfplotstablenew[
    create on use/cut/.style={create col/function graph cut y=
        {2.5e-4} % search for fixed L2 = 2.5e-4
        {x=Basis,y=L2,ymode=log,xmode=log} % double log, each function is L2(Basis)
        % now, provide each single function f_i(Basis):
        {{table=plotdata/newexperiment1.dat},{table=plotdata/newexperiment2.dat}}
    },
    columns={cut}]
    {2}
    \loadedtable

% Show the data:
\pgfplotstabletypeset{\loadedtable}

\begin{tikzpicture}
\begin{loglogaxis}
    \addplot table[x=Basis,y=L2] {plotdata/newexperiment1.dat};
    \addplot table[x=Basis,y=L2] {plotdata/newexperiment2.dat};
    \draw[blue!30!white] (1,2.5e-4) -- (1e5,2.5e-4);
    \node[pin=-90:{$x=53.66$}] at (53.66,2.5e-4) {};
    \node[pin=45:{$x=601.83$}] at (601.83,2.5e-4) {};
\end{loglogaxis}
\end{tikzpicture}
\end{codeexample}
    %
    在上面的例子中，我们在寻找 $x_1$ 和 $x_2$，使得
    $f_1(x_1) = \pgfmathprintnumber{2.5e-4}$ 且
    $f_2(x_2) =\pgfmathprintnumber{2.5e-4}$。左边是自动计算出的结果。右边是一个
    问题示意图，用 \PGFPlots{} 加上适当标注来把结果可视化。\meta{cut value} 设为
    |2.5e-4|。\meta{common options} 包含问题设置；在我们这里是对数刻度和列名。
    第三个参数是一个逗号分隔列表。每个元素 $i$ 是一组键，描述如何取得 $f_i(\cdot)$。

    在 \meta{common options} 和 \meta{one key--value set for each plot} 中，
    都可以使用下列键：
    %
    \begin{itemize}
        \item |table=|\marg{table file or \textbackslash macro}：要么是文件名，
            要么是一个已加载的表格，作为数据点的来源，
        \item |x=|\marg{col name}：$x$ 轴的列名，
        \item |y=|\marg{col name}：$y$ 轴的列名。
        \item |foreach=|\marg{\textbackslash foreach loop
            head}\marg{file name pattern} 这个稍微高级的语法
            允许在一个循环中自动收集表格：
            %
\begin{codeexample}[]
\pgfplotstablenew[
    % same as above...
    create on use/cut/.style={create col/function graph cut y=
        {2.5e-4}% search for fixed L2 = 2.5e-4
        {x=Basis,y=L2,ymode=log,xmode=log,
         foreach={\i in {1,2}}{plotdata/newexperiment\i.dat}}%
        {}% just leave this empty.
    },
    columns={cut}]
    {2}
    \loadedtable
% Show the data:
\pgfplotstabletypeset{\loadedtable}
\end{codeexample}
            %
            \PGFPlotstable{} 会调用 |\foreach |\meta{\textbackslash foreach loop
            head}，并在每次迭代时展开 \meta{file name pattern}。每次迭代都会生成一个
            形如

            |table=|\marg{expanded pattern}|,x=|\marg{value of x}|,y=|\marg{value of y}

            的更简单列表项。

            也可以在 \meta{one key--value set for each plot} 内部提供 |foreach=|。
            |foreach| 键优先于 |table|。|\foreach| 所接受语法的细节见 \pgfname\ 手册。
    \end{itemize}
    %
    键 |xmode| 和 |ymode| 可取 |log| 或 |linear|。上面提到的所有键都有共同的键路径

    \textcolor{red!75!black}{\texttt{/pgfplots/table/create col/function graph cut/}}。
\end{stylekey}

\begin{stylekey}{/pgfplots/table/create col/function graph cut x=\marg{cut value}\marg{common options}\marg{one key--value set for each plot}}
    与上面相同，只是把 $x$ 和 $y$ 互换。
\end{stylekey}


\section{杂项}

\subsection{把（修改后的）表格写入磁盘}

\begin{key}{/pgfplots/table/outfile=\marg{file name} (initially empty)}
    把完全处理后的表格作为 \TeX{} 文件写入 \meta{file name}。这个键在第
    ~\pageref{page:outfile} 页有详细描述。
\end{key}

\begin{command}{\pgfplotstablesave\oarg{options}\marg{\textbackslash macro or input file name}\marg{output file name}}
    这个命令取一个表格并把它写入一个新的数据文件（不进行任何排版）。

    如果第一个参数是文件名，就先加载该文件。

    这个命令只是以清空的输出参数调用 |\pgfplotstabletypeset|。也就是说，任何列创建
    方法在此同样适用，包括任何后处理步骤（但不做最终排版）。

    |\pgfplotstablesave| 使用键 |reset styles| 和 |disable rowcol styles|
    来清除任何排版相关选项。

    此外，它设置 |string type| 以允许逐字输出。如果你只有数值数据，可能想改用
    |numeric as string type|——这样会把算术得到的整数以非科学记数法显示。\footnote{不过要注意，\texttt{string type} 也不会对整数做舍入或截断，
    即便它们以浮点数显示。}
    %
\begin{codeexample}[]
\pgfplotstablesave[
    create on use/postproc1/.style={create col/dyadic refinement rate=error1},
    columns={dof,error1,postproc1}
]
    {pgfplotstable.example1.dat}
    {pgfplotstable.example1.out.dat}
\end{codeexample}
    %
    现在，|pgfplotstable.example1.out.dat| 的内容是
    %
    \lstinputlisting[basicstyle=\ttfamily\footnotesize,tabsize=8]{pgfplotstable.example1.out.dat}

    你可以在 \meta{options} 内部用 |col sep| 键为输出文件定义列分隔符。如果你需要不同的
    输入列分隔符，请用 |in col sep| 代替 |col sep|。


    \paragraph{说明}

    \begin{itemize}
        \item 如果 |col sep=space|，空单元格会被填为 |{}|。可用 |empty cells with|
            样式来改变这一点。
        \item 如果你需要改变基于列/行的样式，请在 \meta{options} 内部使用
            |disable rowcol styles=false|。
    \end{itemize}
\end{command}


\subsection{杂项键}

\begin{key}{/pgfplots/table/disable rowcol styles=\mchoice{true,false} (initially false)}
    如果希望 |\pgfplotstabletypeset| \emph{不}设置任何只适用于特定列或只适用于
    特定行的样式，就把它设为 |true|。

    这会禁用下列样式
    %
    \begin{itemize}
        \item |columns/|\meta{column name},
        \item |display columns/|\meta{column index},
        \item |every col no |\meta{column index},
        \item |every row no |\meta{row index}.
    \end{itemize}
\end{key}

\begin{key}{/pgfplots/table/reset styles}
    重置所有不显式依赖列名/行名和列索引/行索引的表格排版样式。受影响的样式有
    %
    \begin{itemize}
        \item |every table|,
        \item |every even row|, |every odd row|, |every column|,
            |every even column|, |every odd column|,
        \item |every first column|, |every last column|, |every first row|,
            |every last row|,
        \item |every head row|,
        \item |postproc cell content|, |preproc cell content|.
    \end{itemize}
    %
    如果你想全部重置，还应考虑键 |disable rowcol styles|。
\end{key}


\subsection{样式与键的定义和使用小结}

本节小结 |pgfkeys| 的特性。完整文档见 \pgfname{} 手册~\cite{tikz}。

\begin{handler}{{.style}=\marg{key-value-list}}
    定义或重定义样式 \meta{key}。样式是一个普通键，在被设置时会设置
    \meta{key-value-list} 中的所有选项。

    用 |\pgfplotstableset{|\meta{key}|/.style={|\meta{key-value-list}|}}| 在
    命名空间 |/pgfplots/table| 中（重）定义样式 \meta{key}。
\end{handler}

\begin{handler}{{.append style}=\marg{key-value-list}}
    把 \meta{key-value-list} 追加到一个已存在的样式 \meta{key}。这是修改预定义样式的
    首选方式：只做追加可保持与未来版本的兼容。

    用 |\pgfplotstableset{|\meta{key}|/.append style={|\meta{key-value-list}|}}| 把
    \meta{key-value-list} 追加到样式 \meta{key}。这会假定前缀 |/pgfplots/table|。
\end{handler}

\begin{handler}{{.initial}=\marg{value}}
    定义一个新 \meta{key} 并赋值 \meta{value}。
\end{handler}

\begin{handler}{{.add}=\marg{before}\marg{after}}
    通过前插 \meta{before}、后追加 \meta{after} 来修改 \meta{key}。
    %
\begin{codeexample}[]
\pgfplotstableset{columns={a column}}
`\pgfkeysvalueof{/pgfplots/table/columns}';
\pgfplotstableset{columns/.add={}{,another}}
`\pgfkeysvalueof{/pgfplots/table/columns}';
\pgfplotstableset{columns/.add={}{,and one more}}
`\pgfkeysvalueof{/pgfplots/table/columns}'.
\end{codeexample}
    %
    这可在 |\pgfplotsinvokeforeach| 或类似的（无分组！）循环结构内部使用。
\end{handler}

\begin{handler}{{.code}=\marg{\TeX{} code}}
    有时，\PGFPlots{} 用户界面会提供替换其
    例程。这是通过所谓 ``代码键'' 实现的。其含义是用新的 \meta{\TeX{} code}
    替换原来的键及其行为。在 \meta{\TeX{} code} 内部可以使用任何命令。此外，
    模式 |#1| 就是提供给该键的参数。
    %
\begin{codeexample}[]
\pgfplotsset{
    My Code/.code={This is a pgfkeys feature. Argument=`#1'}}
\pgfplotsset{My Code={is here}}
\end{codeexample}
    %
    这个例子定义了一个（新的）名为 |My Code| 的键。本质上它不过是一个接入键值接口的
    |\newcommand|。第二条语句 ``调用'' 这个代码键。
\end{handler}

\begin{handler}{{.append code}=\marg{\TeX{} code}}
    把 \meta{\TeX{} code} 追加到一个已存在的、名为 \meta{key} 的 |/.code| 键。
\end{handler}

\begin{handler}{{.code 2 args}=\marg{\TeX{} code}}
    与 |/.code| 类似，但这个处理程序定义的键接受两个参数。使用这样定义的键时，
    这两个参数可作为 |#1| 和 |#2| 使用。
\end{handler}


\subsection{原生 \TeX{} 与 Con\TeX{}t 支持}
\label{sec:pgfplotstable:context}

表格代码生成器被初始化为生成 \LaTeX{} 的 |tabular| 环境。不过，它只依赖于~`|&|'
作为列分隔符、~`|\\|' 作为行结束符。|column type| 特性或多或少是 |tabular| 特有的，
但你可以完全禁用它。把 |begin table| 和 |end table| 换成合适的 \TeX{} 或
Con\TeX{}t 命令即可改变这一点。如果你有有用的默认样式（或 bug 报告），请告诉我。


\subsection{基本层级的表格访问与修改}

\PGFPlotstable{} 提供若干在基本层级访问和操作表格的方法。

请记住，\PGFPlotstable{} 是作为表格可视化工具编写的。因此它针对行数较少
（尽管列可能很多）的情形做了优化。目前表格创建和修改的运行时间为 $O(N^2)$，
其中 $N$ 是行数。\footnote{\texttt{plot table} 的运行时间借助一个特殊例程随行数线性增长。}
对于行数少的表格这完全可接受，因为 \TeX{} 处理这些结构相对较快。请保持表格小巧！
\PGFPlotstable{} \emph{不是}用于大规模矩阵运算的工具。

表格总是以列向量序列的形式存储。因此，遍历某一列中的所有值很简单，而遍历某一行中的
所有值则复杂且代价高昂。

\begin{command}{\pgfplotstableforeachcolumn\meta{table}\textbackslash as\marg{\textbackslash macro}\marg{code}}
    遍历 \meta{table} 的每个列名。\meta{\textbackslash macro} 会被设为当前访问到的
    列名，然后执行 \meta{code}。在 \meta{code} 期间，|\pgfplotstablecol| 表示当前列索引
    （从 0 开始）。
    %
\begin{codeexample}[]
\begin{minipage}{0.8\linewidth}
\pgfplotstableread{pgfplotstable.example1.dat}\loadedtable
\pgfplotstableforeachcolumn\loadedtable\as\col{%
    column name is `\col'; index is \pgfplotstablecol;\par
}
\end{minipage}
\end{codeexample}

    这个例程不引入 \TeX{} 分组，\meta{code} 内部的变量不受作用域限制。
\end{command}

\begin{command}{\pgfplotstableforeachcolumnelement\meta{column name}\textbackslash of\meta{table}\textbackslash as\meta{\textbackslash cellcontent}\marg{code}}
    以只读模式报告固定列 $j$ 的每个表格单元格 $t_{ij}$。

    对于名为 \meta{column name} 的列中的每个单元格，都会执行 \meta{code}。在这次调用
    期间，宏 \meta{\textbackslash cellcontent} 会包含该单元格的内容，
    |\pgfplotstablerow| 会包含当前行的索引。
    %
\begin{codeexample}[]
\begin{minipage}{0.8\linewidth}
\pgfplotstableread{pgfplotstable.example1.dat}\loadedtable
\pgfplotstableforeachcolumnelement{error1}\of\loadedtable\as\cell{%
    I have now cell element `\cell' at row index `\pgfplotstablerow';\par
}
\end{minipage}
\end{codeexample}
    %
    参数 \meta{column name} 也可以是列索引。此时它应包含 |[index]|\meta{integer}，
    例如 |[index]4|。此外，列别名以及应即时生成的列（见 |create on use|）也可用作
    \meta{column name}。

    这个例程不引入 \TeX{} 分组，\meta{code} 内部的变量不受作用域限制。
\end{command}

\begin{command}{\pgfplotstablemodifyeachcolumnelement\meta{column name}\textbackslash of\meta{table}\textbackslash as\meta{\textbackslash cellcontent}\marg{code}}
    一个与 |\pgfplotstableforeachcolumnelement| 类似的例程，但在 \meta{code} 期间对
    \meta{\textbackslash cellcontent} 所做的任何改动都会写回相应的单元格。
    %
\begin{codeexample}[]
\pgfplotstableread{pgfplotstable.example1.dat}\loadedtable
\pgfplotstablemodifyeachcolumnelement{error1}\of\loadedtable\as\cell{%
    \edef\cell{\#\pgfplotstablerow: \cell}%
}
\pgfplotstabletypeset[columns=error1,string type]{\loadedtable}
\end{codeexample}

    如果 \meta{column name} 是列别名或已即时创建的列，就会创建一个名为
    \meta{column name} 的新列。
\end{command}

\begin{command}{\pgfplotstablegetelem\marg{row}\marg{col}\textbackslash of\meta{table}}
    选取行 \meta{row}、列 \meta{col} 处的单个表格元素。第二个参数的格式与上一段所述
    相同：它应是列名或列索引（若是后者需写成 |[index]|\meta{number}）。

    返回值会写入 |\pgfplotsretval|。
    %
\begin{codeexample}[]
\pgfplotstableread{pgfplotstable.example1.dat}{\loadedtable}
\pgfplotstablegetelem{4}{error1}\of{\loadedtable}
The value (4,error1) is `\pgfplotsretval'.

\pgfplotstablegetelem{2}{[index]0}\of{\loadedtable}
The value (2,0) is `\pgfplotsretval'.
\end{codeexample}

    \paragraph{注意：}

    如果可能，避免在循环内部使用这个命令。它相当慢。
\end{command}

\begin{command}{\pgfplotstablegetcolumnnamebyindex\marg{col index}\textbackslash of\meta{table}\textbackslash to\meta{macro}}
    从 \meta{table} 中取出 \meta{col index} 处的列名，并把它存入 \meta{macro}。
    %
\begin{codeexample}[]
\pgfplotstableread{pgfplotstable.example1.dat}{\loadedtable}
\pgfplotstablegetcolumnnamebyindex{4}\of{\loadedtable}\to\pgfplotsretval
The name of column 4 `\pgfplotsretval'.
\end{codeexample}
    %
    \noindent Note that \meta{table} must be a \emph{loaded} table, i.e.\@ it
    must be a macro resulting from |\pgfplotstableread|. Column indices start
    at~$0$.
\end{command}

\begin{commandlist}{%
    \pgfplotstablegetrowsof\marg{file name or \textbackslash loadedtable},
    \pgfplotstablegetcolsof\marg{file name or \textbackslash loadedtable}%
}
    把 |\pgfplotsretval| 定义为一个表格中的行数或列数。\footnote{它也会把
    \texttt{\textbackslash pgfmathresult} 赋为同一值。} 参数既可以是文件名，
    也可以是一个已加载的表格（即 |\pgfplotstableread| 的 \meta{\textbackslash macro}）。
\end{commandlist}

\begin{command}{\pgfplotstablevertcat\marg{\textbackslash table1}\marg{\textbackslash table2 or filename}}
    关于这个命令的细节见第~\pageref{table:vertcat} 页。
\end{command}

\begin{command}{\pgfplotstablenew\oarg{options}\marg{row count}\marg{\textbackslash table}}
    关于这个命令的细节见节~\ref{pgfplotstable:createcol}。
\end{command}
\begin{command}{\pgfplotstablecreatecol\oarg{options}\marg{row count}\marg{\textbackslash table}}
    关于这个命令的细节见节~\ref{pgfplotstable:createcol}。
\end{command}

\begin{commandlist}{%
    \pgfplotstabletranspose\oarg{options}\marg{\textbackslash outtable}\marg{\textbackslash table or filename},
    \pgfplotstabletranspose*\oarg{options}\marg{\textbackslash outtable}\marg{\textbackslash table or filename}%
}
    把 \meta{\textbackslash outtable} 定义为 \meta{\textbackslash table or filename} 的
    转置。输入参数既可以是文件名，也可以是一个已加载的表格。

    带 `|*|' 的版本只有配合 |columns| 选项时才有意义，见下文。
    %
\begin{codeexample}[]
\pgfplotstabletypeset[string type]{pgfplotstable.example3.dat}
\end{codeexample}

\begin{codeexample}[]
\pgfplotstabletranspose\loadedtable{pgfplotstable.example3.dat}
\pgfplotstabletypeset[string type]\loadedtable
\end{codeexample}

    可选参数 \meta{options} 可以包含影响转置的选项：
    %
    \begin{pgfplotskey}{table/colnames from=\marg{colname} (initially empty)}
        在 |\pgfplotstabletranspose| 内部，这个键处理如何定义输出列名。

        如果 \meta{colname} \emph{为}空（初始值），输出列名就直接是原来的行索引，
        从~$0$ 开始。

        如果 \meta{colname} 非空，它表示一个输入列名，其单元格值将构成输出列名：
        %
\begin{codeexample}[]
\pgfplotstabletranspose[colnames from=c]\loadedtable{pgfplotstable.example3.dat}
\pgfplotstabletypeset[string type]\loadedtable
\end{codeexample}
        %
        参数 \meta{colname} 不会作为单元格内容出现。如果 \meta{colname} 中的单元格不能
        给出唯一的列名，就是错误。
    \end{pgfplotskey}

    \begin{pgfplotskey}{table/input colnames to=\marg{name} (initially colnames)}
        在 |\pgfplotstabletranspose| 内部，这个键处理如何对待\emph{输入}列名。

        这个键会再创建一个名为 \meta{name} 的列，其中填入输入列名（作为字符串类型）。
        %
\begin{codeexample}[]
\pgfplotstabletranspose[input colnames to=Input]\loadedtable{pgfplotstable.example3.dat}
\pgfplotstabletypeset[string type]\loadedtable
\end{codeexample}
        %
        把 \meta{name} 设为空字符串即可禁用这一列。
        %
\begin{codeexample}[]
\pgfplotstabletranspose[input colnames to=]\loadedtable{pgfplotstable.example3.dat}
\pgfplotstabletypeset[string type]\loadedtable
\end{codeexample}
    \end{pgfplotskey}

    \begin{pgfplotskey}{table/columns=\marg{list} (initially empty)}
        在 |\pgfplotstabletranspose| 内部，这个键处理转置时应考虑哪些输入列。

        如果 \meta{list} 为空，就使用输入表格的所有列（这是初始配置）。

        如果 \meta{list} 非空，它应是一个列名列表。只有这些列会作为转置的输入，
        就像其余列不存在一样。在 \meta{list} 内提供列别名或 |create on use| 参数
        是允许的。
        %
\begin{codeexample}[]
\pgfplotstabletranspose[columns={a,b}]\loadedtable{pgfplotstable.example3.dat}
\pgfplotstabletypeset[string type]\loadedtable
\end{codeexample}

        这就是 |\pgfplotstabletranspose| 与 |\pgfplotstabletranspose*| 之间唯一的区别：
        不带 `|*|' 的版本在开始前会\emph{重置} |columns| 键，而带 `|*|' 的版本直接使用
        |columns| 的当前内容。
    \end{pgfplotskey}
\end{commandlist}

\begin{command}{\pgfplotstablesort\oarg{options}\meta{\textbackslash resulttable}\meta{\textbackslash table or filename}}
    按 \meta{options} 对 \meta{\textbackslash table or filename} 排序，并把排序后的
    表格写入 \meta{\textbackslash resulttable}。

    使用高层的 |sort| 键可在 |\pgfplotstabletypeset| 期间自动启用排序。
    %
\begin{codeexample}[]
\pgfplotstablesort\result{%
    a      b  c
    19     2 [a]
    -6   -14 [b]
     4   -14 [c]
   -11    -9 [d]
    11    14 [e]
    -9    -9 [f]
     1    13 [g]
     8   -10 [h]
    16    18 [i]
    19    -6 [j]
}
\pgfplotstabletypeset[columns/c/.style={string type}]{\result}%
\end{codeexample}

    排序键和比较函数可用下列键定制：

    \begin{key}{/pgfplots/table/sort key=\marg{column} (initially [index]0)}
        指定含有排序键的列。参数 \meta{column} 可以是输入表格中的任何一列，
        包括 |create on use|、|alias| 或 |[index]|\meta{integer} 规范。初始设置使用
        第一个可用列。
    \end{key}

    \begin{key}{/pgfplots/table/sort key from=\marg{table} (initially empty)}
        允许从另一个 \meta{table} 加载 |sort key|，该 \meta{table} 既可以是
        \meta{\textbackslash macro}，也可以是 \meta{file name}。
    \end{key}

    \begin{key}{/pgfplots/table/sort cmp=\marg{less than routine} (initially float <)}
        允许使用不同的比较函数。

        \begin{pgfplotskeylist}{%
            fixed <,
            fixed >,
            int <,
            int >,
            float <,
            float >,
            date <,
            date >,
            string <,
            string >%
        }
            这些样式构成预定义的比较函数。|fixed <|、|int <| 和 |float <| 例程作用于
            数值型数据，其中 |int <| 期待正整数或负整数，另两个期待实数。|fixed <|
            的数值范围小得多，但比 |float <| 略快。

            |date <| 比较形如 |YYYY-MM-DD| 的日期。|string <| 基于排序键的 ASCII
            字符码做字典序字符串比较。|string <| 例程还会求值控制序列或活动字符的
            ASCII 码。\footnote{只要它们尚未被 \TeX 的预处理消耗掉。}
            %
\begin{codeexample}[]
\pgfplotstablesort[sort cmp=string <]\result{% 'Header' is the column name:
    Header
    the
    quick
    brown
    fox
    jumps
    over
    the
    lazy
    dog
}
\pgfplotstabletypeset[string type]{\result}%
\end{codeexample}
        \end{pgfplotskeylist}

        参数 \meta{less than routine} 实际上是一个样式：它应是带 |/pgfplots/| 键前缀的
        一个或多个键。因此，所有这些预定义的排序比较其实都是定义 |iflessthan| 的键。
        例如，|string <| 的定义如下：
        %
\begin{codeexample}[code only]
\makeatletter
\pgfplotsset{
    string </.style={%
        iflessthan/.code args={##1##2##3##4}{%
            \t@pgfplots@toka=\expandafter{##1}%
            \t@pgfplots@tokb=\expandafter{##2}%
            \edef\pgfplots@loc@TMPa{{\the\t@pgfplots@toka}{\the\t@pgfplots@tokb}}%
            \expandafter\pgfplotsutilstrcmp\pgfplots@loc@TMPa
            \if1\pgfplotsretval ##3\else ##4\fi
        }%
    },%
}
\end{codeexample}

        \begin{codeargskey}{/pgfplots/iflessthan}{#1#2#3#4}
            允许定义自定义比较函数（严格序）。它比较 |#1 < #2|，若比较为真则调用 |#3|，
            否则调用 |#4|。比较在局部作用域中求值（局部变量之后会被释放）。

            为了能用作 |sort cmp| 的值，|iflessthan| 例程需要在值
            |sort cmp=|\meta{less than routine} 内部定义：
            %
\begin{codeexample}[code only]
\pgfplotsset{
    sort cmp={
        iflessthan/.code args={#1#2#3#4}{%
            % here comes your code
        }%
    },
}
\end{codeexample}
        \end{codeargskey}
    \end{key}
\end{command}


\subsection{重复做事：循环}

\begin{command}{\foreach \meta{variables} |in| \meta{list} \marg{commands}}
    一个由 \Tikz 提供的强大循环命令，见~\cite[Section ``Utilities'']{tikz}。
    %
\begin{codeexample}[]
\foreach \x in {1,2,...,4} {Iterating \x. }%
\end{codeexample}

    一个与 \PGFPlots{} 相关的例子可以是
    %
\begin{codeexample}[code only]
\foreach \i in {1,2,...,10} {\addplot table {datafile\i}; }%
\end{codeexample}
\end{command}

\noindent 下面这些循环命令随 \PGFPlots{} 提供。它们与 \Tikz{} 强大的 |\foreach|
循环命令类似，但后者并非总能用于表格处理：它的循环体在每次迭代后作用即结束。

下面这些 \PGFPlots{} 循环宏是一种替代。

\begin{command}{\pgfplotsforeachungrouped \meta{variable} |in| \meta{list} \marg{command}}
    |\foreach| 的一个专用变体，它能做两件事：执行 \meta{command} 时不引入额外分组，
    并且允许对（简单的！）\meta{$x_0$}|,|\meta{$x_1$}|,...,|\meta{$x_n$}
    表达式调用数学解析器。
    %
\begin{codeexample}[]
\def\allcollected{}
\pgfplotsforeachungrouped \x in {1,2,...,4} {Iterating \x. \edef\allcollected{\allcollected, \x}}%
All collected = \allcollected.
\end{codeexample}

    一个更有用的例子可能是处理表格：
    %
\begin{codeexample}[code only]
\pgfplotsforeachungrouped \i in {1,2,...,10} {%
    \pgfplotstablevertcat{\output}{datafile\i} % appends `datafile\i' -> `\output'
}%
% since it was ungrouped, \output is still defined (would not work
% with \foreach)
\end{codeexample}

    \paragraph{说明：}

    特殊语法
    \meta{list}=\meta{$x_0$}|,|\meta{$x_1$}|,...,|\meta{$x_n$}，
    即前面两个元素、后跟省略号和一个末尾元素，会为该循环调用数学解析器。因此，
    如果 |/pgf/fpu| 处于激活状态，它比任何其他语法都允许更大的数值范围。在其他所有
    情形下，|\pgfplotsforeachungrouped| 调用 |\foreach| 并不带 \TeX{} 分组地提供结果。
\end{command}

\begin{command}{\pgfplotsinvokeforeach\marg{list} \marg{command}}
    |\pgfplotsforeachungrouped|（因而也是 |\foreach|）的一个变体，它对 \meta{list} 中
    每个元素，把 \meta{command} 内部出现的每个 |#1| 都替换一次。也就是说，它实际上
    假定 \meta{command} 像一个 |\newcommand| 的函数体。

    换句话说，对 \meta{list} 的每个元素都会调用 \meta{command}。\meta{list} 的当前元素
    可作为 |#1| 使用。

    与 |\pgfplotsforeachungrouped| 一样，这个命令\emph{不}引入额外作用域
    （也就是说它也是无分组的）。

    它与 |\foreach \x in |\meta{list}\marg{command} 的区别很细微：|\x| 不会被展开，
    而 |#1| 会。
    %
\begin{codeexample}[]
\pgfkeys{
  otherstyle a/.code={[a]},
  otherstyle b/.code={[b]},
  otherstyle c/.code={[c]},
  otherstyle d/.code={[d]}}
\pgfplotsinvokeforeach{a,b,c,d}
    {\pgfkeys{key #1/.style={otherstyle #1}}}
Invoke them: \pgfkeys{key a} \pgfkeys{key b} \pgfkeys{key c} \pgfkeys{key d}
\end{codeexample}
    %
    反例则用一个宏（这里是 |\x|）作为循环参数：
    %
\begin{codeexample}[]
\pgfkeys{
  otherstyle a/.code={[a]},
  otherstyle b/.code={[b]},
  otherstyle c/.code={[c]},
  otherstyle d/.code={[d]}}
\pgfplotsforeachungrouped \x in {a,b,c,d}
    {\pgfkeys{key \x/.style={otherstyle \x}}}
Invoke them: \pgfkeys{key a} \pgfkeys{key b} \pgfkeys{key c} \pgfkeys{key d}
\end{codeexample}

    \paragraph{限制：}

    你目前还不能嵌套这个命令（因为它不提供作用域保护）。
\end{command}


\printindex

\bibliographystyle{abbrv} %gerapali} %gerabbrv} %gerunsrt.bst} %gerabbrv}% gerplain}
\bibliography{pgfplots}
\end{document}
% vi: spell spelllang=en
